% Lesson 35 — Vocabulary: Words and expressions related to garbage collection, disposal, and recycling services — Anglais Tle A — Cours signé OnBuch+
% Programme officiel MINESEC. Compiler avec : tectonic EN35-vocabulary-words-and-expressions-related-to-garbage-col.tex
\newcommand{\DOCMATIERE}{Anglais}
\newcommand{\DOCNIVEAU}{Tle A}
\newcommand{\DOCMODULE}{Chapter 3 : Environment, Well-being and Health}
\newcommand{\DOCLECON}{EN35}
\newcommand{\DOCTITRE}{Lesson 35 — Vocabulary: Words and expressions related to garbage collection, disposal, and recycling services}
% OnBuch+ — préambule « cours signés » ENRICHI (compilé avec tectonic / XeLaTeX)
% Système visuel « Pop » d'OnBuch+ : fond crème (jamais blanc), bordures encre
% épaisses, ombres « dures » décalées, titres Archivo Black en capitales.
% Version MATHÉMATIQUES : graphes (pgfplots/tikz), boîtes pédagogiques
% (définition, propriété, méthode, attention, le savais-tu, exercice résolu,
% exercice, TP, bilan), page de garde et signature OnBuch+ ancrée en bas.
%
% Macros attendues dans main.tex AVANT \input{preamble} :
%   \DOCMATIERE  \DOCNIVEAU  \DOCMODULE  \DOCLECON (numéro)  \DOCTITRE
\documentclass[11pt]{article}

\usepackage{fontspec}
\usepackage[english,french]{babel} % français = langue principale ; l'anglais via \eng{..} / textebox
\usepackage[a4paper,margin=2cm,top=3.4cm,bottom=2.8cm,headheight=1.4cm,footskip=1.3cm]{geometry}
\usepackage{amsmath,amssymb,amsfonts}
\usepackage{mathtools} % \xmapsto, \coloneqq... souvent utilisés par les modèles
\usepackage{cancel} % simplifications barrées (bilans de forces)
% \abs{z} (module, valeur absolue) et \norm{u} (norme), utilisés par les modèles
\providecommand{\abs}{}\let\abs\relax\DeclarePairedDelimiter{\abs}{\lvert}{\rvert}
\providecommand{\norm}{}\let\norm\relax\DeclarePairedDelimiter{\norm}{\lVert}{\rVert}
\usepackage[table]{xcolor} % [table] : fournit aussi \rowcolors (alternance de lignes)
\usepackage{pagecolor}
\usepackage{tikz}
\usetikzlibrary{arrows.meta,positioning,calc,shapes.geometric,shapes.misc,shapes.arrows,decorations.pathmorphing,decorations.pathreplacing,decorations.markings,patterns,shadows,mindmap,trees,fit,backgrounds,matrix,intersections,angles,quotes}
\usepackage{pgfplots}
\usepackage[version=4]{mhchem} % équations nucléaires \ce{^{235}_{92}U}
\usepackage[european,siunitx]{circuitikz} % schémas électriques
\pgfplotsset{compat=1.18}
\usepgfplotslibrary{fillbetween,groupplots} % aires entre courbes (calcul intégral)
\usepackage{enumitem}
\usepackage{fancyhdr}
% [most] charge toutes les bibliothèques tcolorbox (listings, minted, raster,
% documentation...) et gonfle inutilement la mémoire TeX sur un document long
% (~300 boîtes) : on ne charge que ce qui est réellement utilisé.
\usepackage[skins,breakable]{tcolorbox}
\usepackage{graphicx}
\usepackage{colortbl}
\usepackage{booktabs}
\usepackage{tabularx}
\usepackage{diagbox} % case d'en-tête diagonale des tableaux à double entrée
% Colonnes extensibles centrée / gauche / droite (C, L, R), souvent utilisées par les modèles
\newcolumntype{C}{>{\centering\arraybackslash}X}
\newcolumntype{L}{>{\raggedright\arraybackslash}X}
\newcolumntype{R}{>{\raggedleft\arraybackslash}X}
\usepackage{array}
\usepackage{multicol}
\usepackage{multirow}
\usepackage{siunitx}
% parse-numbers=false : accepte aussi \SI{2,5 \times 10^{-3}}{...}
\sisetup{locale=FR,output-decimal-marker={,},group-minimum-digits=4,parse-numbers=false}
\DeclareSIUnit\ohm{\ensuremath{\symup{\Omega}}} % U+2126 absent de la police
\DeclareSIUnit\division{div} % réglages d'oscilloscope (ms/div, V/div)
\DeclareSIUnit\tr{tr} % tours (vitesse de rotation, ex. \SI{1500}{\tr\per\minute})
\DeclareSIUnit\molar{M} % concentration molaire (ex. \SI{3}{\molar}, \SI{1}{\micro\molar})
\DeclareSIUnit\an{an} % année (le modèle confond parfois avec \year, un compteur TeX) — ex. \SI{5}{\kilo\meter\per\an}
\DeclareSIUnit\FCFA{FCFA} % franc CFA (ex. \SI{500}{\FCFA\per\kilo\gram})
\DeclareSIUnit\degreeBrix{\textdegree Brix} % teneur en sucre d'un sirop/jus (réfractométrie)
\DeclareSIUnit\month{mois} % le modèle utilise parfois \month, un compteur TeX, comme unité de durée
\usepackage{qrcode}
% \degree : commande fréquemment utilisée par les modèles pour un angle ou une
% température en dehors de \SI{}{} (non fournie par les paquets chargés).
\providecommand{\degree}{^{\circ}}
% \qed (fin de démonstration), utilisé par les modèles sans amsthm
\providecommand{\qed}{\hfill$\square$}
% \barre : double barre des valeurs interdites dans un tableau de variations (tkz-tab)
\providecommand{\barre}{\ensuremath{\|}}
% environnement proof (amsthm non chargé) utilisé par les modèles
\ifdefined\proof\else\newenvironment{proof}[1][Démonstration]{\par\noindent\textit{#1.}\ }{\qed\par}\fi
\usepackage{lastpage}
\usepackage{hyperref}
\hypersetup{hidelinks}

% ============================================================
% Palette « Pop » OnBuch+ — mêmes hex que la classe P (plus_ui.dart)
% ============================================================
\definecolor{poporange}{HTML}{F59321}
\definecolor{popdark}{HTML}{E07A0C}
\definecolor{popcream}{HTML}{FFF6E8}
\definecolor{popink}{HTML}{1C1714}
\definecolor{poppurple}{HTML}{7A5AE0}
\definecolor{popgreen}{HTML}{2E9E6A}
\definecolor{poppink}{HTML}{E0457B}
\definecolor{popblue}{HTML}{2F7DE1}
\definecolor{popgold}{HTML}{C98A12}
\definecolor{popmuted}{HTML}{8A7F76}
\definecolor{popline}{HTML}{EADFD2}
\definecolor{popgray}{HTML}{8A7F76}
\colorlet{popgrey}{popgray}
% Alias tolérants (le modèle nomme parfois les couleurs différemment)
\colorlet{popwhite}{white}
\colorlet{popblack}{popink}
\colorlet{popred}{poppink}
\colorlet{popyellow}{popgold}
% Teintes claires (fonds de tableaux/encadrés) — le modèle nomme parfois
% les couleurs avec un suffixe L (ex. popblueL) sans les avoir définies.
\colorlet{poporangeL}{poporange!20}
\colorlet{popdarkL}{popdark!20}
\colorlet{poppurpleL}{poppurple!20}
\colorlet{popgreenL}{popgreen!20}
\colorlet{poppinkL}{poppink!20}
\colorlet{popblueL}{popblue!20}
\colorlet{popgoldL}{popgold!20}
\colorlet{popredL}{popred!20}
\colorlet{popgrayL}{popgray!20}
% Teintes claires pour fonds de tableaux / zones
\colorlet{popblueL}{popblue!10!white}
\colorlet{poporangeL}{poporange!14!white}
\colorlet{popgreenL}{popgreen!12!white}
\colorlet{poppurpleL}{poppurple!12!white}
\colorlet{poppinkL}{poppink!12!white}
\colorlet{popgoldL}{popgold!14!white}

\pagecolor{popcream}

% ============================================================
% Typographie
% ============================================================
\defaultfontfeatures{Ligatures=TeX}
\setmainfont{PlusJakartaSans-Medium.ttf}[Path=../../../fonts/,BoldFont=PlusJakartaSans-Bold.ttf]
% Symboles, API et emojis absents de la police principale : repli sur DejaVu Sans (caractères actifs)
\newfontface\symfont{DejaVuSans.ttf}[Path=../../../fonts/]
\makeatletter
\newcommand\symactive[1]{\catcode#1=13 \begingroup\lccode`\~=#1 \lowercase{\endgroup\def~}{{\symfont\char#1}}}
\newcommand\symblank[1]{\catcode#1=13 \begingroup\lccode`\~=#1 \lowercase{\endgroup\def~}{}}
\newcommand\symhyph[1]{\catcode#1=13 \begingroup\lccode`\~=#1 \lowercase{\endgroup\def~}{\mbox{-}}}
\newcount\symc
\newcommand\symrange[2]{\symc=#1 \@whilenum\symc<#2 \do{\expandafter\symactive\expandafter{\the\symc}\advance\symc by 1 }}
\newcommand\symblankrange[2]{\symc=#1 \@whilenum\symc<#2 \do{\expandafter\symblank\expandafter{\the\symc}\advance\symc by 1 }}
\makeatother
\symrange{"0250}{"02FF}\symactive{"2713}\symactive{"2714}\symactive{"2717}\symactive{"2718}\symactive{"2610}\symactive{"2611}\symactive{"2612}
\symhyph{"2011}\symhyph{"2E1A}\symblankrange{"1F300}{"1FAFF}\symblank{"FE0F}
% Maths en sans-serif (Fira Math) avec chiffres et lettres droites de Plus
% Jakarta, pour que les formules se fondent dans le texte.
\usepackage[math-style=ISO,bold-style=ISO]{unicode-math}
\setmathfont{FiraMath-Regular.otf}[Path=../../../fonts/]
\setmathfont{PlusJakartaSans-Medium.ttf}[Path=../../../fonts/,range={up/num,up/{Latin,latin},\mathup}]
\usepackage{icomma} % virgule décimale sans espace en mode math
% Caractères mathématiques Unicode absents des polices de texte (Plus Jakarta,
% Archivo Black) : rendus par leur équivalent mathématique, en texte comme en
% formule — sinon ils s'affichent en carré vide (ex. « Arithmétique dans ℤ »).
\usepackage{newunicodechar}
\newunicodechar{ℝ}{\ensuremath{\mathbb{R}}}
\newunicodechar{ℤ}{\ensuremath{\mathbb{Z}}}
\newunicodechar{ℂ}{\ensuremath{\mathbb{C}}}
\newunicodechar{ℕ}{\ensuremath{\mathbb{N}}}
\newunicodechar{ℚ}{\ensuremath{\mathbb{Q}}}
\newunicodechar{𝔻}{\ensuremath{\mathbb{D}}}
\newunicodechar{α}{\ensuremath{\alpha}}
\newunicodechar{β}{\ensuremath{\beta}}
\newunicodechar{γ}{\ensuremath{\gamma}}
\newunicodechar{δ}{\ensuremath{\delta}}
\newunicodechar{ε}{\ensuremath{\varepsilon}}
\newunicodechar{θ}{\ensuremath{\theta}}
\newunicodechar{λ}{\ensuremath{\lambda}}
\newunicodechar{μ}{\ensuremath{\mu}}
\newunicodechar{σ}{\ensuremath{\sigma}}
\newunicodechar{τ}{\ensuremath{\tau}}
\newunicodechar{φ}{\ensuremath{\varphi}}
\newunicodechar{ψ}{\ensuremath{\psi}}
\newunicodechar{ω}{\ensuremath{\omega}}
\newunicodechar{Σ}{\ensuremath{\Sigma}}
% majuscules produites par \MakeUppercase dans les titres (α-aminés -> Α)
\newunicodechar{Α}{\ensuremath{\alpha}}
\newunicodechar{Β}{\ensuremath{\beta}}
\newunicodechar{Γ}{\ensuremath{\Gamma}}
\newunicodechar{Δ}{\ensuremath{\Delta}}
\newunicodechar{Ε}{\ensuremath{\varepsilon}}
\newunicodechar{Θ}{\ensuremath{\Theta}}
\newunicodechar{Λ}{\ensuremath{\Lambda}}
\newunicodechar{Μ}{\ensuremath{\mu}}
\newunicodechar{Τ}{\ensuremath{\tau}}
\newunicodechar{Φ}{\ensuremath{\Phi}}
\newunicodechar{Ψ}{\ensuremath{\Psi}}
\newunicodechar{Ω}{\ensuremath{\Omega}}
\newunicodechar{↦}{\ensuremath{\mapsto}}
\newunicodechar{∈}{\ensuremath{\in}}
\newunicodechar{∉}{\ensuremath{\notin}}
\newunicodechar{∧}{\ensuremath{\wedge}}
\newunicodechar{⇔}{\ensuremath{\Leftrightarrow}}
\newunicodechar{⟺}{\ensuremath{\Longleftrightarrow}}
\newunicodechar{⇒}{\ensuremath{\Rightarrow}}
\newunicodechar{⟹}{\ensuremath{\Longrightarrow}}
\newunicodechar{∘}{\ensuremath{\circ}}
\newunicodechar{∪}{\ensuremath{\cup}}
\newunicodechar{∩}{\ensuremath{\cap}}
\newunicodechar{≡}{\ensuremath{\equiv}}
\newunicodechar{⊂}{\ensuremath{\subset}}
\newunicodechar{⊥}{\ensuremath{\perp}}
\newunicodechar{∃}{\ensuremath{\exists}}
\newunicodechar{∀}{\ensuremath{\forall}}
\newunicodechar{∛}{\ensuremath{\sqrt[3]{\,}}}
\newunicodechar{∜}{\ensuremath{\sqrt[4]{\,}}}
\newunicodechar{⃗}{\ensuremath{{}^{\to}}}
\newunicodechar{ⁿ}{\ifmmode^{n}\else\textsuperscript{\ensuremath{n}}\fi}
\newunicodechar{ˣ}{\ifmmode^{x}\else\textsuperscript{\ensuremath{x}}\fi}
\newunicodechar{ᵇ}{\ifmmode^{b}\else\textsuperscript{\ensuremath{b}}\fi}
\newunicodechar{ʳ}{\ifmmode^{r}\else\textsuperscript{\ensuremath{r}}\fi}
\newunicodechar{ᵘ}{\ifmmode^{u}\else\textsuperscript{\ensuremath{u}}\fi}
\newunicodechar{ʸ}{\ifmmode^{y}\else\textsuperscript{\ensuremath{y}}\fi}
\newunicodechar{ᵃ}{\ifmmode^{a}\else\textsuperscript{\ensuremath{a}}\fi}
\newunicodechar{ᵉ}{\ifmmode^{e}\else\textsuperscript{\ensuremath{e}}\fi}
\newunicodechar{ᵏ}{\ifmmode^{k}\else\textsuperscript{\ensuremath{k}}\fi}
\newunicodechar{ᵖ}{\ifmmode^{p}\else\textsuperscript{\ensuremath{p}}\fi}
\newunicodechar{ᴺ}{\ifmmode^{N}\else\textsuperscript{\ensuremath{N}}\fi}
\newunicodechar{ᶜ}{\ifmmode^{c}\else\textsuperscript{\ensuremath{c}}\fi}
\newunicodechar{ʰ}{\ifmmode^{h}\else\textsuperscript{\ensuremath{h}}\fi}
\newunicodechar{ᵠ}{\ifmmode^{q}\else\textsuperscript{\ensuremath{q}}\fi}
\newunicodechar{ₙ}{\ifmmode_{n}\else\textsubscript{\ensuremath{n}}\fi}
\newunicodechar{ₐ}{\ifmmode_{a}\else\textsubscript{\ensuremath{a}}\fi}
\newunicodechar{ᵢ}{\ifmmode_{i}\else\textsubscript{\ensuremath{i}}\fi}
\newunicodechar{ᵣ}{\ifmmode_{r}\else\textsubscript{\ensuremath{r}}\fi}
\newunicodechar{ₖ}{\ifmmode_{k}\else\textsubscript{\ensuremath{k}}\fi}
\newunicodechar{ᵦ}{\ifmmode_{\beta}\else\textsubscript{\ensuremath{\beta}}\fi}
\newunicodechar{ₓ}{\ifmmode_{x}\else\textsubscript{\ensuremath{x}}\fi}
\newfontfamily\popdisplay{ArchivoBlack-Regular.ttf}[Path=../../../fonts/]
\newfontfamily\poplogo{SpaceGrotesk-Bold.ttf}[Path=../../../fonts/]
\newfontfamily\popsemi{PlusJakartaSans-SemiBold.ttf}[Path=../../../fonts/]
\newfontfamily\popxbold{PlusJakartaSans-ExtraBold.ttf}[Path=../../../fonts/]
\newfontfamily\popxboldls{PlusJakartaSans-ExtraBold.ttf}[Path=../../../fonts/,LetterSpace=3.0]

\setlist[enumerate]{leftmargin=1.7em,itemsep=5pt,topsep=4pt}
\setlist[itemize]{leftmargin=1.7em,itemsep=4pt,topsep=4pt,label={\color{poporange}\textbullet}}
\setlength{\parindent}{0pt}
\setlength{\parskip}{5pt}
\renewcommand{\arraystretch}{1.25}

% pgfplots : style Pop par défaut
\pgfplotsset{
  every axis/.append style={axis line style={popink,line width=1pt},tick style={popink},
    grid style={popline,line width=0.5pt},label style={font=\small},tick label style={font=\footnotesize}},
  popcurve/.style={poporange,line width=1.8pt,no markers,smooth},
}
% Même style disponible dans un \draw TikZ simple (hors pgfplots)
\tikzset{popcurve/.style={poporange,line width=1.8pt}}
\tikzset{popgrid/.style={popline,very thin}}
\colorlet{popcurve}{poporange} % le nom du style est parfois utilisé comme couleur

% ============================================================
% Pill / squiggle
% ============================================================
\newcommand{\poppill}[3]{%
  \tikz[baseline=-0.55ex]\node[draw=popink,line width=1.2pt,rounded corners=2.6pt,%
    fill=#2,inner xsep=6.5pt,inner ysep=3pt]{%
    \popxboldls\fontsize{7}{7}\selectfont\color{#3}\MakeUppercase{#1}};%
}
\newcommand{\popsquiggle}[1]{%
  \tikz[baseline=-0.4ex]{
    \draw[#1,line width=2.6pt,line cap=round]
      plot [smooth] coordinates {(0,0.09) (0.34,0.01) (0.68,0.21) (1.02,0.01) (1.36,0.10)};
  }%
}
% Mot-clé surligné (marqueur orange sous le texte)
\newcommand{\cle}[1]{{\popxbold\color{popdark}#1}}

% ============================================================
% En-tête / pied
% ============================================================
\pagestyle{fancy}
\fancyhf{}
\renewcommand{\headrulewidth}{0pt}
\renewcommand{\footrulewidth}{0pt}
\lhead{\raisebox{-0.15ex}{\poplogo\fontsize{13}{13}\selectfont\color{popink}OnBuch\color{poporange}+}}
\rhead{\poppill{\DOCMATIERE\ \textbullet\ \DOCNIVEAU\ \textbullet\ Leçon \DOCLECON}{white}{popink}}
\lfoot{\popxbold\fontsize{7.6}{7.6}\selectfont\color{popmuted}\MakeUppercase{Cours signé OnBuch+}\ \textbullet\ {\popsemi\DOCTITRE}}
\rfoot{\tikz[baseline=-0.6ex]\node[rounded corners=8.5pt,fill=popink,minimum height=17pt,minimum width=17pt,inner xsep=5pt,inner sep=0]{\popxbold\fontsize{8}{8}\selectfont\color{popcream}\thepage\,/\,\pageref*{LastPage}};}

% ============================================================
% Titres
% ============================================================
\newcommand{\chaptitre}[2]{%
  \begin{tcolorbox}[enhanced,colback=poporange,colframe=popink,boxrule=2.6pt,arc=11pt,
      drop shadow=popink,left=15pt,right=15pt,top=12pt,bottom=11pt,before skip=3pt,after skip=18pt]
    \poppill{#2}{popink}{popcream}\par\vspace{9pt}
    {\raggedright\hyphenpenalty=10000\exhyphenpenalty=10000\popdisplay\fontsize{22}{24}\selectfont\color{popcream}\MakeUppercase{#1}\par}
  \end{tcolorbox}
}
% Section numérotée : pastille orange avec le numéro + titre Archivo
\newcounter{popsec}
\newcommand{\coursec}[1]{%
  \stepcounter{popsec}\setcounter{popsub}{0}%
  \par\vspace{16pt}\noindent
  \begin{minipage}{\linewidth}
  \tikz[baseline=-0.7ex]\node[draw=popink,line width=1.6pt,fill=poporange,rounded corners=4pt,minimum size=22pt,inner sep=0]{\popdisplay\fontsize{12}{12}\selectfont\color{popcream}\thepopsec};%
  \hspace{8pt}{\popdisplay\fontsize{15}{17}\selectfont\color{popink}\MakeUppercase{#1}}\par
  \vspace{4pt}\noindent{\color{poporange}\rule{36pt}{3.6pt}}
  \end{minipage}\par\nopagebreak\vspace{8pt}
}
\newcounter{popsub}
\newcommand{\courssub}[1]{%
  \stepcounter{popsub}%
  \par\vspace{9pt}\noindent{\popxbold\fontsize{11.5}{13}\selectfont\color{popink}{\color{poporange}\thepopsec.\thepopsub}\hspace{6pt}#1}\par\nopagebreak\vspace{4pt}
}

% ============================================================
% Boîtes pédagogiques — même recette : fond blanc, bordure épaisse colorée,
% ombre dure encre, badge plein. Titre optionnel [#1].
% ============================================================
\tcbset{popbase/.style={breakable,enhanced,colback=white,boxrule=2.3pt,arc=9pt,
  shadow={3pt}{-3pt}{0pt}{popink},left=14pt,right=14pt,top=11pt,bottom=11pt,before skip=11pt,after skip=15pt,
  fontupper=\popsemi\fontsize{10.6}{14.5}\selectfont\color{popink},
  fontlower=\popsemi\fontsize{10.6}{14.5}\selectfont\color{popink}}}
% Coche dessinée (le glyphe ✓ n'existe pas dans les polices du document)
\newcommand{\popcheck}[1]{\tikz[baseline=-0.5ex]\draw[#1,line width=1.6pt,line cap=round,line join=round](-0.13,0.0)--(-0.04,-0.09)--(0.13,0.1);}
\providecommand{\coche}{\popcheck{popgreen}}
\providecommand{\resultat}[1]{\par\smallskip\noindent\fcolorbox{popgreen}{popgreenL}{\parbox{\dimexpr\linewidth-2\fboxsep-2\fboxrule\relax}{\textbf{Résultat.} #1}}\par\smallskip}
\newcommand{\popboxhead}[3]{\poppill{#1}{#2}{white}\ifx\relax#3\relax\else\hspace{6pt}{\popxbold\fontsize{10.6}{12}\selectfont\color{popink}#3}\fi\par\vspace{7pt}}

\newtcolorbox{aretenir}[1][]{popbase,colframe=popgreen,before upper={\popboxhead{À retenir}{popgreen}{#1}}}
% Texte en anglais (document de lecture, dialogue, extrait) : cadre
% d'étude. Un titre optionnel (source, auteur) s'affiche dans l'en-tête.
\newtcolorbox{textebox}[1][]{popbase,colframe=popblue,before upper={\popboxhead{Text}{popblue}{#1}\begin{otherlanguage}{english}},after upper={\end{otherlanguage}}}
% Marques ✓ / ✗ (forme correcte / fautive) souvent utilisées par les modèles
\providecommand{\cross}{\textcolor{poppink}{\ensuremath{\times}}}
\providecommand{\xmark}{\cross}\providecommand{\crossmark}{\cross}\providecommand{\cmark}{\ensuremath{\checkmark}}
% Ligne de réponse / justification à compléter (inventée par les modèles)
\providecommand{\justification}[1]{\par\noindent\textit{Justification~:} \dotfill\par}
% \textipa (package tipa non chargé) : la police ne couvre pas l'API -> simple passage
\providecommand{\textipa}[1]{#1}
% Soulignement (ulem non chargé) : \uline, \ul
\providecommand{\uline}[1]{\underline{#1}}\providecommand{\ul}[1]{\underline{#1}}
% \glossaire{\item ...} inventée par les modèles : liste de glossaire
\providecommand{\glossaire}[1]{\begin{itemize}#1\end{itemize}}
% \alert (beamer) inventée par les modèles : texte mis en évidence
\providecommand{\alert}[1]{\textbf{\textcolor{poppink}{#1}}}
% variantes ASCII d'ordinaux écrites en macro
\providecommand{\iere}{\textsuperscript{re}}\providecommand{\ieme}{\textsuperscript{e}}\providecommand{\ere}{\textsuperscript{re}}\providecommand{\eme}{\textsuperscript{e}}\providecommand{\er}{\textsuperscript{er}}\providecommand{\ie}{\textsuperscript{e}}
% Ordinaux français écrits en macro par les modèles : 1\ière, 3\ième
\newcommand{\ière}{\textsuperscript{re}}\newcommand{\ième}{\textsuperscript{e}}
% \exemple (inventée par les modèles) : étiquette « Exemple : »
\providecommand{\exemple}{\textbf{Exemple~:}\ }
% \square utilisable aussi hors mode math (cases à cocher)
\AtBeginDocument{\let\squareMath\square\renewcommand{\square}{\ensuremath{\squareMath}}}
% Mot ou phrase en anglais à l'intérieur d'un paragraphe français
\newcommand{\eng}[1]{\ifmmode\text{\textit{#1}}\else\begingroup\selectlanguage{english}\itshape #1\endgroup\fi}
\let\esp\eng % alias toléré
% flèches d'intonation inventées par les modèles
\providecommand{\textsearrow}{}\renewcommand{\textsearrow}{\ensuremath{\searrow}}\providecommand{\textnearrow}{}\renewcommand{\textnearrow}{\ensuremath{\nearrow}}
% \fr{..} : mot français (inventé par les modèles)
\providecommand{\fr}[1]{\textit{#1}}
\newtcolorbox{exemplebox}[1][]{popbase,colframe=poporange,before upper={\popboxhead{Exemple}{poporange}{#1}}}
\newtcolorbox{definition}[1][]{popbase,colframe=popblue,before upper={\popboxhead{Définition}{popblue}{#1}}}
\newtcolorbox{propriete}[1][]{popbase,colframe=poppurple,before upper={\popboxhead{Propriété}{poppurple}{#1}}}
\newtcolorbox{methode}[1][]{popbase,colframe=popgold,before upper={\popboxhead{Méthode}{popgold}{#1}}}
\newtcolorbox{attention}[1][]{popbase,colframe=poppink,before upper={\popboxhead{Attention}{poppink}{#1}}}
\newtcolorbox{experience}[1][]{popbase,colframe=popblue,colback=popblueL,before upper={\popboxhead{Expérience}{popblue}{#1}}}
% « Le savais-tu ? » : carte violette pleine (contexte, culture, Cameroun)
\newtcolorbox{savaistu}[1][]{popbase,colback=poppurple,colframe=popink,
  fontupper=\popsemi\fontsize{10.4}{14.2}\selectfont\color{white},
  before upper={\poppill{Le savais-tu\,?}{white}{poppurple}\ifx\relax#1\relax\else\hspace{6pt}{\popxbold\color{white}#1}\fi\par\vspace{7pt}}}
% Exercice résolu : énoncé en haut, \tcblower puis la solution sur fond crème
\newcounter{popexo}\newcounter{popexr}
\newtcolorbox{exoresolu}[1][]{popbase,colframe=poporange,colbacklower=popcream,
  segmentation style={popink,line width=1.2pt,dashed},
  before upper={\stepcounter{popexr}\popboxhead{Exercice résolu \thepopexr}{poporange}{#1}},
  before lower={\poppill{Solution}{popink}{popcream}\par\vspace{6pt}}}
% Exercice (non résolu en place) — la correction est dans la partie Corrigés
\newtcolorbox{exercice}[1][]{popbase,colframe=popink,
  before upper={\stepcounter{popexo}\popboxhead{Exercice \thepopexo}{popink}{#1}}}
% Corrigé
\newtcolorbox{corrige}[1][]{popbase,colframe=popgreen,colback=popgreenL,
  before upper={\popboxhead{Corrigé}{popgreen}{#1}}}
% Objectifs / prérequis / bilan
\newtcolorbox{objectifs}{popbase,colframe=popink,colback=poporangeL,
  before upper={\popboxhead{Objectifs de la leçon}{popink}{}}}
\newtcolorbox{prerequis}{popbase,colframe=popink,colback=popblueL,
  before upper={\popboxhead{Prérequis}{popblue}{}}}
\newtcolorbox{bilan}{popbase,colframe=popink,colback=white,boxrule=2.8pt,
  before upper={\popboxhead{Fiche bilan}{poporange}{L'essentiel à connaître pour l'examen}}}
% Figure encadrée avec légende
\newtcolorbox{popfigure}[1][]{enhanced,breakable=false,colback=white,colframe=popline,boxrule=1.4pt,arc=8pt,
  left=10pt,right=10pt,top=10pt,bottom=8pt,before skip=10pt,after skip=12pt,halign=center,
  lower separated=false,halign lower=center,
  fontlower=\popsemi\fontsize{9}{11.5}\selectfont\color{popmuted},
  #1}
\newcommand{\legende}[1]{\par\vspace{6pt}{\popsemi\fontsize{9}{11.5}\selectfont\color{popmuted}#1}}

% ============================================================
% Signature OnBuch+
% ============================================================
\newcommand{\poplogoinline}[1]{{\poplogo\fontsize{#1}{#1}\selectfont\color{popink}OnBuch\color{poporange}+}}
\newcommand{\signatureonbuch}{%
  \begin{tcolorbox}[enhanced,colback=popink,colframe=popink,boxrule=2.2pt,arc=10pt,
      drop shadow=poporange,left=14pt,right=14pt,top=10pt,bottom=10pt,before skip=0pt,after skip=0pt]
    \tikz[baseline=-0.7ex]\node[circle,fill=popgreen,draw=popcream,line width=1.3pt,minimum size=22pt,inner sep=0]{\popcheck{white}};%
    \hspace{9pt}\begin{minipage}[c]{0.62\linewidth}
      {\popxbold\fontsize{12}{13}\selectfont\color{popcream}Signé\ \textbullet\ {\poplogo OnBuch\color{poporange}+}}\par\vspace{2pt}
      {\raggedright\popsemi\fontsize{8}{10.5}\selectfont\color{popline}\MakeUppercase{Équipe pédagogique OnBuch+}\\\MakeUppercase{Conforme au programme officiel MINESEC}\par}
    \end{minipage}\hfill
    {\poplogo\fontsize{22}{22}\selectfont\color{popcream}OnBuch\color{poporange}+}
  \end{tcolorbox}%
}

% ============================================================
% Page de garde
% #1 titre · #2 matière · #3 niveau · #4 module · #5 n° leçon · #6 durée ·
% #7 description / objectifs (texte ou itemize)
% ============================================================
\newcommand{\pagegarde}[7]{%
  \thispagestyle{empty}
  \newgeometry{a4paper,margin=1.7cm,top=1.6cm,bottom=1.4cm}%
  \begin{tikzpicture}[remember picture,overlay]
    \fill[poporange!18] ([xshift=-3.2cm,yshift=-3.6cm]current page.north east) circle (4.6cm);
    \fill[poppurple!14] ([xshift=2.4cm,yshift=7.6cm]current page.south west) circle (3.8cm);
  \end{tikzpicture}%
  {\poplogo\fontsize{30}{30}\selectfont\color{popink}OnBuch\color{poporange}+}\hfill
  \raisebox{0.6ex}{\poppill{Cours signé \textbullet\ Premium}{popink}{popcream}}\par
  \vspace{1pt}\popsquiggle{poppurple}\par
  \vspace{8pt}
  \begin{tcolorbox}[enhanced,colback=poporange,colframe=popink,boxrule=2.8pt,arc=13pt,
      drop shadow=popink,left=18pt,right=18pt,top=12pt,bottom=13pt,after skip=10pt]
    \poppill{#2 \textbullet\ #3}{popink}{popcream}\hspace{5pt}\poppill{#4}{white}{popink}\par\vspace{8pt}
    {\popdisplay\fontsize{14}{14}\selectfont\color{popink}LEÇON #5}\par\vspace{4pt}
    {\raggedright\hyphenpenalty=10000\exhyphenpenalty=10000\popdisplay\fontsize{25}{27}\selectfont\color{popcream}\MakeUppercase{#1}\par}
    \vspace{8pt}
    {\popxbold\fontsize{9.5}{9.5}\selectfont\color{popink}\MakeUppercase{Durée conseillée : #6}}
  \end{tcolorbox}
  \begin{tcolorbox}[enhanced,colback=white,colframe=popink,boxrule=2.2pt,arc=10pt,
      drop shadow=popink,left=16pt,right=16pt,top=10pt,bottom=10pt,after skip=8pt]
    \poppill{Dans cette leçon}{poppurple}{white}\par\vspace{5pt}
    \popsemi\fontsize{9.3}{12.6}\selectfont\color{popink}#7
  \end{tcolorbox}
  \noindent\begin{minipage}[t]{0.487\linewidth}
    \begin{tcolorbox}[enhanced,colback=popgreen,colframe=popink,boxrule=2.2pt,arc=10pt,
        drop shadow=popink,left=11pt,right=11pt,top=10pt,bottom=10pt,height=3.1cm,valign=center]
      \tcbox[colback=white,colframe=popink,boxrule=1.3pt,arc=5pt,boxsep=0pt,nobeforeafter,
        left=3pt,right=3pt,top=3pt,bottom=3pt,valign=center]{\qrcode[height=1.9cm]{https://play.google.com/store/apps/details?id=com.luvvix.onbuch&hl=fr}}%
      \hspace{8pt}\begin{minipage}[c]{0.5\linewidth}
        {\popdisplay\fontsize{11}{12.5}\selectfont\color{white}\MakeUppercase{Télécharge OnBuch}}\par\vspace{4pt}
        {\raggedright\popsemi\fontsize{8.4}{11}\selectfont\color{white}Gratuit \textbullet\ Google Play\\Tuteur IA, annales, concours, résultats\par}
      \end{minipage}
    \end{tcolorbox}
  \end{minipage}\hfill
  \begin{minipage}[t]{0.487\linewidth}
    \begin{tcolorbox}[enhanced,colback=poppurple,colframe=popink,boxrule=2.2pt,arc=10pt,
        drop shadow=popink,left=11pt,right=11pt,top=10pt,bottom=10pt,height=3.1cm,valign=center]
      \tikz[baseline=-0.7ex]\node[circle,fill=white,draw=popink,line width=1.3pt,minimum size=1.6cm,inner sep=0]{\popdisplay\fontsize{24}{24}\selectfont\color{poppurple}+};%
      \hspace{8pt}\begin{minipage}[c]{0.6\linewidth}
        {\popdisplay\fontsize{11}{12.5}\selectfont\color{white}\MakeUppercase{Passe à OnBuch+}}\par\vspace{4pt}
        {\raggedright\popsemi\fontsize{8.4}{11}\selectfont\color{white}Tous les cours signés, Tuteur IA illimité, fiches PDF — onglet OnBuch+ de l'app\par}
      \end{minipage}
    \end{tcolorbox}
  \end{minipage}
  \vspace{8pt}
  \popcontenu
  \vfill
  \signatureonbuch
  \restoregeometry
  \newpage
}

% Rangée « Ce cours contient » (page de garde)
\newcommand{\poptile}[3]{% #1 couleur #2 gros texte #3 légende
  \begin{tcolorbox}[enhanced,colback=#1,colframe=popink,boxrule=2pt,arc=9pt,shadow={2.5pt}{-2.5pt}{0pt}{popink},
      left=6pt,right=6pt,top=9pt,bottom=9pt,halign=center,valign=center,height=1.75cm,width=\linewidth,nobeforeafter]
    {\popdisplay\fontsize{12.5}{13}\selectfont\color{white}\MakeUppercase{#2}}\par\vspace{3pt}
    {\centering\popsemi\fontsize{7.8}{9.5}\selectfont\color{white}#3\par}
  \end{tcolorbox}}
\newcommand{\popcontenu}{%
  \par\noindent{\popxboldls\fontsize{7.5}{7.5}\selectfont\color{popmuted}CE COURS SIGNÉ CONTIENT}\par\vspace{7pt}
  \noindent\begin{minipage}[t]{0.235\linewidth}\poptile{popblue}{Cours}{complet, illustré, pas à pas}\end{minipage}\hfill
  \begin{minipage}[t]{0.235\linewidth}\poptile{poporange}{Méthodes}{et exercices résolus}\end{minipage}\hfill
  \begin{minipage}[t]{0.235\linewidth}\poptile{poppink}{Exercices}{type examen + corrigés}\end{minipage}\hfill
  \begin{minipage}[t]{0.235\linewidth}\poptile{popgreen}{Fiche bilan}{l'essentiel à retenir}\end{minipage}\par}

% Sommaire
\newtcolorbox{sommaire}{enhanced,colback=white,colframe=popink,boxrule=2.2pt,arc=10pt,
  drop shadow=popink,left=18pt,right=18pt,top=13pt,bottom=13pt,before skip=6pt,after skip=20pt,
  before upper={\poppill{Sommaire}{poppurple}{white}\par\vspace{9pt}\popsemi\fontsize{10.6}{14.5}\selectfont\color{popink}}}

% Signature de fin de cours — poussée tout en bas de la dernière page
\newcommand{\findecours}{%
  \par\vfill\nopagebreak
  \begin{center}\popsquiggle{poporange}\end{center}\vspace{4pt}
  \signatureonbuch
}
\AtBeginDocument{\def\llbracket{\mathopen{[\mkern-3.5mu[}}\def\rrbracket{\mathclose{]\mkern-3.5mu]}}} % intervalles d'entiers (glyphe ⟦ absent de la police)
% Glyphes absents de Fira Math : redéfinis à partir de glyphes disponibles
\AtBeginDocument{%
  \def\setminus{\mathbin{\backslash}}\let\smallsetminus\setminus
  \def\vdots{\mathord{\vbox{\baselineskip3.2pt\lineskiplimit0pt\kern4pt\hbox{.}\hbox{.}\hbox{.}}}}%
  \def\ddots{\mathinner{\mkern1mu\raise6.5pt\hbox{.}\mkern2mu\raise3.6pt\hbox{.}\mkern2mu\raise0.7pt\hbox{.}\mkern1mu}}%
  \def\triangle{\mathord{\scalebox{0.9}{$\symup{\Delta}$}}}%
  \def\nabla{\mathord{\raisebox{\depth}{\scalebox{1}[-1]{$\symup{\Delta}$}}}}%
  \def\checkmark{\popcheck{popdark}}%
  \def\star{\mathbin{\ast}}\def\bigstar{\mathord{\ast}}%
  \def\diamond{\mathbin{\scalebox{0.6}{$\Diamond$}}}%
  \def\bigcup{\mathop{\mathchoice{\raisebox{-0.3em}{\scalebox{1.7}{$\cup$}}}{\scalebox{1.25}{$\cup$}}{\cup}{\cup}}\displaylimits}%
  \def\bigcap{\mathop{\mathchoice{\raisebox{-0.3em}{\scalebox{1.7}{$\cap$}}}{\scalebox{1.25}{$\cap$}}{\cap}{\cap}}\displaylimits}%
  \def\lesseqqgtr{\mathrel{\lessgtr}}\def\bigoplus{\mathop{\oplus}\displaylimits}%
}
\providecommand{\ssi}{si et seulement si} % abréviation fréquente
\AtBeginDocument{\providecommand{\wideparen}{\overparen}} % arc de cercle

%% FIGFIX-BEGIN v5 — correctifs globaux des figures (docs/onbuch-generation/tools/figfix)
\usepackage{adjustbox}\usepackage{etoolbox}\tcbuselibrary{raster}
% toute figure TikZ/pgfplots plus large que la ligne est réduite à la largeur disponible
\BeforeBeginEnvironment{tikzpicture}{\begin{adjustbox}{max width=\linewidth}}
\AfterEndEnvironment{tikzpicture}{\end{adjustbox}}
% légendes pgfplots lisibles par-dessus les courbes
\pgfplotsset{every axis/.append style={legend style={fill=white,fill opacity=0.92,text opacity=1,draw=black!15,inner sep=3pt}}}
% étiquettes d'arêtes (syntaxe edge["..."]) avec fond blanc
\tikzset{every edge quotes/.append style={fill=white,inner sep=1.2pt}}
%% FIGFIX-END
\begin{document}

\pagegarde{Lesson 35 — Vocabulary: Words and expressions related to garbage collection, disposal, and recycling services}{Anglais}{Tle A}{Chapter 3 : Environment, Well-being and Health}{EN35}{2 h}{Cette leçon de vocabulaire présente le lexique thématique de la collecte, de l'élimination et du recyclage des déchets : champs lexicaux, collocations, expressions idiomatiques, faux amis et familles de mots. Les élèves apprennent à employer ce lexique avec précision à l'oral et à l'écrit, en contexte camerounais et anglophone.
\vspace{4pt}
\begin{multicols}{2}\small
\begin{itemize}
\item À la fin de cette leçon, je sais nommer les acteurs, lieux et équipements de la collecte des déchets (dustbin, bin lorry, landfill...)
\item je sais distinguer les différents types de déchets en anglais (rubbish, garbage, litter, waste, refuse)
\item je sais utiliser les collocations essentielles du thème (take out the rubbish, sort waste, dump illegally...)
\item je sais employer le vocabulaire du recyclage et de la gestion des déchets (recycle, reuse, compost, dispose of...)
\item je sais éviter les faux amis et les calques du français (décharge ≠ 'discharge', ordures ≠ 'orders'...)
\item je sais former des mots de la même famille (collect/collection/collector, recycle/recyclable/recycling...)
\end{itemize}\end{multicols}}

\begin{sommaire}
\begin{multicols}{2}
\begin{enumerate}
\item Waste and its types: the core vocabulary
\item Collection services: people, vehicles and places
\item Disposal, recycling and the 3 Rs
\item Collocations and idiomatic expressions
\item Word families, spelling and false friends
\item Using the vocabulary in context
\item Activité d'intégration
\item Méthodes et exercices résolus
\item Exercices
\item Corrigés des exercices
\item Fiche bilan
\end{enumerate}
\end{multicols}
\end{sommaire}

\begin{objectifs}
\begin{itemize}
\item À la fin de cette leçon, je sais nommer les acteurs, lieux et équipements de la collecte des déchets (dustbin, bin lorry, landfill...)
\item je sais distinguer les différents types de déchets en anglais (rubbish, garbage, litter, waste, refuse)
\item je sais utiliser les collocations essentielles du thème (take out the rubbish, sort waste, dump illegally...)
\item je sais employer le vocabulaire du recyclage et de la gestion des déchets (recycle, reuse, compost, dispose of...)
\item je sais éviter les faux amis et les calques du français (décharge ≠ 'discharge', ordures ≠ 'orders'...)
\item je sais former des mots de la même famille (collect/collection/collector, recycle/recyclable/recycling...)
\item je sais prononcer et orthographier correctement les mots-clés du thème
\item je sais réemployer ce lexique dans des phrases, des dialogues et un court texte argumentatif
\end{itemize}
\end{objectifs}

\begin{prerequis}
\begin{itemize}
\item Vocabulaire général de l'environnement vu en Première (pollution, environment, protect)
\item Les verbes suivis de la préposition 'of' ou 'to' (dispose of, consist of)
\item La formation des noms en -tion, -ment, -er (dérivation de base)
\item Les modaux must, should, have to pour exprimer l'obligation (utile pour parler des règles de tri)
\item Le present simple pour les faits et habitudes
\end{itemize}
\end{prerequis}

\newpage

\coursec{Waste and its types: the core vocabulary}

Every morning in Yaoundé and Douala, HYSACAM trucks collect tons of household waste, yet many bins overflow before collection day. In London or New York, citizens sort their rubbish into colour-coded bins and can be fined for mixing recyclables with general waste. What do we call the people who empty our bins? What is the difference between \eng{rubbish}, \eng{garbage} and \eng{litter}? In this lesson, you will learn the exact English words and expressions to talk about waste collection, disposal and recycling — and avoid the traps of French-English false friends.

\textbf{Problématique :} How can we name, classify and talk about waste accurately in English, choosing the right word for the right context (British or American English, formal or everyday language)?

\courssub{Discovering the vocabulary through a text}

Read the following short report about waste management in a Cameroonian town. Underline all the words related to waste, then try to guess their meaning from the context.

\begin{textebox}[A report from Buea — original text]
In Buea, the council provides wheelie bins for household waste, but litter is still dropped in the streets. The town has no proper landfill, so rubbish is often dumped in open sites.

Every Monday and Thursday, the collection truck passes through the neighbourhoods. Families put their bin bags outside their gates the night before. Most of the waste is food waste — banana peels, cassava leaves and leftover fufu — but there is also a growing amount of plastic waste: empty water bottles, plastic bags and broken buckets. Near the market, traders throw away rotten fruit and vegetables, and the smell attracts flies.

The situation becomes dangerous when hazardous waste is mixed with ordinary rubbish. Last month, a worker was injured by broken glass and rusty nails hidden in a bin liner. The hospital, for its part, must deal with medical waste such as used syringes and bandages, which should never be dumped with household refuse.

“We need more litter bins in the streets,” says Mrs Efange, a shopkeeper. “People drop litter because there is nowhere to put it. And the council should empty the skips near the market more often — they overflow after two days.”

The council has promised to build a proper disposal site outside the town and to start a recycling programme for plastic and electronic waste. Until then, the people of Buea continue to live with overflowing bins and dirty streets.
\end{textebox}

\textbf{Glossaire :}
\begin{itemize}
\item \cle{council} (n.) : la municipalité, le conseil municipal
\item \cle{landfill} (n.) : une décharge (contrôlée), un site d'enfouissement
\item \cle{to dump} (v.) : déverser, déposer sauvagement
\item \cle{peel} (n.) : épluchure, pelure
\item \cle{leftover} (n.) : reste (de nourriture)
\item \cle{rusty} (adj.) : rouillé
\item \cle{syringe} (n.) : seringue
\item \cle{to overflow} (v.) : déborder
\end{itemize}

\textbf{Comprehension questions :}
\begin{enumerate}
\item What two types of waste are mentioned as the most common in Buea?
\item Why was the worker injured last month?
\item According to Mrs Efange, why do people drop litter?
\item What has the council promised to do?
\end{enumerate}

\courssub{The general words: waste, rubbish, garbage, trash, refuse, litter, junk}

English has several words for « les déchets », and they are \emph{not} interchangeable. Each word has its own register (everyday, official, American, British) and its own meaning. Observe the examples first:

\eng{The factory produces a lot of industrial waste.} « L'usine produit beaucoup de déchets industriels. »

\eng{Don't forget to take out the rubbish!} « N'oublie pas de sortir les ordures ! »

\eng{Someone dropped litter on the pavement.} « Quelqu'un a jeté des détritus sur le trottoir. »

\eng{My garage is full of old junk.} « Mon garage est plein de vieilleries. »

\begin{definition}[The seven key words]
\begin{itemize}
\item \cle{waste} (n., uncountable) : the general and rather formal word for « déchets ». Used in official and scientific contexts: \eng{waste management}, \eng{nuclear waste}. Prononciation : « ouèiste ».
\item \cle{rubbish} (n., uncountable, GB) : household waste, everyday British word. Prononciation : « reubiche ». \eng{Put the rubbish in the bin.}
\item \cle{garbage} / \cle{trash} (n., uncountable, US) : the American equivalents of \eng{rubbish}. Prononciation : « gâabidj » / « trache ». \eng{Take out the trash, please.}
\item \cle{refuse} (n., uncountable) : the official, administrative word, used by councils and in regulations. Prononciation : « ré-fiouze ». \eng{refuse collection} = « la collecte des ordures ».
\item \cle{litter} (n., uncountable) : waste dropped in public places (streets, parks, beaches). Prononciation : « liteur ». \emph{Never} use it for household waste.
\item \cle{junk} (n., uncountable) : old, useless objects, « vieilleries, bric-à-brac ». \eng{junk food} = « la malbouffe ».
\end{itemize}
All these nouns are \textbf{uncountable}: no plural, no \eng{a/an}. Say \eng{some rubbish}, \eng{a piece of litter}, never \eng{a rubbish} or \eng{litters}.
\end{definition}

\begin{attention}[Faux ami : « les ordures » n'est pas « orders » !]
Ne traduisez jamais « les ordures » par \eng{orders} ! \eng{Orders} signifie « les commandes » ou « les ordres » (militaires). On dit \eng{household waste}, \eng{rubbish} (GB) ou \eng{garbage} (US).

Autre piège : \eng{waste} se prononce « ouèiste » (comme \eng{taste}, \eng{haste}), et non « ouaste ». Et attention au verbe \eng{to waste} = « gaspiller » : \eng{Don't waste water!} « Ne gaspille pas l'eau ! »
\end{attention}

\begin{exemplebox}[The words in context]
\eng{Please put your litter in the bin.} « Jetez vos détritus dans la poubelle, s'il vous plaît. »

\eng{The bins are emptied twice a week.} « Les poubelles sont vidées deux fois par semaine. »

\eng{The city spends millions on refuse collection.} « La ville dépense des millions pour la collecte des ordures. »

\eng{He cleared all the junk out of his room.} « Il a enlevé toutes les vieilleries de sa chambre. »

\eng{Plastic waste pollutes our rivers and oceans.} « Les déchets plastiques polluent nos rivières et nos océans. »
\end{exemplebox}

\courssub{Categories of waste}

Waste can be classified by its origin or its nature. Learn these compound nouns — they are extremely frequent in news reports and exams:

\begin{center}
\begin{tabularx}{\linewidth}{|X|X|X|}
\hline
\rowcolor{popblueL} English & Français & Exemple \\
\hline
\cle{household waste} & déchets ménagers & \eng{Household waste is collected twice a week.} \\
\hline
\cle{food waste} / \cle{kitchen waste} & déchets alimentaires & \eng{Food waste can be turned into compost.} \\
\hline
\cle{plastic waste} & déchets plastiques & \eng{Plastic waste takes centuries to disappear.} \\
\hline
\cle{electronic waste} (\cle{e-waste}) & déchets électroniques & \eng{Old phones and computers are e-waste.} \\
\hline
\cle{hazardous} / \cle{toxic waste} & déchets dangereux / toxiques & \eng{Batteries are hazardous waste.} \\
\hline
\cle{medical waste} & déchets médicaux & \eng{Hospitals must burn their medical waste.} \\
\hline
\cle{industrial waste} & déchets industriels & \eng{Factories produce industrial waste.} \\
\hline
\cle{green} / \cle{garden waste} & déchets verts & \eng{Grass and leaves are garden waste.} \\
\hline
\end{tabularx}
\end{center}

\begin{propriete}[Formation : noun + waste]
Ces mots composés suivent le schéma \textbf{adjectif ou nom + waste} : \eng{household waste}, \eng{toxic waste}. Le premier élément précise l'origine ou la nature du déchet. En français, on utilise souvent « déchets + de/à » : \eng{kitchen waste} = « déchets de cuisine ». Tous restent indénombrables : \eng{E-waste \textbf{is} a growing problem.} (singulier !)
\end{propriete}

\courssub{Containers: bins, cans, skips and bags}

\begin{definition}[The containers]
\begin{itemize}
\item \cle{dustbin} / \cle{wheelie bin} (GB) : la poubelle (domestique) ; la \eng{wheelie bin} est la poubelle à roulettes. Prononciation de \eng{bin} : « bine ».
\item \cle{trash can} / \cle{garbage can} (US) : la poubelle (américaine).
\item \cle{litter bin} : la poubelle publique (dans la rue, les parcs).
\item \cle{skip} (GB) : la grande benne (pour les gravats, les chantiers).
\item \cle{bin bag} / \cle{bin liner} : le sac poubelle.
\item \cle{dumpster} (US) : la benne à ordures (américaine).
\end{itemize}
\end{definition}

\begin{exemplebox}[Containers in context]
\eng{Throw your empty bottle in the litter bin.} « Jetez votre bouteille vide dans la poubelle publique. »

\eng{They hired a skip to get rid of the rubble.} « Ils ont loué une benne pour se débarrasser des gravats. »

\eng{Line the bin with a bin liner.} « Mettez un sac poubelle dans la poubelle. »
\end{exemplebox}

\courssub{British English vs American English}

\begin{center}
\begin{tabularx}{\linewidth}{|X|X|X|}
\hline
\rowcolor{popblueL} British English & American English & Français \\
\hline
rubbish & garbage / trash & les ordures, les déchets \\
\hline
dustbin / (wheelie) bin & trash can / garbage can & la poubelle \\
\hline
dustman / binman & garbage collector / garbage man & l'éboueur \\
\hline
rubbish dump / tip & garbage dump / landfill & la décharge \\
\hline
skip & dumpster & la benne \\
\hline
lorry (refuse lorry) & garbage truck & le camion à ordures \\
\hline
\end{tabularx}
\end{center}

\begin{attention}[Mélanger GB et US]
Évitez de mélanger les deux variétés dans une même production : ne dites pas \eng{Put the garbage in the dustbin}. Choisissez un système : \eng{Put the rubbish in the dustbin} (GB) ou \eng{Put the garbage in the trash can} (US). En revanche, \eng{litter} et \eng{waste} s'emploient dans les deux variétés.
\end{attention}

\courssub{Mind map: organising the vocabulary}

\begin{popfigure}
\begin{center}
\begin{center}\tcbox[colback=popink,colframe=popink,coltext=white,arc=9pt,boxsep=4pt,left=6pt,right=6pt]{\bfseries\small WASTE}\end{center}
\vspace{-2pt}
\begin{tcbraster}[raster columns=2,raster equal height=rows,raster column skip=6pt,raster row skip=6pt,enhanced,blankest]
\begin{tcolorbox}[enhanced,colback=white,colframe=poporange!80,boxrule=0.9pt,arc=3pt,title={Types},fonttitle=\bfseries\scriptsize,coltitle=white,colbacktitle=poporange!80,left=4pt,right=4pt,top=2pt,bottom=2pt,boxsep=1pt,toptitle=1.5pt,bottomtitle=1.5pt,halign=left]
\begin{itemize}[nosep,leftmargin=*,itemsep=1pt,font=\scriptsize]
\item household food plastic
\item e-waste hazardous medical
\end{itemize}
\end{tcolorbox}
\begin{tcolorbox}[enhanced,colback=white,colframe=popgreen!70,boxrule=0.9pt,arc=3pt,title={Containers},fonttitle=\bfseries\scriptsize,coltitle=white,colbacktitle=popgreen!70,left=4pt,right=4pt,top=2pt,bottom=2pt,boxsep=1pt,toptitle=1.5pt,bottomtitle=1.5pt,halign=left]
\begin{itemize}[nosep,leftmargin=*,itemsep=1pt,font=\scriptsize]
\item bin, trash can
\item skip, bin bag
\end{itemize}
\end{tcolorbox}
\begin{tcolorbox}[enhanced,colback=white,colframe=poppurple!70,boxrule=0.9pt,arc=3pt,title={People},fonttitle=\bfseries\scriptsize,coltitle=white,colbacktitle=poppurple!70,left=4pt,right=4pt,top=2pt,bottom=2pt,boxsep=1pt,toptitle=1.5pt,bottomtitle=1.5pt,halign=left]
\begin{itemize}[nosep,leftmargin=*,itemsep=1pt,font=\scriptsize]
\item dustman (GB)
\item garbage collector (US)
\end{itemize}
\end{tcolorbox}
\begin{tcolorbox}[enhanced,colback=white,colframe=popgold!90,boxrule=0.9pt,arc=3pt,title={Places},fonttitle=\bfseries\scriptsize,coltitle=white,colbacktitle=popgold!90,left=4pt,right=4pt,top=2pt,bottom=2pt,boxsep=1pt,toptitle=1.5pt,bottomtitle=1.5pt,halign=left]
\begin{itemize}[nosep,leftmargin=*,itemsep=1pt,font=\scriptsize]
\item landfill, dump
\item recycling centre
\end{itemize}
\end{tcolorbox}
\end{tcbraster}

\end{center}
\legende{Carte mentale du vocabulaire des déchets : quatre branches à compléter au fil de la leçon.}
\end{popfigure}

\begin{methode}[How to learn thematic vocabulary]
\begin{enumerate}
\item \textbf{Group} words by theme (types, containers, people, places) — never learn isolated lists.
\item \textbf{Note} the variety (GB/US) and the register (formal/informal) next to each word.
\item \textbf{Learn collocations}, not single words: \eng{drop litter}, \eng{empty a bin}, \eng{dump rubbish}, \eng{sort waste}.
\item \textbf{Test yourself} with translation both ways: French to English, then English to French.
\item \textbf{Reuse} the words in your own sentences about your town (Yaoundé, Douala, Buea...).
\end{enumerate}
\end{methode}

\begin{exoresolu}[Choosing the right word]
Énoncé — Complete each sentence with the correct word: \eng{litter}, \eng{rubbish}, \eng{junk}, \eng{waste}, \eng{bin}.
\begin{enumerate}
\item Please don't drop \dots\ in the school playground.
\item Our \dots\ is collected every Tuesday morning.
\item Nuclear \dots\ must be stored very carefully.
\item The attic is full of old \dots\ we never use.
\item Throw that empty can in the \dots\ !
\end{enumerate}
\tcblower
Solution détaillée :
\begin{enumerate}
\item \textbf{litter} — il s'agit de détritus jetés dans un lieu public (la cour de récréation).
\item \textbf{rubbish} — déchets ménagers ramassés chaque semaine (contexte britannique ; \eng{garbage} serait accepté en contexte américain).
\item \textbf{waste} — seul \eng{waste} convient dans le composé officiel \eng{nuclear waste} (registre soutenu).
\item \textbf{junk} — de vieux objets sans valeur entreposés dans un grenier.
\item \textbf{bin} — le contenant ; \eng{throw ... in the bin} (GB) / \eng{in the trash can} (US).
\end{enumerate}
\end{exoresolu}

\begin{savaistu}[Dropping litter costs money!]
Au Royaume-Uni, jeter des détritus par terre (\eng{dropping litter}) est une infraction punie d'une amende (\eng{a fine}) pouvant atteindre £150. Les panneaux \eng{“No littering — £150 fine”} sont visibles dans les rues, les parcs et les gares. Le mot \eng{litter} vient du vieux français \emph{litière} — la paille sur laquelle on dormait au Moyen Âge, qu'on laissait par terre !
\end{savaistu}

\begin{exercice}[Your turn — translate into English]
Traduisez les phrases suivantes en anglais, en choisissant le mot juste :
\begin{enumerate}
\item Les poubelles débordent avant le jour de la collecte.
\item Ne jetez pas vos détritus dans la rue !
\item Les déchets électroniques contiennent des métaux dangereux.
\item L'éboueur vide notre poubelle tous les lundis.
\end{enumerate}
\end{exercice}

\begin{corrige}[Exercice « Your turn »]
\begin{enumerate}
\item \eng{The bins overflow before collection day.}
\item \eng{Don't drop your litter in the street!}
\item \eng{Electronic waste contains dangerous metals.} (singulier : \eng{contains}, car \eng{waste} est indénombrable)
\item \eng{The dustman empties our bin every Monday.} (GB) / \eng{The garbage collector empties our trash can every Monday.} (US)
\end{enumerate}
\end{corrige}

\begin{aretenir}[Section 1 — Key points]
\begin{itemize}
\item \eng{waste} = terme général et soutenu ; \eng{rubbish} (GB) = \eng{garbage/trash} (US) ; \eng{refuse} = terme administratif ; \eng{litter} = détritus dans les lieux publics ; \eng{junk} = vieilleries.
\item Tous ces noms sont \textbf{indénombrables} : jamais de pluriel ni de \eng{a/an}.
\item Catégories : \eng{household, food, plastic, electronic (e-waste), hazardous, medical, industrial waste}.
\item Contenants : \eng{dustbin/wheelie bin} (GB), \eng{trash can} (US), \eng{litter bin} (publique), \eng{skip} (benne), \eng{bin bag/bin liner} (sac poubelle).
\item Faux ami : « les ordures » n'est jamais \eng{orders} !
\end{itemize}
\end{aretenir}

\coursec{Collection services: people, vehicles and places}

Dans la section précédente, nous avons appris à nommer les déchets eux-mêmes (\eng{waste, rubbish, garbage, litter}) et à classer leurs types (\eng{household waste, organic waste, recyclable waste}). Mais une fois la poubelle remplie, que se passe-t-il ? Qui ramasse les ordures, avec quels véhicules, et où sont-elles emmenées ? C'est tout le vocabulaire des \emph{services de collecte} que nous allons maintenant construire : les \cle{acteurs} (les éboueurs), les \cle{véhicules et équipements} (le camion à ordures, les gants), les \cle{lieux} (la décharge, le centre de recyclage) et les \cle{services} (la collecte porte-à-porte).

\courssub{Observation : a text about waste collection in Douala}

Lisons d'abord ce court texte, qui décrit une scène quotidienne dans la capitale économique du Cameroun. Soulignez mentalement tous les mots liés à la collecte des déchets. Le texte est rédigé en \textbf{anglais britannique}, variété de référence au Cameroun.

\begin{textebox}[Waste collection in Douala]
In Douala, HYSACAM bin lorries collect household waste door to door. The refuse collectors wear gloves and load the bins onto the lorry. The waste is then taken to a landfill on the outskirts of the city.

Early in the morning, before the traffic becomes heavy, the crew starts its round in the residential neighbourhoods. One worker drives the truck while two others walk behind it, emptying the dustbins that families have left in front of their gates. The work is hard and sometimes dangerous, so the collectors wear protective gloves, boots and bright jackets. In some districts, street sweepers clean the roads and the markets with brooms and shovels, and they carry the rubbish away in wheelbarrows.

Every household pays refuse collection fees to the waste management company. Thanks to this service, the streets stay cleaner and diseases are less common. However, not all the waste goes to the landfill: scavengers often search the dump for plastic bottles, scrap metal and old electrical appliances, which they sell to small recycling businesses. In this way, informal workers also play a role in the cleansing system of the city.
\end{textebox}

\textbf{Glossaire}

\begin{center}
\begin{tabularx}{\linewidth}{|l|l|X|}
\hline
\rowcolor{popblueL} English & Nature & Français \\
\hline
bin lorry & nom (GB) & camion à ordures \\
\hline
door to door & expression & porte-à-porte \\
\hline
refuse collector & nom (GB) & éboueur \\
\hline
bin / dustbin & nom & poubelle \\
\hline
landfill & nom & décharge contrôlée, site d'enfouissement \\
\hline
outskirts & nom (pluriel) & banlieue, périphérie \\
\hline
crew & nom & équipe \\
\hline
street sweeper & nom & balayeur de rue \\
\hline
wheelbarrow & nom & brouette \\
\hline
refuse collection fees & nom & redevance d'enlèvement des ordures \\
\hline
waste management company & nom & société de gestion des déchets \\
\hline
scavenger & nom & récupérateur (informel) \\
\hline
scrap metal & nom & ferraille \\
\hline
\end{tabularx}
\end{center}

\textbf{Comprehension questions}

\begin{enumerate}
\item Who collects household waste in Douala?
\item What do the refuse collectors wear, and why?
\item Where is the waste taken after collection?
\item What do scavengers look for at the dump?
\item How do households pay for the service?
\end{enumerate}

\courssub{The people: naming the workers}

\begin{definition}[Les acteurs de la collecte]
\begin{itemize}
\item \eng{dustman} (GB, pluriel \eng{dustmen}) et \eng{refuse collector} (GB, terme officiel) : l'éboueur. Prononciation de \eng{collector} : « keu-lèk-teur » (accent sur \emph{lec}).
\item \eng{garbage collector} ou \eng{garbage man} (US) : l'éboueur (variété américaine).
\item \eng{street cleaner} / \eng{street sweeper} : le balayeur de rue.
\item \eng{scavenger} : le récupérateur informel, celui qui trie la décharge pour revendre des matériaux (mot très utile pour décrire le contexte africain).
\item \eng{waste management company} : la société de gestion des déchets (ex. HYSACAM au Cameroun).
\end{itemize}
\end{definition}

\begin{exemplebox}[Exemples traduits]
\eng{The dustmen come on Mondays and Thursdays.} = « Les éboueurs passent les lundis et jeudis. »

\eng{Refuse collectors in Yaoundé start work very early in the morning.} = « Les éboueurs de Yaoundé commencent le travail très tôt le matin. »

\eng{Scavengers recover plastic and metal from the dump.} = « Les récupérateurs récupèrent le plastique et le métal dans la décharge. »

\eng{HYSACAM is a waste management company.} = « HYSACAM est une société de gestion des déchets. »
\end{exemplebox}

\courssub{Vehicles, equipment and places}

\begin{definition}[Véhicules et équipements]
\begin{itemize}
\item \eng{dustcart} / \eng{bin lorry} (GB) = \eng{garbage truck} (US) : le camion à ordures.
\item \eng{wheelbarrow} : la brouette ; \eng{shovel} : la pelle ; \eng{broom} : le balai.
\item \eng{protective gloves} : les gants de protection ; \eng{boots} : les bottes ; \eng{high-visibility jacket} : gilet haute visibilité.
\end{itemize}
\end{definition}

\begin{definition}[Landfill]
A \cle{landfill} (ou \eng{landfill site}) is a large area of land where waste is buried. Prononciation : « land-file » (accent sur \emph{land}). C'est le terme précis et moderne pour « décharge contrôlée / site d'enfouissement », plus précis que \eng{dump}, qui désigne une décharge sauvage ou ancienne.
\end{definition}

\begin{definition}[Les lieux de traitement]
\begin{itemize}
\item \eng{dump} / \eng{rubbish dump} : la décharge (souvent sauvage ou terme familier).
\item \eng{tip} (GB) : la déchetterie (lieu où les particuliers apportent leurs encombrants).
\item \eng{recycling centre} : le centre de recyclage.
\item \eng{incinerator} : l'incinérateur (prononciation : « in-si-neu-réi-teur », accent sur \emph{neu}).
\item \eng{composting site} : le site de compostage.
\end{itemize}
\end{definition}

\begin{attention}[Faux ami : « une décharge » n'est pas « a discharge »]
« Une décharge » (d'ordures) ne se dit PAS \eng{a discharge}. Le mot anglais \eng{discharge} signifie « décharge électrique », « démission / libération » ou « écoulement ». On dit \eng{a dump} ou \eng{a landfill}.

Faux : \eng{The rubbish is taken to a discharge.}

Correct : \eng{The rubbish is taken to a landfill.} = « Les ordures sont emmenées à la décharge. »
\end{attention}

\begin{exemplebox}[Exemples traduits]
\eng{Old fridges are taken to the recycling centre.} = « Les vieux réfrigérateurs sont amenés au centre de recyclage. »

\eng{The bin lorry empties the dustbins twice a week.} = « Le camion à ordures vide les poubelles deux fois par semaine. »

\eng{Street sweepers use brooms and wheelbarrows to clean the market.} = « Les balayeurs utilisent des balais et des brouettes pour nettoyer le marché. »
\end{exemplebox}

\courssub{The services and the journey of waste}

\begin{definition}[Les services de collecte]
\begin{itemize}
\item \eng{waste collection service} : le service d'enlèvement des ordures.
\item \eng{door-to-door collection} : la collecte porte-à-porte.
\item \eng{kerbside collection} (GB) / \eng{curbside collection} (US) : la collecte en bordure de trottoir.
\item \eng{public cleansing department} : le service municipal de propreté.
\item \eng{refuse collection fees} / \eng{garbage collection fees} : la redevance d'enlèvement des ordures (contexte camerounais : redevance payée à HYSACAM).
\end{itemize}
\end{definition}

\begin{popfigure}
\begin{center}
\begin{tikzpicture}[node distance=0.55cm and 0.9cm,
box/.style={draw=popblue, fill=popblueL, rounded corners, align=center, minimum height=0.9cm, font=\small, text width=2.2cm},
arr/.style={-{Stealth[length=3mm]}, thick, poporange}]
\node[box, align=center] (bin){Household\\bin};
\node[box, right=of bin, align=center] (truck){Collection\\truck};
\node[box, right=of truck, align=center] (sort){Sorting\\facility};
\node[box, above right=0.35cm and 0.9cm of sort] (land) {Landfill};
\node[box, right=of sort, align=center] (rec){Recycling\\centre};
\node[box, below right=0.35cm and 0.9cm of sort] (inc) {Incinerator};
\draw[arr] (bin) -- (truck);
\draw[arr] (truck) -- (sort);
\draw[arr] (sort) -- (land);
\draw[arr] (sort) -- (rec);
\draw[arr] (sort) -- (inc);
\end{tikzpicture}
\end{center}
\legende{The journey of waste: from the household bin to its final destination.}
\end{popfigure}

\courssub{Word families and derivation (Mots de la même famille)}

\begin{aretenir}[Dérivation nominale et verbale]
Maîtriser les familles de mots permet d'enrichir son vocabulaire et de varier ses productions (écrites/orales). Attention à l'orthographe et à l'accentuation.
\begin{center}
\begin{tabularx}{\linewidth}{|l|l|l|X|}
\hline
\rowcolor{popblueL} Verbe & Nom (action/processus) & Nom (agent/lieu) & Adjectif / Participe \\
\hline
\eng{collect} & \eng{collection} & \eng{collector} & \eng{collective} \\
\hline
\eng{recycle} & \eng{recycling} & \eng{recycler} & \eng{recyclable} \\
\hline
\eng{dispose (of)} & \eng{disposal} & -- & \eng{disposable} \\
\hline
\eng{manage} & \eng{management} & \eng{manager} & \eng{manageable} \\
\hline
\eng{incinerate} & \eng{incineration} & \eng{incinerator} & \eng{incinerated} \\
\hline
\eng{compost} & \eng{composting} & \eng{compost} & \eng{compostable} \\
\hline
\eng{pollute} & \eng{pollution} & \eng{polluter} & \eng{polluted} \\
\hline
\end{tabularx}
\end{center}
\end{aretenir}

\begin{attention}[Pièges d'orthographe et de prononciation dans les familles]
\begin{itemize}
\item \eng{collect} $\to$ \eng{collection} (double \textbf{l}, accent sur \emph{lec} : « keu-lèk-chn »).
\item \eng{recycle} $\to$ \eng{recycling} (conserver le \textbf{y}, pas de \textbf{i} : \emph{re-sai-kling}).
\item \eng{dispose} $\to$ \eng{disposal} (perte du \textbf{e} final, \textbf{s} se prononce \textbf{z} : « dis-po-zeul »).
\item \eng{incinerate} $\to$ \eng{incinerator} (pas de \textbf{e} après \textbf{c}, accent sur la 3e syllabe).
\item \eng{manage} $\to$ \eng{management} (ajout de \textbf{ment}, \textbf{g} dur « gué »).
\end{itemize}
\end{attention}

\courssub{Key collocations and idiomatic expressions (Collocations et expressions figées)}

\begin{propriete}[Collocations incontournables du thème]
Ces associations de mots sont figées ; il faut les apprendre par cœur pour parler naturellement.
\begin{center}
\begin{tabularx}{\linewidth}{|X|X|}
\hline
\rowcolor{popblueL} Collocation (Anglais) & Traduction / Contexte \\
\hline
\eng{collect waste / collect rubbish} & ramasser les ordures (verbe + objet) \\
\hline
\eng{dispose of waste} & éliminer / se débarrasser des déchets (\textbf{of} obligatoire) \\
\hline
\eng{sort waste / separate waste} & trier les déchets \\
\hline
\eng{recycle materials} & recycler des matériaux \\
\hline
\eng{empty the bin / empty the dustbin} & vider la poubelle \\
\hline
\eng{load the lorry / load the truck} & charger le camion \\
\hline
\eng{pay collection fees} & payer la redevance d'enlèvement \\
\hline
\eng{illegal dumping / fly-tipping (GB)} & dépôt sauvage d'ordures \\
\hline
\eng{reduce, reuse, recycle} & les 3 R : réduire, réutiliser, recycler \\
\hline
\end{tabularx}
\end{center}
\end{propriete}

\begin{exemplebox}[Collocations en contexte]
\eng{Please \textbf{dispose of} your batteries properly.} = « Veuillez éliminer vos piles correctement. »

\eng{The municipality urges residents to \textbf{sort their waste} before collection.} = « La municipalité exhorte les habitants à trier leurs déchets avant la collecte. »

\eng{\textbf{Illegal dumping} is a major problem in the outskirts of Bamenda.} = « Les dépôts sauvages sont un problème majeur en périphérie de Bamenda. »
\end{exemplebox}

\courssub{Tableau récapitulatif GB / US / français}

\begin{center}
\begin{tabularx}{\linewidth}{|X|X|X|}
\hline
\rowcolor{popblueL} British English & American English & Français \\
\hline
dustman / refuse collector & garbage collector / garbage man & éboueur \\
\hline
dustcart / bin lorry & garbage truck & camion à ordures \\
\hline
dustbin / rubbish bin & garbage can / trash can & poubelle \\
\hline
rubbish dump / tip & dump / landfill & décharge / déchetterie \\
\hline
kerbside collection & curbside collection & collecte en bordure de trottoir \\
\hline
rubbish & garbage / trash & ordures \\
\hline
\end{tabularx}
\end{center}

\begin{methode}[Retenir les couples GB / US]
\begin{enumerate}
\item Apprends chaque mot avec sa variante : \eng{dustman} (GB) — \eng{garbage collector} (US) ; \eng{bin lorry} (GB) — \eng{garbage truck} (US).
\item Choisis une seule variété (britannique ou américaine) dans une production écrite, et reste cohérent du début à la fin.
\item Ne mélange jamais : évite \eng{The dustman drives a garbage truck} dans une même phrase soutenue.
\item Pour le contexte camerounais, retiens les expressions figées : \eng{HYSACAM is a waste management company} ; \eng{households pay refuse collection fees}.
\end{enumerate}
\end{methode}

\begin{savaistu}[Le saviez-vous ? HYSACAM et le secteur informel]
Au Cameroun, \textbf{HYSACAM} (\emph{Hygiène et Salubrité du Cameroun}) est la principale société déléguée pour la collecte des ordures ménagères dans les grandes villes (Yaoundé, Douala, Bafoussam, etc.). Elle gère la collecte porte-à-porte et l'enfouissement en \eng{landfill} (ex. le site de Nkolfoulou près de Yaoundé).
Parallèlement, le \textbf{secteur informel} joue un rôle majeur : les \eng{scavengers} (souvent appelés \eng{waste pickers} dans les rapports internationaux) récupèrent plastiques, métaux, cartons sur les décharges ou dans les rues pour les revendre aux industries de recyclage. Ce travail, bien que précaire, alimente une économie circulaire locale.
\end{savaistu}

\begin{exoresolu}[Compléter avec le vocabulaire de la collecte]
Énoncé — Complete the sentences with: \eng{landfill, scavengers, refuse collection fees, refuse collectors, recycling centre, wheelbarrow}.

1. The \ldots\ empty the dustbins every Tuesday morning.
2. Every family in the neighbourhood pays \ldots\ to the waste management company.
3. \ldots\ search the dump for scrap metal and plastic bottles.
4. The street sweeper pushed his \ldots\ full of rubbish.
5. Old newspapers and bottles are sent to the \ldots\ .
6. Most of the city's waste is buried in a \ldots\ outside Buea.

\tcblower
Solution détaillée :

1. \eng{refuse collectors} — ce sont les éboueurs qui vident les poubelles (terme GB officiel).
2. \eng{refuse collection fees} — la redevance payée par les ménages (terme GB cohérent).
3. \eng{Scavengers} — les récupérateurs informels (majuscule en début de phrase).
4. \eng{wheelbarrow} — la brouette du balayeur (un seul mot, deux \textbf{r}).
5. \eng{recycling centre} — le centre de recyclage pour les matériaux réutilisables (orthographe GB \textbf{centre}).
6. \eng{landfill} — le site d'enfouissement où les déchets sont enterrés.
\end{exoresolu}

\begin{attention}[Pièges fréquents des francophones — Synthèse]
\begin{itemize}
\item \textbf{Faux ami} : Ne dis pas \eng{a discharge} pour « une décharge » : dis \eng{a dump} ou \eng{a landfill}.
\item \textbf{Pluriel irrégulier} : \eng{dustman} $\to$ \eng{dustmen} (comme \eng{man/men}).
\item \textbf{Noms indénombrables} : \eng{garbage}, \eng{rubbish}, \eng{waste}, \eng{litter}, \eng{scrap metal} ne prennent jamais de \textbf{s} au pluriel et ne s'utilisent pas avec \eng{a/an}. Faux : \eng{a rubbish, two garbages}.
\item \textbf{Orthographe} : \eng{wheelbarrow} (un seul mot, deux \textbf{r}) ; \eng{incinerator} (pas de \textbf{e} après \textbf{c}) ; \eng{kerbside} (GB) contre \eng{curbside} (US) ; \eng{recycling centre} (GB \textbf{re}).
\item \textbf{Préposition} : \eng{dispose \textbf{of}} (ne pas oublier \textbf{of}).
\item \textbf{Prononciation} (approx. française) : \eng{collector} « keu-lèk-teur » ; \eng{landfill} « land-file » ; \eng{incinerator} « in-si-neu-réi-teur » ; \eng{disposal} « dis-po-zeul ».
\end{itemize}
\end{attention}

\begin{exercice}[Your turn: describe your neighbourhood]
Write five sentences about waste collection in your town or quarter. Use at least five of these words: \eng{refuse collectors, bin lorry, door-to-door collection, landfill, street sweeper, refuse collection fees, recycling centre}.
\end{exercice}

\coursec{Disposal, Recycling and the 3 Rs}

Dans la section précédente, nous avons suivi les déchets depuis la poubelle domestique jusqu'au camion : \eng{the dustman}, \eng{the refuse truck}, \eng{the landfill}. Mais une fois collectés, que deviennent-ils ? C'est toute la question de la \cle{disposal} (l'élimination) et du \cle{recycling} (le recyclage). C'est aussi le thème favori des sujets d'examen sur l'environnement : la fameuse hiérarchie des \cle{3 Rs} — \eng{Reduce, Reuse, Recycle}.

\courssub{Observing the vocabulary in context}

Lisons d'abord un court texte de sensibilisation, du type de ceux que l'on trouve sur les affiches des mairies anglophones ou dans les brochures des ONG environnementales.

\begin{textebox}[A poster from the Bamenda City Council]
To reduce waste, sort your rubbish: glass, paper and plastic go into separate bins. Organic waste can be composted and used in gardens. Never dump waste in rivers or gutters — it causes floods and diseases.
\end{textebox}

\textbf{Glossaire}

\begin{center}
\begin{tabularx}{\linewidth}{|l|l|X|}
\hline
\rowcolor{popblueL} English & Nature & Français \\
\hline
to reduce waste & verbe & réduire les déchets \\
\hline
to sort (rubbish) & verbe & trier (les ordures) \\
\hline
separate bins & nom & poubelles distinctes / de tri \\
\hline
organic waste & nom & déchets organiques / biodéchets \\
\hline
to compost & verbe & composter (transformer en compost) \\
\hline
to dump & verbe & déverser, jeter (illégalement) \\
\hline
gutter & nom & caniveau, rigole \\
\hline
flood & nom & inondation \\
\hline
\end{tabularx}
\end{center}

\textbf{Questions de compréhension}
\begin{enumerate}
\item Which three materials should go into separate bins?
\item What can we do with organic waste?
\item Why is it dangerous to dump waste in rivers and gutters?
\end{enumerate}

Ce texte de trois phrases contient déjà l'essentiel du lexique de cette section : \eng{reduce}, \eng{sort}, \eng{compost}, \eng{dump}. Observons maintenant chaque notion en détail.

\courssub{The 3 Rs: Reduce, Reuse, Recycle}

\begin{definition}[The waste hierarchy]
La \cle{waste hierarchy} (hiérarchie du traitement des déchets) classe les actions de la plus écologique à la moins écologique : d'abord \cle{reduce} (réduire à la source), ensuite \cle{reuse} (réutiliser), puis \cle{recycle} (recycler), et enfin \cle{recover} (valoriser, par exemple produire de l'énergie en brûlant les déchets). Jeter à la décharge (\eng{landfill}) est la dernière option.
\end{definition}

\begin{itemize}
\item \cle{Reduce} = réduire la quantité de déchets produits. \eng{Reduce your consumption of plastic bottles.} « Réduisez votre consommation de bouteilles en plastique. »
\item \cle{Reuse} = réutiliser un objet au lieu de le jeter. \eng{We should reuse plastic bags instead of throwing them away.} « Nous devrions réutiliser les sacs plastiques au lieu de les jeter. »
\item \cle{Recycle} = transformer un déchet en nouveau produit. \eng{Old newspapers are recycled into cardboard.} « Les vieux journaux sont recyclés en carton. »
\item \cle{Recover} (parfois ajouté, on parle alors des 4 Rs) = valoriser, récupérer de l'énergie ou des matières. \eng{Some plants recover energy by burning waste.} « Certaines usines valorisent les déchets en les brûlant pour produire de l'énergie. »
\end{itemize}

\begin{popfigure}
\begin{tikzpicture}[node distance=2.6cm]
\node[draw=popgreen, fill=popgreenL, circle, minimum size=2.2cm, align=center, line width=1.2pt, font=\bfseries] (reduce) {REDUCE\\ \footnotesize buy less};
\node[draw=poporange, fill=poporangeL, circle, minimum size=2.2cm, align=center, line width=1.2pt, font=\bfseries, right=of reduce] (reuse) {REUSE\\ \footnotesize use again};
\node[draw=popblue, fill=popblueL, circle, minimum size=2.2cm, align=center, line width=1.2pt, font=\bfseries, below=of reuse] (recycle) {RECYCLE\\ \footnotesize make new};
\draw[-{Stealth[length=3mm]}, line width=1.5pt, popdark] (reduce) -- (reuse);
\draw[-{Stealth[length=3mm]}, line width=1.5pt, popdark] (reuse) -- (recycle);
\draw[-{Stealth[length=3mm]}, line width=1.5pt, popdark] (recycle.west) to[bend right=30] (reduce.south);
\end{tikzpicture}
\legende{Le cercle vertueux des 3 Rs : réduire, réutiliser, recycler.}
\end{popfigure}

\begin{attention}[Prononciation des 3 Rs]
\eng{Recycle} se prononce « ri-saï-keul » avec un /s/ (et non « ré-si-cler » à la française). \eng{Reuse} se prononce « ri-iouze » (le préfixe \eng{re-} est long : « rii »). \eng{Reduce} se prononce « ri-diousse ». Attention aussi au nom \eng{recycling} : « ri-saï-kling ».
\end{attention}

\courssub{Key verbs: getting rid of waste}

Voici les verbes essentiels pour parler de l'élimination des déchets, classés de l'action la plus propre à la plus polluante.

\begin{center}
\begin{tabularx}{\linewidth}{|l|X|X|}
\hline
\rowcolor{popblueL} English & Français & Exemple \\
\hline
to sort / separate waste & trier les déchets & Sort your waste into categories. \\
\hline
to recycle & recycler & We recycle glass bottles. \\
\hline
to reuse & réutiliser & She reuses old jars. \\
\hline
to compost & composter & They compost kitchen waste. \\
\hline
to dispose of waste & éliminer les déchets & The factory disposes of waste safely. \\
\hline
to throw away & jeter (à la poubelle) & Don't throw away old clothes. \\
\hline
to get rid of & se débarrasser de & We got rid of the old fridge. \\
\hline
to burn / incinerate & brûler / incinérer & The plant incinerates medical waste. \\
\hline
to bury & enfouir, enterrer & Waste is buried in landfills. \\
\hline
to dump & déverser (illégalement) & People dump rubbish in the river. \\
\hline
\end{tabularx}
\end{center}

\begin{propriete}[to dispose of : la préposition OF est obligatoire]
\eng{To dispose of} est \textbf{toujours} suivi de la préposition \cle{of} : \eng{to dispose of waste safely} = « éliminer les déchets en toute sécurité ». Ne dites jamais \eng{to dispose the waste} (sans \eng{of}). De la même façon, le nom est \cle{waste disposal} : \eng{safe waste disposal} = « l'élimination sûre des déchets ». Même logique pour \eng{to get rid of} : le \eng{of} est obligatoire.
\end{propriete}

\begin{exemplebox}[Exemples traduits]
\eng{We should reuse plastic bags instead of throwing them away.} « Nous devrions réutiliser les sacs plastiques au lieu de les jeter. »

\eng{The city launched a clean-up campaign last month.} « La ville a lancé une campagne de nettoyage le mois dernier. »

\eng{Hospitals must dispose of medical waste very carefully.} « Les hôpitaux doivent éliminer les déchets médicaux avec beaucoup de précautions. »

\eng{Dumping waste in gutters blocks the drains and causes floods.} « Déverser des déchets dans les caniveaux bouche les canalisations et provoque des inondations. »
\end{exemplebox}

\courssub{Sorting waste: the vocabulary of recycling}

Le \cle{sorting} (le tri) consiste à séparer les déchets en catégories : \eng{to sort waste into categories}. On distingue :

\begin{itemize}
\item \cle{recyclable waste} : les déchets recyclables (verre, papier, plastique, métal) ;
\item \cle{non-recyclable waste} : les déchets non recyclables ;
\item \cle{organic waste} ou \cle{organic matter} : les déchets organiques / la matière organique (épluchures, restes de nourriture), qui peut être \eng{composted}.
\end{itemize}

Dans les pays anglophones, on utilise souvent des \cle{colour-coded bins} (poubelles à code couleur) : une \eng{sorting bin} (poubelle de tri) pour chaque matériau — \eng{glass containers}, \eng{paper containers}, \eng{plastic containers}, \eng{metal containers} (conteneurs à verre, à papier, à plastique, à métal).

Les déchets triés partent ensuite dans une \cle{recycling plant} (usine de recyclage), où ils sont \eng{processed} (traités, transformés) : \eng{to process waste}. On obtient des \cle{recycled materials} (matériaux recyclés) : \eng{recycled paper} (papier recyclé), \eng{recycled plastic}, etc.

\begin{definition}[Useful expressions]
\begin{itemize}
\item \cle{waste disposal} : l'élimination des déchets ;
\item \cle{waste management} : la gestion des déchets ;
\item \cle{waste treatment} : le traitement des déchets ;
\item \cle{zero waste} : zéro déchet (mode de vie qui vise à ne produire aucun déchet) ;
\item \cle{a clean-up campaign / operation} : une campagne / opération de nettoyage ;
\item \cle{compost} (nom) : le compost, engrais naturel obtenu par décomposition des déchets organiques.
\end{itemize}
\end{definition}

\begin{attention}[Faux ami : composter]
En français, « composter » signifie souvent « valider un billet » (composter son ticket). En anglais, \eng{to compost} signifie uniquement « transformer en compost » : \eng{We compost fruit peels in the garden.} « Nous compostons les épluchures de fruits au jardin. » Le nom \eng{compost} se prononce « kome-poste ». Autre piège : \eng{disposal} se prononce « dis-peu-zeul » (accent sur « peu »), et non « disposale ».
\end{attention}

\courssub{Exercice résolu et entraînement}

\begin{exoresolu}[Complétez avec le mot qui convient]
Énoncé — Complete the sentences with: \eng{recycle, dump, compost, dispose of, sort, clean-up}.

1. Please \ldots\ your rubbish: paper goes in the blue bin.
2. It is forbidden to \ldots\ waste in the Wouri River.
3. Factories must \ldots\ toxic waste safely.
4. We \ldots\ vegetable peels to make natural fertilizer.
5. Our school organised a \ldots\ campaign in the neighbourhood.
6. Supermarkets should \ldots\ cardboard boxes.

\tcblower
Solution détaillée :
\begin{enumerate}
\item \eng{sort} — la phrase parle de séparer les déchets par poubelle : « trier ».
\item \eng{dump} — jeter illégalement dans une rivière = \eng{to dump}.
\item \eng{dispose of} — attention à la préposition \eng{of} obligatoire : « éliminer ».
\item \eng{compost} — transformer des épluchures en engrais naturel = \eng{to compost}.
\item \eng{clean-up} — \eng{a clean-up campaign} = « une campagne de nettoyage ».
\item \eng{recycle} — transformer le carton en nouveau produit = \eng{to recycle}.
\end{enumerate}
\end{exoresolu}

\begin{exercice}[À votre tour]
1. Traduisez en anglais : « Nous devrions réduire les déchets, réutiliser les bouteilles et recycler le papier. »
2. Complétez : \eng{Glass, paper and plastic are r\ldots\ materials, but plastic bags are often non-r\ldots .}
3. Trouvez l'erreur : \eng{The company must dispose the chemical waste.}
\end{exercice}

\begin{corrige}[Exercice « À votre tour »]
1. \eng{We should reduce waste, reuse bottles and recycle paper.}
2. \eng{recyclable} ; \eng{non-recyclable}.
3. Il manque la préposition : \eng{The company must dispose of the chemical waste.}
\end{corrige}

\begin{savaistu}[E-waste in West Africa]
Le Ghana et le Nigeria accueillent d'énormes décharges de \cle{e-waste} (déchets électroniques : vieux ordinateurs, téléphones, téléviseurs) venus d'Europe et d'Amérique. La décharge d'Agbogbloshie, à Accra, est l'une des plus connues au monde. Ces appareils contiennent des métaux précieux mais aussi des substances toxiques qui polluent le sol et l'eau. C'est un sujet très fréquent dans les épreuves du Baccalauréat sur l'environnement : pensez aux expressions \eng{electronic waste}, \eng{toxic fumes} (fumées toxiques) et \eng{to pollute the soil} (polluer le sol) !
\end{savaistu}

\begin{aretenir}[L'essentiel de la section]
\begin{itemize}
\item La hiérarchie des déchets : \eng{Reduce} (réduire) → \eng{Reuse} (réutiliser) → \eng{Recycle} (recycler) → parfois \eng{Recover} (valoriser).
\item \eng{To dispose of} et \eng{to get rid of} : la préposition \eng{of} est obligatoire.
\item \eng{To dump} = jeter illégalement ; \eng{to sort} = trier ; \eng{to compost} = transformer en compost (faux ami avec « composter » un billet).
\item Prononciations : \eng{recycle} « ri-saï-keul », \eng{reuse} « ri-iouze », \eng{compost} « kome-poste », \eng{disposal} « dis-peu-zeul ».
\item Expressions clés : \eng{waste management}, \eng{waste treatment}, \eng{zero waste}, \eng{a clean-up campaign}.
\end{itemize}
\end{aretenir}

\coursec{Collocations and idiomatic expressions}

Dans la section précédente, nous avons découvert le vocabulaire de l'élimination des déchets et la règle des \eng{3 Rs} (\eng{reduce, reuse, recycle}). Mais connaître des mots isolés ne suffit pas pour parler un anglais naturel : il faut savoir \emph{avec quels mots ils s'associent}. Un francophone qui traduit mot à mot dira peut-être \eng{make the garbage out} au lieu de \eng{take out the rubbish} : la phrase sera comprise, mais elle ne sonnera jamais anglaise. Cette section est donc consacrée aux \cle{collocations} (associations habituelles de mots) et aux \cle{expressions idiomatiques} du domaine des déchets — un point régulièrement évalué à l'épreuve d'anglais du Baccalauréat, dans les questions de grammaire et de vocabulaire comme dans la composition.

\courssub{What is a collocation? Observation first}

\begin{definition}[Collocation]
Une \cle{collocation} est une association de mots qui vont naturellement ensemble dans une langue, par habitude et non par logique. En anglais, on \eng{takes out} the rubbish, on \eng{drops} litter, on \eng{empties} a bin : aucun autre verbe ne convient vraiment, même si un synonyme existe. Les collocations ne se traduisent pas mot à mot : elles s'apprennent comme des blocs.
\end{definition}

Observons d'abord un échange authentique entre deux voisins de Bonabéri, à Douala, la veille du jour de collecte. Lis le dialogue une première fois pour le sens général, puis une seconde fois en soulignant les associations verbe + nom.

\begin{textebox}[At the gate — the night before collection day]
— Could you take out the rubbish, please? The bin is overflowing.

— Sure. By the way, the dustmen come early tomorrow, so I'll put the bin out tonight.

— Good idea. Last week the lorry passed before seven and we missed the collection. Our household waste stayed in front of the house for days.

— Terrible. And have you seen the street corner? People keep dropping litter everywhere. It's a litter-strewn mess.

— I know. The council should introduce a fine for littering and a ban on plastic bags.

— Absolutely. We also need a real solution to our waste problems, not just promises. Anyway, I'll empty the kitchen bin and sort the rubbish — bottles on one side, the rest on the other.

— Thanks, neighbour. One day this city will clean up its act!
\end{textebox}

\textbf{Glossaire}

\begin{center}
\begin{tabularx}{\linewidth}{|l|l|X|}
\hline
\rowcolor{popblueL} English & Nature & Français \\
\hline
to overflow & verbe & déborder \\
\hline
dustmen & nom (plur.) & éboueurs (GB) \\
\hline
household waste & nom & déchets ménagers \\
\hline
to drop litter & verbe + nom & jeter des détritus \\
\hline
a fine & nom & une amende \\
\hline
the council & nom & la municipalité \\
\hline
to sort & verbe & trier \\
\hline
\end{tabularx}
\end{center}

\textbf{Questions de compréhension}
\begin{enumerate}
\item Why will the neighbour put the bin out tonight?
\item What happened last week?
\item What two measures does the first speaker suggest?
\end{enumerate}

\begin{corrige}[Questions de compréhension]
1. \eng{Because the dustmen come early the next morning.} (« Parce que les éboueurs passent tôt le lendemain matin. »)

2. \eng{Last week the lorry passed before seven and they missed the collection, so their household waste stayed in front of the house for days.}

3. \eng{He suggests a fine for littering and a ban on plastic bags.} (« Il propose une amende pour jet de détritus et une interdiction des sacs plastiques. »)
\end{corrige}

\courssub{Verb + noun collocations}

\begin{propriete}[Les collocations verbe + nom du domaine des déchets]
Le verbe et le nom forment une unité figée. Voici les collocations essentielles à connaître en Terminale :
\begin{itemize}
\item \eng{to take out the rubbish} = sortir la poubelle / les ordures
\item \eng{to put the bin out} = mettre la poubelle dehors (pour la collecte)
\item \eng{to empty the bin} = vider la poubelle
\item \eng{to collect waste} = ramasser / collecter les déchets
\item \eng{to drop litter} = jeter des détritus (par terre)
\item \eng{to pick up litter} = ramasser des détritus
\item \eng{to dump waste illegally} = déverser des déchets illégalement
\item \eng{to sort rubbish} = trier les ordures
\item \eng{to burn garbage} = brûler les ordures
\end{itemize}
\end{propriete}

\begin{exemplebox}[Exemples en contexte]
\eng{Could you take out the rubbish before you go to school?} « Peux-tu sortir la poubelle avant d'aller à l'école ? »

\eng{The dustmen collect waste twice a week in our neighbourhood.} « Les éboueurs ramassent les déchets deux fois par semaine dans notre quartier. »

\eng{He was fined for dropping litter.} « Il a été condamné à une amende pour avoir jeté des détritus. »

\eng{Volunteers picked up litter along the Wouri riverbank on Saturday.} « Des volontaires ont ramassé des détritus le long des berges du Wouri samedi. »

\eng{Some factories dump waste illegally at night.} « Certaines usines déversent des déchets illégalement la nuit. »

\eng{We sort our rubbish: plastic here, organic waste there.} « Nous trions nos ordures : le plastique ici, les déchets organiques là. »

\eng{Burning garbage in the open air pollutes the atmosphere.} « Brûler les ordures à l'air libre pollue l'atmosphère. »
\end{exemplebox}

\begin{attention}[Ne traduis pas mot à mot !]
On dit \eng{to drop litter} (jeter des détritus) et non \eng{to throw litter on the floor}, qui est un calque du français « jeter des détritus par terre ». De plus, \eng{to litter} existe aussi comme verbe à part entière : \eng{Don't litter!} = « Ne jetez pas de détritus ! » (c'est la phrase que l'on lit sur les panneaux). Autre piège : on dit \eng{to take out the rubbish} et non \eng{to go out with the rubbish}. Enfin, rappelle-toi que \eng{rubbish}, \eng{litter}, \eng{waste} et \eng{garbage} sont \textbf{indénombrables} : jamais de \eng{a} ni de pluriel — on dira \eng{some rubbish}, \eng{a piece of litter}.
\end{attention}

\courssub{Adjective + noun collocations}

Les adjectifs (ou noms employés comme adjectifs) s'associent eux aussi à des noms précis :

\begin{center}
\begin{tabularx}{\linewidth}{|X|X|}
\hline
\rowcolor{popblueL} Collocation & Traduction française \\
\hline
household waste & déchets ménagers \\
\hline
hazardous waste & déchets dangereux \\
\hline
illegal dumping & dépôt sauvage (de déchets) \\
\hline
overflowing bins & poubelles débordantes \\
\hline
a litter-strewn street & une rue jonchée de détritus \\
\hline
a rubbish-strewn beach & une plage jonchée d'ordures \\
\hline
\end{tabularx}
\end{center}

\begin{exemplebox}[Exemples]
\eng{Hazardous waste from hospitals must be treated separately.} « Les déchets dangereux des hôpitaux doivent être traités séparément. »

\eng{Illegal dumping has become a serious problem on the outskirts of Yaoundé.} « Le dépôt sauvage est devenu un grave problème à la périphérie de Yaoundé. »

\eng{After the festival, the field was a litter-strewn mess.} « Après la fête, le terrain était jonché de détritus. »
\end{exemplebox}

Note la prononciation de \eng{hazardous} : « a-zeur-deuss » (trois syllabes, accent sur la première), et de \eng{strewn} : « stroune » (une seule syllabe ; le \eng{w} final ne se prononce pas).

\courssub{Noun + preposition collocations}

En anglais comme en français, certains noms réclament une préposition précise — mais ce n'est pas toujours la même dans les deux langues :

\begin{center}
\begin{tabularx}{\linewidth}{|X|l|X|}
\hline
\rowcolor{popblueL} English & Préposition & Français \\
\hline
a ban on plastic bags & \eng{on} & une interdiction des sacs plastiques \\
\hline
a fine for littering & \eng{for} & une amende pour jet de détritus \\
\hline
a shortage of landfills & \eng{of} & une pénurie de décharges \\
\hline
a solution to waste problems & \eng{to} & une solution aux problèmes de déchets \\
\hline
\end{tabularx}
\end{center}

\begin{attention}[La préposition change tout]
Attention à \eng{a solution \textbf{to} a problem} : en anglais, la préposition est \eng{to}, alors que le français dit « une solution \emph{au} problème ». Ne dis jamais \eng{a solution of the problem}. De même : \eng{a ban \textbf{on} smoking}, \eng{a fine \textbf{for} speeding} — mémorise toujours le nom avec sa préposition.
\end{attention}

\begin{exemplebox}[Exemples]
\eng{The government has announced a ban on plastic bags.} « Le gouvernement a annoncé une interdiction des sacs plastiques. »

\eng{There is a fine for littering in the city centre.} « Il y a une amende pour jet de détritus au centre-ville. »

\eng{Recycling is part of the solution to our waste problems.} « Le recyclage fait partie de la solution à nos problèmes de déchets. »
\end{exemplebox}

\courssub{Idiomatic expressions}

\begin{definition}[Expression idiomatique]
Une \cle{expression idiomatique} est une phrase figée dont le sens global ne peut pas se deviner à partir des mots pris un à un. Elle se traduit par un équivalent, jamais mot à mot.
\end{definition}

\begin{exemplebox}[Les quatre idiomes à connaître]
\eng{One man's trash is another man's treasure.} « Les déchets des uns font le bonheur des autres. » (ce qui ne vaut rien pour l'un peut être précieux pour l'autre)

\eng{to throw money down the drain} « jeter l'argent par les fenêtres » (gaspiller de l'argent)

\eng{a litter bug} « un pollueur, quelqu'un qui jette des détritus partout » (terme familier et légèrement péjoratif)

\eng{to clean up one's act} « changer de conduite, se reprendre en main » (adopter un meilleur comportement)
\end{exemplebox}

\begin{exemplebox}[En contexte]
\eng{The old bottles I threw away were sold by a recycler — one man's trash is another man's treasure!} « Les vieilles bouteilles que j'ai jetées ont été revendues par un récupérateur — les déchets des uns font le bonheur des autres ! »

\eng{Buying bottled water when tap water is safe is throwing money down the drain.} « Acheter de l'eau en bouteille quand l'eau du robinet est potable, c'est jeter l'argent par les fenêtres. »

\eng{Don't be a litter bug — use the bin!} « Ne sois pas un pollueur — utilise la poubelle ! »

\eng{Our town must clean up its act and stop illegal dumping.} « Notre ville doit changer de conduite et mettre fin au dépôt sauvage. »
\end{exemplebox}

Note la conjugaison du verbe irrégulier employé dans le premier exemple : \eng{throw — threw — thrown} (\eng{I threw away}, \eng{I have thrown away}).

\begin{savaistu}[The “litter bug”]
\eng{A litter bug} est un terme informel et légèrement péjoratif pour désigner quelqu'un qui salit l'espace public. Il est très utilisé dans les campagnes de sensibilisation britanniques : dès les années 1950, des affiches montraient un insecte caricatural (un \eng{bug}) jetant des papiers, avec le slogan \eng{Don't be a litter bug!}. L'expression est ensuite entrée dans la langue courante, et l'on trouve aujourd'hui le même type de campagnes dans de nombreux pays anglophones, de l'Australie au Nigeria.
\end{savaistu}

\courssub{Summary table and mind map}

\begin{popfigure}
\begin{center}
\begin{tikzpicture}[
  col/.style={rectangle, rounded corners=2pt, draw=popline, fill=popcream, align=left, text width=6.4cm, inner sep=4pt, font=\small},
  trad/.style={rectangle, rounded corners=2pt, draw=popline, fill=popblueL, align=left, text width=5.4cm, inner sep=4pt, font=\small},
  head/.style={rectangle, rounded corners=2pt, draw=popblue, fill=popblue, text=white, align=center, font=\bfseries\small, inner sep=4pt}
]
\matrix[column sep=0.4cm, row sep=0.15cm] {
\node[head, text width=6.4cm] {Collocation}; & \node[head, text width=5.4cm] {Traduction française}; \\
\node[col] {to take out the rubbish}; & \node[trad] {sortir la poubelle}; \\
\node[col] {to empty the bin}; & \node[trad] {vider la poubelle}; \\
\node[col] {to drop / pick up litter}; & \node[trad] {jeter / ramasser des détritus}; \\
\node[col] {to sort rubbish}; & \node[trad] {trier les ordures}; \\
\node[col] {a ban on plastic bags}; & \node[trad] {une interdiction des sacs plastiques}; \\
\node[col] {a fine for littering}; & \node[trad] {une amende pour jet de détritus}; \\
\node[col] {a solution to waste problems}; & \node[trad] {une solution aux problèmes de déchets}; \\
\node[col] {one man's trash is another man's treasure}; & \node[trad] {les déchets des uns font le bonheur des autres}; \\
};
\end{tikzpicture}
\end{center}
\legende{Tableau récapitulatif : les collocations et idiomes essentiels de la leçon}
\end{popfigure}

\courssub{How to learn collocations}

\begin{methode}[Mémoriser les collocations efficacement]
\begin{enumerate}
\item N'apprends jamais un mot seul : mémorise-le \textbf{avec son partenaire habituel}. Apprends \eng{take out the rubbish} comme un bloc, pas \eng{rubbish} tout seul.
\item Fais attention aux variétés d'anglais : en Grande-Bretagne on dit \eng{take out the \textbf{rubbish}} et \eng{bin} ; aux États-Unis, \eng{take out the \textbf{garbage / trash}} et \eng{trash can}. Ne mélange pas les deux dans une même rédaction.
\item Note chaque collocation dans une phrase complète qui te parle : \eng{I take out the rubbish every evening in Mvog-Ada.}
\item Réemploie la collocation à l'oral et à l'écrit dans la semaine : une collocation utilisée trois fois est une collocation acquise.
\item Pour les noms, apprends la préposition avec le nom : \eng{a fine \textbf{for}}, \eng{a ban \textbf{on}}, \eng{a solution \textbf{to}}.
\end{enumerate}
\end{methode}

\courssub{Using the collocations in context: mini-dialogues}

\begin{textebox}[Dialogue 1 — Between neighbours on collection day]
— Morning, Mr Ekane! Have you put your bin out? The refuse collectors are already on our street.

— Oh no, I completely forgot! Could you help me carry it? It's full of household waste.

— Of course. You know, you should sort your rubbish — the council now collects plastic separately.

— You're right. I promise I'll clean up my act from next week!
\end{textebox}

\begin{textebox}[Dialogue 2 — A student and a dustman in Buea]
— Excuse me, sir. Why is there so much litter on this street?

— Because people drop litter instead of using the bins, young man. And when the bins are overflowing, the wind spreads everything.

— That's awful. What can students like me do?

— Simple: don't be a litter bug, pick up litter when you see it, and tell your friends. One clean street can change a whole neighbourhood.

— Thank you, sir. We'll organise a clean-up day at our school.
\end{textebox}

\courssub{Solved exercise}

\begin{exoresolu}[Complete with the correct collocation]
Complète chaque phrase avec la collocation qui convient : \eng{take out the rubbish}, \eng{a fine for littering}, \eng{drop litter}, \eng{overflowing}, \eng{a solution to}, \eng{one man's trash is another man's treasure}.

1. Please \dots\ before the dustmen arrive.
2. The bin in the school yard is \dots ; we must empty it.
3. In Singapore, you pay \dots\ if you throw paper on the ground.
4. Recycling is \dots\ the waste problems of big cities.
5. Don't \dots ! Use the bin over there.
6. The scrap dealer bought my old radio — well, \dots !
\tcblower
\textbf{Solution détaillée}

1. \eng{Please \textbf{take out the rubbish} before the dustmen arrive.} — collocation verbe + nom : « sortir la poubelle ».

2. \eng{The bin in the school yard is \textbf{overflowing}; we must empty it.} — collocation adjectif + nom : \eng{an overflowing bin}, « une poubelle débordante ».

3. \eng{In Singapore, you pay \textbf{a fine for littering} if you throw paper on the ground.} — collocation nom + préposition : \eng{a fine \textbf{for}} + nom en \eng{-ing}.

4. \eng{Recycling is \textbf{a solution to} the waste problems of big cities.} — attention à la préposition \eng{to} (et non \eng{of}).

5. \eng{Don't \textbf{drop litter}! Use the bin over there.} — on dit \eng{drop litter}, jamais \eng{throw litter on the floor}.

6. \eng{The scrap dealer bought my old radio — well, \textbf{one man's trash is another man's treasure}!} — expression idiomatique : « les déchets des uns font le bonheur des autres ».
\end{exoresolu}

\begin{exercice}[Your turn — Translate into English]
Traduis en anglais en utilisant les collocations de la leçon :
\begin{enumerate}
\item Les éboueurs ramassent les déchets très tôt le matin.
\item Il a payé une amende pour avoir jeté des détritus.
\item Brûler les ordures est interdit en ville.
\item Notre quartier doit changer de conduite.
\end{enumerate}
\end{exercice}

\begin{corrige}[Exercice « Your turn »]
1. \eng{The dustmen collect waste very early in the morning.}

2. \eng{He paid a fine for dropping litter.} (ou \eng{He was fined for littering.})

3. \eng{Burning garbage is forbidden in town.} — le verbe \eng{burn} admet deux formes au prétérit et au participe passé : \eng{burn — burnt — burnt} (surtout GB) ou \eng{burn — burned — burned} (surtout US).

4. \eng{Our neighbourhood must clean up its act.}
\end{corrige}

\begin{aretenir}[L'essentiel de la section]
\begin{itemize}
\item Une collocation s'apprend en bloc : \eng{take out the rubbish}, \eng{empty the bin}, \eng{drop litter}, \eng{pick up litter}, \eng{sort rubbish}, \eng{dump waste illegally}.
\item Noms + prépositions figées : \eng{a ban \textbf{on}}, \eng{a fine \textbf{for}}, \eng{a shortage \textbf{of}}, \eng{a solution \textbf{to}}.
\item Idiomes : \eng{one man's trash is another man's treasure}, \eng{to throw money down the drain}, \eng{a litter bug}, \eng{to clean up one's act}.
\item Pièges de francophone : jamais \eng{throw litter on the floor}, jamais \eng{a solution of the problem} ; GB \eng{rubbish/bin}, US \eng{garbage/trash can} ; \eng{rubbish}, \eng{litter}, \eng{waste} sont indénombrables.
\end{itemize}
\end{aretenir}

\coursec{Word families, spelling and false friends}

Après avoir maîtrisé les collocations et expressions idiomatiques du secteur de la gestion des déchets (\emph{Collocations and idiomatic expressions}), nous devons maintenant consolider les bases morphologiques et orthographiques. Connaître les familles de mots permet de multiplier son vocabulaire à partir d'une seule racine ; éviter les faux amis et les pièges d'orthographe garantit une communication précise, à l'écrit comme à l'oral.

\courssub{Word families and derivation : building your lexicon}

\begin{textebox}[Observation : from verb to noun and adjective]
Read the following sentences and notice how the same root changes its form to play different grammatical roles.

\eng{The municipality \textbf{collects} household waste every Tuesday.} (verb)
\eng{The \textbf{collection} of recyclable materials has improved.} (noun)
\eng{A \textbf{collector} came early this morning.} (noun, person)
\eng{It is a \textbf{collective} responsibility.} (adjective)

\eng{We must \textbf{recycle} more plastic.} (verb)
\eng{\textbf{Recycling} creates jobs.} (noun)
\eng{This bottle is \textbf{recyclable}.} (adjective)

\eng{They \textbf{dispose} of toxic waste carefully.} (verb + preposition)
\eng{The \textbf{disposal} of nuclear waste is a major issue.} (noun)
\eng{These cups are \textbf{disposable}.} (adjective)

\eng{Factories \textbf{pollute} the river.} (verb)
\eng{\textbf{Pollution} kills fish.} (noun)
\eng{Carbon monoxide is a dangerous \textbf{pollutant}.} (noun)
\eng{The \textbf{polluted} water is unfit for drinking.} (adjective)
\end{textebox}

\begin{propriete}[Les principaux suffixes de dérivation]
En anglais, on forme les noms et les adjectifs à partir des verbes grâce à des suffixes réguliers. Les connaître permet de deviner le sens de mots inconnus.

\begin{center}
\begin{tabularx}{\linewidth}{|X|X|X|X|}
\hline
\rowcolor{popblueL} \textbf{Verbe (racine)} & \textbf{Nom d'action / état (-tion, -sion, -al, -ing)} & \textbf{Nom d'agent / instrument (-er, -or)} & \textbf{Adjectif (-able, -ive, -al, -ed, -ing)} \\
\hline
to \cle{collect} & \cle{collection} & \cle{collector} & \cle{collective}, \cle{collectible} \\
\hline
to \cle{recycle} & \cle{recycling} & \cle{recycler} (person/machine) & \cle{recyclable} \\
\hline
to \cle{dispose} (of) & \cle{disposal} (\emph{attention : -al}) & \cle{disposer} (rare) & \cle{disposable} \\
\hline
to \cle{pollute} & \cle{pollution} & \cle{polluter} & \cle{polluted}, \cle{polluting} \\
\hline
to \cle{reduce} & \cle{reduction} & \cle{reducer} (technique) & \cle{reducible} \\
\hline
to \cle{reuse} & \cle{reuse} (identique) & \cle{reuser} (rare) & \cle{reusable} \\
\hline
to \cle{incinerate} & \cle{incineration} & \cle{incinerator} & \cle{incinerable} \\
\hline
to \cle{compost} & \cle{composting} & \cle{composter} (bin) & \cle{compostable} \\
\hline
to \cle{separate} & \cle{separation} & \cle{separator} (machine) & \cle{separate}, \cle{separated} \\
\hline
to \cle{treat} & \cle{treatment} & \cle{treatment plant} (installation) & \cle{treatable}, \cle{treated} \\
\hline
\end{tabularx}
\end{center}

\textbf{Règles clés :}
\begin{itemize}
    \item \textbf{-tion / -sion} : action ou résultat (\eng{collection, pollution, incineration, separation}). Notez l'orthographe de \eng{disposal} (terminaison \emph{-al} et non \emph{-tion}) et de \eng{treatment} (\emph{-ment}).
    \item \textbf{-er / -or} : personne ou machine qui fait l'action (\eng{collector, incinerator, recycler, separator}). Souvent \emph{-or} pour les mots d'origine latine technique (\eng{incinerator, compactor, separator}).
    \item \textbf{-able} : « qui peut être… » (\eng{recyclable, reusable, biodegradable, compostable, disposable, treatable}). C'est un suffixe productif : on peut l'ajouter à presque tout verbe transitif.
    \item \textbf{Participles comme adjectifs} : \eng{polluted} (subi), \eng{polluting} (actif), \eng{treated}, \eng{separated}.
\end{itemize}
\end{propriete}

\begin{popfigure}
\begin{tikzpicture}[
    mindmap,
    grow cyclic,
    every node/.style={concept, execute at begin node=\hskip0pt},
    root concept/.append style={concept color=popblue, minimum size=2.5cm, text=white, font=\large\bfseries},
    level 1 concept/.append style={concept color=popblue!70, sibling angle=90, minimum size=1.8cm, text=white, font=\bfseries},
    level 2 concept/.append style={concept color=poporange!80, sibling angle=45, minimum size=1.5cm, text=white, font=\small},
    level 3/.style={level distance=3.5cm},
    concept/.append style={text width=2.2cm, align=center}
]
\node {collect} [clockwise from=30]
    child { node {collection} }
    child { node {collector} }
    child { node {collective} }
    child { node {collectible} };
\end{tikzpicture}
\legende{Arbre de dérivation de la racine \eng{collect}.}
\end{popfigure}

\begin{methode}[Construire et mémoriser un tableau de famille de mots]
Pour chaque nouveau verbe clé du thème, suivez ces 4 étapes :
\begin{enumerate}
    \item \textbf{Identifiez la racine} (ex. \eng{recycl-}).
    \item \textbf{Remplissez un tableau 4 colonnes} : Verbe / Nom d'action / Nom d'agent / Adjectif.
    \item \textbf{Écrivez une phrase exemple par case} (contexte personnel ou camerounais).
    \item \textbf{Révisez par colonnes} : dites tous les noms d'action, puis tous les adjectifs, etc.
\end{enumerate}
\emph{Exemple personnel :} \eng{My father works as a \textbf{collector} for HYSACAM. The \textbf{collection} starts at 5 AM. It is a \textbf{collective} effort to keep Yaoundé clean.}
\end{methode}

\begin{exoresolu}[Compléter le tableau de dérivation]
Complétez le tableau ci-dessous pour les verbes \eng{to separate}, \eng{to treat}, \eng{to dump}. Indiquez la forme nom (action), nom (agent/instrument) et adjectif si elles existent.

\begin{center}
\begin{tabular}{|l|l|l|l|}
\hline
\rowcolor{popblueL} \textbf{Verb} & \textbf{Noun (action/state)} & \textbf{Noun (agent/tool)} & \textbf{Adjective} \\
\hline
to separate & \dots & \dots & \dots \\
\hline
to treat & \dots & \dots & \dots \\
\hline
to dump & \dots & \dots & \dots \\
\hline
\end{tabular}
\end{center}

\tcblower
\textbf{Solution détaillée :}

\begin{center}
\begin{tabular}{|l|l|l|l|}
\hline
\rowcolor{popgreenL} \textbf{Verb} & \textbf{Noun (action/state)} & \textbf{Noun (agent/tool)} & \textbf{Adjective} \\
\hline
to separate & \textbf{separation} & \textbf{separator} (machine) & \textbf{separate} / \textbf{separated} \\
\hline
to treat & \textbf{treatment} & \textbf{treatment plant} (pas d'agent simple \textit{treater}) & \textbf{treatable} / \textbf{treated} \\
\hline
to dump & \textbf{dumping} / \textbf{dump} (lieu) & \textbf{dumper} (camion) / \textbf{dump truck} & \textbf{dumped} \\
\hline
\end{tabular}
\end{center}

\emph{Notes :} Pour \eng{treat}, l'agent est souvent une installation (\eng{waste treatment plant}). Pour \eng{dump}, le nom d'action est \eng{dumping} (acte souvent illégal) et le lieu est \eng{a dump} ou \eng{a landfill}. \eng{Dumper} désigne le camion (dump truck).
\end{exoresolu}

\courssub{Key adjectives and spelling traps}

\begin{propriete}[Adjectifs essentiels du thème « déchets »]
Ces adjectifs qualifient la nature des déchets ou leur devenir. Maîtrisez leur orthographe et leur traduction.

\begin{center}
\begin{tabularx}{\linewidth}{|X|X|X|}
\hline
\rowcolor{popblueL} \textbf{Adjectif} & \textbf{Traduction} & \textbf{Exemple contextualisé} \\
\hline
\cle{biodegradable} & biodégradable & \eng{Food waste is \textbf{biodegradable}.} « Les déchets alimentaires sont biodégradables. » \\
\hline
\cle{non-biodegradable} & non biodégradable & \eng{Plastic bags are \textbf{non-biodegradable}.} \\
\hline
\cle{recyclable} & recyclable & \eng{Is this cardboard \textbf{recyclable}?} \\
\hline
\cle{reusable} & réutilisable & \eng{Bring a \textbf{reusable} bag to the market.} \\
\hline
\cle{disposable} & jetable & \eng{\textbf{Disposable} cups create tons of waste.} \\
\hline
\cle{toxic} / \cle{hazardous} & toxique / dangereux & \eng{\textbf{Hazardous} waste needs special \textbf{disposal}.} \\
\hline
\cle{organic} & organique (d'origine vivante) & \eng{\textbf{Organic} waste can be composted.} \\
\hline
\cle{single-use} & à usage unique (jetable) & \eng{The ban on \textbf{single-use} plastics starts Monday.} \\
\hline
\end{tabularx}
\end{center}
\end{propriete}

\begin{attention}[Pièges d'orthographe fréquents (Francophones)]
\begin{itemize}
    \item \textbf{environment} : un \emph{n} avant le \emph{m} (pas *enviroment). Pensez à « \emph{iron} » dans le milieu.
    \item \textbf{garbage} : un seul \emph{r} (pas *garrbage).
    \item \textbf{rubbish} : deux \emph{b} (pas *rubish).
    \item \textbf{recycling} : un seul \emph{c} après \emph{re-} (pas *recycling avec double c). Règle : préfixe \emph{re-} + base \emph{cycle}.
    \item \textbf{separate} : \emph{a} puis \emph{e} (adj/verbe). Mnémotechnique : « Sep \textbf{A} \textbf{R}at \textbf{E} » (séparer un rat). L'adjectif se prononce \eng{/sep(ə)rət/}, le verbe \eng{/sepəreɪt/}.
    \item \textbf{disposal} : terminaison \emph{-al} (pas *disposition qui signifie « arrangement » ou « tempérament »).
    \item \textbf{litter} (déchets jetés au sol) : deux \emph{t}.
    \item \textbf{waste} (déchet) vs \textbf{waist} (taille) : homophones parfaits \eng{/weɪst/}.
\end{itemize}
\end{attention}

\begin{exemplebox}[Spelling practice]
\eng{The \textbf{environment} ministry launched a campaign against \textbf{litter}.} \\
\eng{Separate your waste: put \textbf{recyclable} items in the blue \textbf{bin}.} \\
\eng{\textbf{Hazardous} \textbf{disposal} procedures are strict.} \\
\eng{Do not \textbf{bury} \textbf{non-biodegradable} plastic.}
\end{exemplebox}

\courssub{Pronunciation traps: homophones and tricky sounds}

\begin{savaistu}[Waste vs Waist : homophones parfaits]
\eng{Waste} (déchet, gaspillage) et \eng{waist} (taille, ceinture) se prononcent \textbf{exactement de la même façon} : \eng{/weɪst/} (« ouèiste »). Seul le contexte (ou l'orthographe) permet de les distinguer.
\begin{itemize}
    \item \eng{Industrial \textbf{waste} pollutes rivers.}
    \item \eng{He wore a belt around his \textbf{waist}.}
\end{itemize}
Autres homophones utiles : \eng{bin} (poubelle) / \eng{been} (participe de be) ; \eng{sort} (trier) / \eng{sought} (cherché, passé de seek).
\end{savaistu}

\begin{attention}[Prononciation piège : \eng{bury}]
Le verbe \cle{bury} (enterrer) ne se prononce \textbf{PAS} « bioury » (comme \eng{busy} ou \eng{fury}).
\begin{itemize}
    \item \textbf{Prononciation correcte :} \eng{/ˈberi/} → se note en français « \textbf{bè-ri} » (comme \eng{berry} : baie, ou \eng{very}).
    \item \textbf{Erreur fréquente :} lire \eng{u} comme \eng{/ʌ/} ou \eng{/juː/}.
\end{itemize}
\emph{Exemple :} \eng{They \textbf{bury} the waste in a \textbf{landfill}.} « Ils enterrent les déchets dans une décharge contrôlée. »
\end{attention}

\courssub{False friends and calques from French}

C'est le point le plus critique pour les francophones. Le transfert direct du français vers l'anglais produit des erreurs systématiques. Le tableau ci-dessous oppose l'erreur calquée (à ne jamais écrire) à la forme anglaise correcte.

\begin{propriete}[Tableau des faux amis et calques — À apprendre par cœur]
\begin{center}
\begin{tabularx}{\linewidth}{|X|X|X|X|}
\hline
\rowcolor{popblueL} \textbf{Français (source de l'erreur)} & \textbf{Calque / Faux ami (INCORRECT)} & \textbf{Anglais correct} & \textbf{Nature / Note} \\
\hline
une décharge (contrôlée) & \eng{a discharge} & \eng{a landfill} / \eng{a dump} & \eng{Discharge} = rejet (liquide/gaz), sortie d'hôpital, décharge électrique. \\
\hline
les ordures / déchets & \eng{orders} & \eng{rubbish} (UK) / \eng{garbage} (US) / \eng{waste} (technique) & \eng{Orders} = commandes (achat) ou ordres (militaire). \\
\hline
une poubelle & \eng{a poubelle} & \eng{a bin} (UK) / \eng{a trash can} / \eng{a garbage can} (US) & \eng{Poubelle} n'existe pas en anglais (nom propre Poubelle = inventeur). \\
\hline
ramasser les ordures & \eng{to ramass} / \eng{to ramasse} & \eng{to collect waste} / \eng{to pick up trash} & \eng{Ramass} n'existe pas. \eng{Collect} est le terme technique. \\
\hline
jeter (qch) & \eng{to jeter} & \eng{to throw away} / \eng{to dump} (sauvage) / \eng{to discard} & \eng{Jeter} n'existe pas. \eng{Throw away} = mettre à la poubelle. \\
\hline
un sac poubelle & \eng{a poubelle bag} & \eng{a bin bag} / \eng{a trash bag} / \eng{a garbage bag} & \\
\hline
tri sélectif & \eng{selective sorting} & \eng{waste separation} / \eng{selective collection} / \eng{sorting} & \eng{Sorting} seul suffit souvent. \\
\hline
biodégradable & \eng{biodegradable} & \eng{biodegradable} & \textbf{Vrai ami !} Orthographe identique. \\
\hline
incinérateur & \eng{an incinerator} & \eng{an incinerator} & \textbf{Vrai ami !} Attention à la prononciation \eng{/ɪnˈsɪnəreɪtə(r)/}. \\
\hline
\end{tabularx}
\end{center}
\end{propriete}

\begin{attention}[Focus : \eng{Disposal} vs \eng{Disposable}]
Ne confondez pas le nom et l'adjectif :
\begin{itemize}
    \item \eng{\textbf{Waste disposal}} = l'élimination des déchets (nom, processus).
    \item \eng{\textbf{Disposable} gloves / cups / income} = gants \textbf{jetables}, gobelets \textbf{jetables}, revenu \textbf{disponible} (adjectif).
\end{itemize}
\eng{The \textbf{disposal} of \textbf{disposable} plastics is a global challenge.} « L'élimination des plastiques jetables est un défi mondial. »
\end{attention}

\begin{exemplebox}[Correction d'erreurs typiques d'élèves]
\begin{tabularx}{\linewidth}{|X|X|}
\hline
\rowcolor{poporangeL} \textbf{Phrase incorrecte (calque français)} & \textbf{Phrase correcte} \\
\hline
\eng{The \textbf{discharge} is full.} & \eng{The \textbf{landfill} is full.} / \eng{The \textbf{dump} is full.} \\
\hline
\eng{Put the \textbf{orders} in the \textbf{poubelle}.} & \eng{Put the \textbf{rubbish} in the \textbf{bin}.} \\
\hline
\eng{The truck comes to \textbf{ramass} the trash.} & \eng{The truck comes to \textbf{collect} the trash.} \\
\hline
\eng{Don't \textbf{jeter} that, it's \textbf{recyclable}.} & \eng{Don't \textbf{throw that away}, it's \textbf{recyclable}.} \\
\hline
\eng{We need more \textbf{dispositions} for waste.} & \eng{We need better \textbf{disposal} methods.} \\
\hline
\end{tabularx}
\end{exemplebox}

\courssub{Consolidation: using the vocabulary in context}

\begin{textebox}[Reading: Waste Management Challenges in Buea]
Buea, the capital of the Southwest Region, faces growing waste management challenges. With a rising population and expanding markets like Molyko and Mile 16, the volume of \textbf{household waste} has doubled in ten years. The municipal \textbf{collection} service, often \textbf{run by HYSACAM} using compactors, struggles to cover all neighbourhoods, especially the hilly areas where \textbf{collection trucks} cannot easily access.

Many residents resort to \textbf{illegal dumping} in streams or unused land, creating \textbf{eyesores} and health hazards. \textbf{Plastic waste}, mostly \textbf{non-biodegradable} bags and bottles, \textbf{clogs} \textbf{drainage systems}, causing floods during the rainy season. Burning waste in the open air releases \textbf{toxic fumes}, worsening air quality.

However, local initiatives are emerging. The ``Clean Buea'' association organises monthly \textbf{clean-up campaigns}. They promote \textbf{waste separation} at source: \textbf{organic waste} for \textbf{composting}, \textbf{recyclables} (metal, glass, PET) for the \textbf{recycling} chain. Schools now teach pupils the \textbf{3 Rs}: \textbf{Reduce}, \textbf{Reuse}, \textbf{Recycle}. A new \textbf{waste treatment plant} is planned to include an \textbf{incinerator} for \textbf{hazardous} medical waste and a \textbf{landfill} with a \textbf{leachate} treatment system.

Citizens must understand that \textbf{proper disposal} is not only the council's duty. Using \textbf{reusable} bags, avoiding \textbf{single-use} plastics, and reporting \textbf{fly-tipping} are simple acts that protect the \textbf{environment}.

\medskip
\textbf{Glossaire :}
\begin{itemize}
    \item \eng{eyesore} (n.) : chose disgracieuse, tare visuelle.
    \item \eng{clog} (v.) : boucher, obstruer.
    \item \eng{drainage system} (n.) : système d'évacuation (caniveaux).
    \item \eng{fly-tipping} (n., UK) : dépôt sauvage d'ordures (équivalent \eng{illegal dumping}).
    \item \eng{leachate} (n.) : lixiviat (liquide pollué s'écoulant d'une décharge).
    \item \eng{single-use} (adj.) : à usage unique (jetable).
    \item \eng{run by} (phr. v.) : géré par, exploité par.
\end{itemize}
\end{textebox}

\begin{exercice}[Exploitation du texte et révision de la section]
\begin{enumerate}
    \item \textbf{Word families :} Trouvez dans le texte les formes dérivées (noms, adjectifs) des verbes suivants : \eng{collect, recycle, dispose, pollute, treat, separate, incinerate}.
    \item \textbf{False friends :} Traduisez en anglais correct (évitez les calques) :
    \begin{enumerate}
        \item La décharge municipale est pleine.
        \item J'ai mis les ordures dans la poubelle.
        \item Les éboueurs ramassent les déchets mardi.
        \item Ne jette pas cette bouteille, elle est recyclable.
        \item Les gobelets jetables sont interdits au marché.
    \end{enumerate}
    \item \textbf{Spelling :} Corrigez l'orthographe de ces mots : \emph{enviroment, garbadge, rubish, recyclling, seperate, disposel, bury (prononciation ?)}.
    \item \textbf{Production écrite :} En 50 mots, écrivez un message WhatsApp pour sensibiliser vos camarades de classe au tri des déchets au lycée, en utilisant au moins 5 mots de la famille \eng{collect/recycle/dispose} et l'adjectif \eng{reusable}.
\end{enumerate}
\end{exercice}

\begin{corrige}[Exercice n°1 -- Exploitation et révision]
\begin{enumerate}
    \item \textbf{Familles de mots (réponses du texte) :}
    \begin{itemize}
        \item \eng{collect} $\rightarrow$ \eng{collection} (n), \eng{collection service} (nom composé).
        \item \eng{recycle} $\rightarrow$ \eng{recycling} (n), \eng{recyclables} (adj/pluriel nom), \eng{recycling chain} (nom composé).
        \item \eng{dispose} $\rightarrow$ \eng{disposal} (n), \eng{proper disposal}.
        \item \eng{pollute} $\rightarrow$ \eng{pollution} (absent mais implicite), \eng{toxic fumes} (polluants), \eng{air quality}.
        \item \eng{treat} $\rightarrow$ \eng{treatment plant}, \eng{leachate treatment}.
        \item \eng{separate} $\rightarrow$ \eng{separation} (n), \eng{waste separation}.
        \item \eng{incinerate} $\rightarrow$ \eng{incinerator} (n).
    \end{itemize}
    \item \textbf{Traductions correctes (pas de calques) :}
    \begin{enumerate}
        \item \eng{The municipal \textbf{landfill} is full.} (ou \eng{dump} si non contrôlée).
        \item \eng{I put the \textbf{rubbish/garbage} in the \textbf{bin}.}
        \item \eng{The \textbf{collectors} \textbf{collect} waste on Tuesdays.} (ou \eng{The garbage truck picks up the trash...})
        \item \eng{Don't \textbf{throw away} that bottle, it's \textbf{recyclable}.}
        \item \eng{\textbf{Disposable} cups are banned at the market.}
    \end{enumerate}
    \item \textbf{Orthographe corrigée :}
    \begin{itemize}
        \item \emph{enviroment} $\rightarrow$ \textbf{environment}
        \item \emph{garbadge} $\rightarrow$ \textbf{garbage}
        \item \emph{rubish} $\rightarrow$ \textbf{rubbish}
        \item \emph{recyclling} $\rightarrow$ \textbf{recycling}
        \item \emph{seperate} $\rightarrow$ \textbf{separate}
        \item \emph{disposel} $\rightarrow$ \textbf{disposal}
        \item \emph{bury} $\rightarrow$ orthographe correcte, mais prononciation \eng{/beri/} (« bè-ri »), pas « bioury ».
    \end{itemize}
    \item \textbf{Production écrite (modèle) :}
    \eng{``Hey friends! Let's keep our school clean. \textbf{Reduce} waste: bring a \textbf{reusable} water bottle. \textbf{Separate} your trash: plastic in the \textbf{recycling} \textbf{bin}, food in the compost. The \textbf{collector} comes Friday. Don't \textbf{throw away} plastic on the ground! \textbf{Proper disposal} protects our \textbf{environment}. \#CleanCampus''}
\end{enumerate}
\end{corrige}

\begin{aretenir}[Synthèse de la Section 5]
\begin{itemize}
    \item \textbf{Familles de mots :} Apprenez les triplets \eng{Verb / Noun / Adjective} (ex. \eng{recycle – recycling – recyclable}). Suffixes clés : \eng{-tion/-sion}, \eng{-al} (\eng{disposal}), \eng{-ment} (\eng{treatment}), \eng{-er/-or}, \eng{-able}.
    \item \textbf{Orthographe :} \eng{environment, garbage, rubbish, recycling, separate, disposal}. Pas de double consonne là où le français en met une (sauf \eng{rubbish, litter}).
    \item \textbf{Prononciation :} \eng{waste = waist} (/weɪst/) ; \eng{bury} = « bè-ri » (/beri/).
    \item \textbf{Faux amis :} \eng{discharge $\neq$ décharge} (→ \eng{landfill/dump}) ; \eng{orders $\neq$ ordures} (→ \eng{rubbish/waste}) ; \eng{poubelle $\neq$ bin} ; \eng{ramasser $\neq$ collect} ; \eng{jeter $\neq$ throw away}.
    \item \textbf{Confusion nom/adjectif :} \eng{disposal} (nom, élimination) vs \eng{disposable} (adj, jetable).
    \item \textbf{Variantes UK/US :} \eng{rubbish/bin} (UK) vs \eng{garbage/trash/garbage can} (US). Le standard camerounais est l'anglais britannique.
\end{itemize}
\end{aretenir}

\coursec{Using the vocabulary in context}

Dans les sections précédentes, nous avons étudié les familles de mots (\eng{waste, wasteful, to waste}), l'orthographe et les faux amis (\eng{discharge} ne signifie pas « décharge » !). Il est maintenant temps de \cle{réemployer tout ce lexique dans des situations réelles} : lire un article de presse, compléter des phrases, corriger des calques du français, décrire une image et rédiger un paragraphe argumentatif. C'est exactement ce que l'on vous demandera au BAC A, aussi bien en \emph{reading comprehension} qu'en \emph{essay writing}.

\courssub{Guided reading: “Keeping Our City Clean”}

Lisons d'abord un article original sur la gestion des déchets à Bamenda. Pendant la lecture, \cle{soulignez mentalement tous les mots du lexique} étudiés dans cette leçon : vous en retrouverez plus de vingt.

\begin{textebox}[Keeping Our City Clean — a newspaper article]
In Bamenda, the city council has introduced colour-coded bins: green for organic waste, blue for recyclables and red for hazardous waste. Refuse collectors empty the bins three times a week, and a new recycling centre has opened near the market. However, some residents still dump rubbish in gutters, which blocks drains and causes floods during the rainy season.

The mayor explained the new policy at a press conference last week. “For too long, our city has suffered from illegal dumping and overflowing dustbins,” she said. “Uncollected waste attracts flies and rats, and it pollutes our streams. A clean city is a healthy city.” She added that households which fail to sort their waste will pay a small fine, while those who recycle regularly will receive a discount on their sanitation tax.

The recycling centre already employs thirty young people. They sort plastic bottles, aluminium cans and cardboard, which are then sold to factories in Douala. “Before, I had no job,” says Claudine, a sorter at the centre. “Now I earn a living by turning rubbish into raw materials.”

Not everyone is satisfied, however. Some traders at the food market complain that the collection truck sometimes arrives late, leaving piles of rotting vegetables in the sun. “The service is better than before,” admits one shopkeeper, “but during the rainy season, we need the bins emptied every day.”

The city council promises to buy two more garbage trucks next year and to launch an awareness campaign in schools. “Children must learn the three Rs — reduce, reuse, recycle — from an early age,” the mayor concluded. “They are the ones who will keep Bamenda clean tomorrow.”
\end{textebox}

\textbf{Glossaire (mots difficiles)}\par
\begin{itemize}
\item \eng{colour-coded} (adj.) : à code couleur ; \eng{bin} (n.) : poubelle, bac
\item \eng{recyclables} (n. pluriel) : matériaux recyclables ; \eng{hazardous} (adj.) : dangereux
\item \eng{refuse collector} (n.) : éboueur ; \eng{to dump} (v.) : déverser, jeter sauvagement
\item \eng{gutter} (n.) : caniveau ; \eng{drain} (n.) : égout, canal d'évacuation
\item \eng{illegal dumping} (n.) : dépôt sauvage de déchets ; \eng{overflowing} (adj.) : débordant
\item \eng{fine} (n.) : amende ; \eng{sanitation tax} (n.) : taxe d'assainissement
\item \eng{sorter} (n.) : trieur/trieuse ; \eng{raw materials} (n.) : matières premières
\item \eng{awareness campaign} (n.) : campagne de sensibilisation
\end{itemize}

\textbf{Comprehension questions}\par
\begin{enumerate}
\item What do the three colours of the bins mean?
\item How often do refuse collectors empty the bins?
\item What problem does dumping rubbish in gutters cause?
\item What reward do households that recycle receive?
\item What does Claudine do at the recycling centre?
\item What complaint do the market traders have?
\end{enumerate}

\begin{corrige}[Questions de compréhension]
\begin{enumerate}
\item Green is for organic waste, blue is for recyclables and red is for hazardous waste.
\item They empty the bins three times a week.
\item It blocks drains and causes floods during the rainy season.
\item They receive a discount on their sanitation tax.
\item She sorts plastic bottles, aluminium cans and cardboard, which are then sold to factories in Douala.
\item They complain that the collection truck sometimes arrives late, leaving piles of rotting vegetables in the sun.
\end{enumerate}
\end{corrige}

\courssub{Gap-filling: activating the target vocabulary}

\begin{exoresolu}[Complete the sentences]
Complétez avec : \eng{bin, landfill, sort, dump, recycle, litter, collection}.
\tcblower
\begin{enumerate}
\item Please throw your empty bottle in the \textbf{bin}, not on the ground. (\emph{Jetez votre bouteille vide à la poubelle.})
\item Old fridges and batteries are taken to a special site, not to the ordinary \textbf{landfill}. (\emph{décharge contrôlée})
\item In Germany, families must \textbf{sort} their waste into different containers. (\emph{trier})
\item It is forbidden to \textbf{dump} chemical waste into rivers. (\emph{déverser})
\item We \textbf{recycle} newspapers and cardboard at school. (\emph{recycler})
\item Do not drop \textbf{litter} in the schoolyard! (\emph{détritus})
\item In our quarter, garbage \textbf{collection} takes place on Mondays and Thursdays. (\emph{ramassage})
\end{enumerate}
\end{exoresolu}

\courssub{From French calques to correct English}

Les francophones traduisent souvent mot à mot. Observons et corrigeons.

\begin{attention}[Les calques à détruire]
\begin{itemize}
\item \eng{The orders are collected by the discharge men.} FAUX (calque de « les ordures sont ramassées par les éboueurs »). Correct : \eng{The rubbish is collected by the refuse collectors.} — \eng{discharge} = décharge électrique ou renvoi d'un employé ; \eng{orders} = commandes !
\item \eng{Throw it in the dust.} FAUX (calque de « jette-le à la poussière »). Correct : \eng{Throw it in the dustbin.} — \eng{dust} = poussière ; la poubelle = \eng{dustbin} (GB) ou \eng{trash can} (US).
\item \eng{We must sensitize the population.} FAUX (calque de « sensibiliser »). Correct : \eng{We must raise awareness among the population.} — \eng{to sensitize} existe mais signifie « rendre sensible » (médecine, photographie).
\end{itemize}
\end{attention}

\begin{exercice}[Correct the calques]
Corrigez : 1. \eng{The garbage men take the ordures every morning.} 2. \eng{Put the papers in the basket of waste.} 3. \eng{They throw their detritus everywhere.}
\end{exercice}

\begin{corrige}[Correct the calques]
1. \eng{The refuse collectors collect the rubbish every morning.} 2. \eng{Put the papers in the wastepaper basket.} 3. \eng{They drop their litter everywhere.}
\end{corrige}

\courssub{Describing a picture: illegal dump vs sorting centre}

\begin{experience}[Guided oral comparison]
Décrivez et comparez deux images (une décharge sauvage, un centre de tri moderne) en suivant la trame :
\begin{enumerate}
\item \textbf{In the first picture}, I can see an illegal dump: piles of rubbish, plastic bags, flies...
\item \textbf{In the second picture}, there is a modern sorting centre: workers sort recyclables on a conveyor belt...
\item \textbf{The main difference is that...} the dump pollutes the soil and water, whereas the recycling centre turns waste into raw materials.
\item \textbf{In my opinion...} our town needs more bins and a proper collection service.
\end{enumerate}
\end{experience}

\courssub{Writing an argumentative paragraph}

\begin{methode}[Rédiger un paragraphe argumentatif (5–8 lignes)]
\begin{enumerate}
\item \textbf{Phrase d'introduction} : posez le problème. \eng{Our school produces too much waste.}
\item \textbf{2–3 idées} avec le lexique du thème : \eng{First, the school should provide colour-coded bins. Second, students must learn to sort their waste. Finally, we could recycle paper and plastic bottles.}
\item \textbf{Phrase de conclusion} : \eng{If we all reduce, reuse and recycle, our school will be cleaner and healthier.}
\item \textbf{Soulignez} les mots de la leçon utilisés (minimum 6).
\end{enumerate}
\end{methode}

\begin{exemplebox}[Modèle de paragraphe]
\eng{Our school should provide more bins and teach students to sort waste.} « Notre école devrait fournir plus de poubelles et apprendre aux élèves à trier les déchets. » \eng{First, colour-coded bins should be placed in every classroom and in the playground. Second, the canteen could stop using plastic bags and reuse metal plates instead. Finally, a small recycling club could collect paper and bottles every Friday. If we reduce, reuse and recycle every day, our school will become a clean and healthy place for everyone.}
\end{exemplebox}

\begin{attention}[Variez le vocabulaire dans un essay]
Évitez de répéter \eng{waste} dix fois. Variez selon le sens exact : \eng{rubbish} (ordures ménagères, GB), \eng{refuse} (terme officiel, soutenu), \eng{litter} (détritus jetés dans l'espace public), \eng{garbage/trash} (US). Chaque mot a son emploi : on dit \eng{to drop litter} mais \eng{household refuse}.
\end{attention}

\begin{center}
\begin{tabularx}{\linewidth}{|X|X|X|}\hline
\rowcolor{popblueL} Français & English & Emploi exact \\ \hline
ordures ménagères & rubbish / garbage & quotidien (GB / US) \\ \hline
déchets (terme général, officiel) & waste / refuse & essays, documents officiels \\ \hline
détritus dans la rue & litter & espace public \\ \hline
décharge sauvage & illegal dump / fly-tipping site & hors de tout contrôle \\ \hline
centre de tri & sorting centre & recyclage organisé \\ \hline
ramassage des ordures & garbage collection & service municipal \\ \hline
\end{tabularx}
\end{center}

\begin{popfigure}
\begin{tikzpicture}[node distance=0.9cm and 1.4cm,
  box/.style={rectangle, rounded corners, draw=popblue, fill=popblueL, align=center, minimum width=3.4cm, minimum height=1cm, font=\small},
  arr/.style={-{Stealth[length=3mm]}, thick, poporange}]
\node[box, align=center] (observe){1. Observe\\ the text or picture};
\node[box, right=of observe, fill=popgreenL, draw=popgreen, align=center] (identify){2. Identify\\ the vocabulary};
\node[box, right=of identify, fill=popgoldL, draw=popgold, align=center] (build){3. Build\\ sentences};
\node[box, right=of build, fill=poppinkL, draw=poppink, align=center] (write){4. Write\\ a paragraph};
\draw[arr] (observe) -- (identify);
\draw[arr] (identify) -- (build);
\draw[arr] (build) -- (write);
\end{tikzpicture}
\legende{Organigramme de production guidée : de l'observation à la rédaction.}
\end{popfigure}

\textbf{Complément Bac.} Au BAC A, les sujets d'essay portent très souvent sur la santé publique et l'environnement : \eng{Discuss the causes and effects of poor waste management in our towns}, \eng{How can young people protect the environment?} Maîtriser ce lexique (collocations, synonymes, verbes précis) améliore directement votre note de \emph{vocabulary range}.

\begin{savaistu}[Vocabulary range au BAC]
Les correcteurs du BAC A attribuent des points spécifiques à la richesse du vocabulaire. Un candidat qui écrit \eng{the authorities should improve garbage collection and build a sanitary landfill instead of tolerating illegal dumping} marque beaucoup plus de points qu'un candidat qui répète \eng{dirt} et \eng{throw}. Entraînez-vous à caser au moins six mots de ce chapitre dans chaque essay sur l'environnement !
\end{savaistu}

\newpage

\coursec{Activité d'intégration}

\courssub{Objectifs de l'activité}

À l'issue de cette activité d'intégration, l'élève sera capable de :
\begin{itemize}
\item réemployer à l'écrit et à l'oral le vocabulaire de la \cle{garbage collection}, du \cle{disposal} et du \cle{recycling} étudié dans la leçon ;
\item produire une affiche de campagne (\emph{poster}) avec un slogan et des règles formulées à l'impératif ;
\item utiliser correctement au moins huit \cle{collocations} du thème : \eng{sort your waste}, \eng{drop litter}, \eng{empty the bins}, \eng{recycle plastic bottles}, \eng{dispose of rubbish}, etc. ;
\item présenter un \emph{pitch} oral de deux minutes avec une prononciation soignée (\eng{garbage} « gar-bidge », \eng{waste} « ouéiste », \eng{bin} « binne ») ;
\item éviter les \cle{faux amis} et les calques du français (\eng{dustbin} et non « poubelle » traduite mot à mot, \eng{rubbish} et non « rubish »).
\end{itemize}

\courssub{Situation}

\begin{textebox}[A letter from the Principal — Government Bilingual High School, Yaoundé]
Dear students,

As you have all noticed, our schoolyard is often full of plastic bags, empty bottles and paper after the break. The dustbins are not always used, and the garbage truck only comes once a week. Last month, the hygiene inspector from the Ministry of Public Health visited our school and was not satisfied.

The Parents' Association has decided to organise a \textbf{“Clean School Week”} from 12 to 16 May. Each class of Terminale must design a campaign poster in English with the title \textbf{“Keep Our School Clean!”}. The poster must include a slogan, five clear rules for students, and a simple diagram of the 3 Rs (Reduce, Reuse, Recycle).

On Friday afternoon, one group from each class will present its poster to the school assembly in a two-minute pitch. The best campaign will win a set of new recycling bins for the winning class, offered by the Yaoundé City Council.

We count on your creativity and on your English!

The Principal
\end{textebox}

\textbf{Glossaire :} \emph{schoolyard} (n., cour de récréation) ; \emph{dustbin} (n., poubelle) ; \emph{garbage truck} (n., camion à ordures) ; \emph{hygiene inspector} (n., inspecteur d'hygiène) ; \emph{assembly} (n., assemblée, rassemblement) ; \emph{bin} (n., bac, poubelle).

Votre groupe de quatre élèves représente la classe de Terminale A. Vous devez concevoir l'affiche et préparer le pitch oral. Le jury est composé des autres groupes de la classe et du professeur d'anglais.

\courssub{Consignes}

\textbf{Tâche 1 — Compréhension et repérage (10 min).} Lisez la lettre du Principal. Soulignez tous les mots et expressions liés au thème de la leçon (\eng{dustbins}, \eng{garbage truck}, \eng{recycling bins}...). Classez-les en trois colonnes : \emph{collection} / \emph{disposal} / \emph{recycling}.

\textbf{Tâche 2 — Brainstorming lexical (10 min).} En groupe, listez au moins douze mots ou collocations de la leçon que vous pourriez utiliser : noms (\eng{waste}, \eng{litter}, \eng{landfill}, \eng{dump}, \eng{bin}, \eng{garbage collector}, \eng{compost}), verbes (\eng{sort}, \eng{throw away}, \eng{dispose of}, \eng{recycle}, \eng{reuse}, \eng{reduce}, \eng{empty}, \eng{collect}) et adjectifs (\eng{recyclable}, \eng{reusable}, \eng{hazardous}, \eng{organic}).

\textbf{Tâche 3 — Rédaction de l'affiche (25 min).} Concevez l'affiche sur une feuille A3 :
\begin{enumerate}
\item un \textbf{titre} : \eng{Keep Our School Clean!} ;
\item un \textbf{slogan} original en anglais (court, rythmé, mémorable) ;
\item \textbf{cinq règles} à l'impératif, chacune contenant une collocation de la leçon ;
\item un \textbf{schéma simple des 3 Rs} (trois flèches en cercle : \eng{Reduce}, \eng{Reuse}, \eng{Recycle}) avec une phrase d'explication pour chaque R.
\end{enumerate}

\textbf{Tâche 4 — Préparation du pitch (15 min).} Rédigez puis répétez un pitch oral de deux minutes : présentation du groupe, lecture du slogan, explication des cinq règles et du schéma des 3 Rs, appel à l'action final. Chaque membre du groupe parle au moins trente secondes. Vérifiez la prononciation des mots difficiles : \eng{garbage} « gar-bidge », \eng{sewage} « sou-idje », \eng{hazardous} « ra-zeur-deusse », \eng{recycling} « ri-saï-kling ».

\textbf{Tâche 5 — Présentation et vote (20 min).} Chaque groupe présente son affiche. Les autres élèves remplissent une grille d'évaluation (lexique : 4 pts ; grammaire et impératifs : 4 pts ; prononciation : 4 pts ; créativité : 4 pts ; respect du temps : 4 pts). La classe vote pour la meilleure campagne.

\textbf{Tâche 6 — Correction collective (10 min).} Le professeur note au tableau les erreurs entendues (faux amis, calques, prononciation) et la classe les corrige ensemble.

\courssub{Ressources à mobiliser}

\begin{aretenir}[Lexique et structures utiles]
\begin{itemize}
\item \textbf{Noms :} \eng{waste / rubbish / garbage / trash / litter} (déchets, ordures), \eng{dustbin / bin / trash can} (poubelle), \eng{garbage collector / dustman} (éboueur), \eng{garbage truck} (camion à ordures), \eng{landfill / dump} (décharge), \eng{recycling centre} (centre de recyclage), \eng{compost} (compost).
\item \textbf{Verbes et collocations :} \eng{to sort (out) waste} (trier les déchets), \eng{to drop litter} (jeter des ordures par terre — au sens négatif), \eng{to throw away} (jeter), \eng{to dispose of} (se débarrasser de), \eng{to empty the bins} (vider les poubelles), \eng{to collect garbage} (ramasser les ordures), \eng{to recycle / reuse / reduce}.
\item \textbf{Adjectifs :} \eng{recyclable} (recyclable), \eng{reusable} (réutilisable), \eng{organic waste} (déchets organiques), \eng{hazardous waste} (déchets dangereux), \eng{household waste} (ordures ménagères).
\item \textbf{Grammaire :} l'\cle{impératif} pour les règles (\eng{Sort your waste! Don't drop litter!}) ; \eng{must} et \eng{should} pour l'obligation et le conseil ; \eng{let's} pour l'appel à l'action collectif.
\item \textbf{Attention :} \eng{waste} est ici indénombrable ; on dit \eng{a waste of time} (une perte de temps) mais \eng{household waste} sans article. Ne confondez pas \eng{to litter} (jeter des détritus) et « litière ».
\end{itemize}
\end{aretenir}

\begin{attention}[Pièges à éviter dans votre production]
\begin{itemize}
\item Ne dites pas « \eng{throw the garbage on the floor is forbidden} » : préférez \eng{Don't drop litter!} ou \eng{Littering is forbidden.}
\item \eng{Garbage} et \eng{rubbish} sont indénombrables : jamais « garbages » ni « a rubbish ».
\item Prononcez \eng{bin} « binne » (et non « bine ») et \eng{waste} « ouéiste ».
\item Faux ami : \eng{a dump} est une décharge sauvage, pas un « dûmp » ; \eng{disposal} signifie « élimination », pas « disposition ».
\end{itemize}
\end{attention}

\courssub{Proposition de production et corrigé commenté}

\begin{exoresolu}[Modèle d'affiche et de pitch — groupe « Green Eagles »]
\textbf{Énoncé.} Produisez l'affiche (slogan + 5 règles + schéma des 3 Rs) et le script du pitch de deux minutes.
\tcblower
\textbf{L'affiche (texte) :}

\emph{Titre :} \eng{KEEP OUR SCHOOL CLEAN!}

\emph{Slogan :} \eng{“Sort it, bin it, recycle it — a clean school is a healthy school!”}

\emph{Les cinq règles :}
\begin{enumerate}
\item \eng{Sort your waste: plastic, paper and organic waste go in different bins.}
\item \eng{Don't drop litter in the schoolyard — use the dustbins!}
\item \eng{Reuse your plastic bottles instead of throwing them away.}
\item \eng{Help the garbage collectors: put full bags next to the gate on collection day.}
\item \eng{Recycle paper and cans, and compost organic waste for the school garden.}
\end{enumerate}

\emph{Schéma des 3 Rs :} trois flèches en cercle avec les légendes \eng{Reduce — buy less, waste less} ; \eng{Reuse — give things a second life} ; \eng{Recycle — turn old waste into new products}.

\textbf{Le pitch (script oral) :}

\eng{Good afternoon, everyone. We are the Green Eagles, and this is our campaign: “Keep Our School Clean!”}

\eng{Every day, students drop litter in the schoolyard: plastic bags, bottles and paper. But the garbage truck only comes once a week, and the dustbins are often full. Our solution is simple: sort it, bin it, recycle it!}

\eng{First, sort your waste: plastic, paper and organic waste must go in different bins. Second, never drop litter — always use the dustbins. Third, reuse your bottles instead of throwing them away. Fourth, help the garbage collectors by putting the bags out on collection day. And fifth, recycle paper and cans, and compost organic waste for our school garden.}

\eng{Remember the 3 Rs: Reduce, Reuse, Recycle. Let's reduce the waste we produce, let's reuse what we can, and let's recycle the rest. A clean school is a healthy school. Let's keep our school clean — together! Thank you.}

\textbf{Commentaire des choix de langue :}
\begin{itemize}
\item \textbf{Lexique de la leçon (11 items)} : \eng{drop litter}, \eng{schoolyard}, \eng{garbage truck}, \eng{dustbins}, \eng{sort your waste}, \eng{organic waste}, \eng{throw away}, \eng{garbage collectors}, \eng{collection day}, \eng{recycle}, \eng{compost}. L'objectif de huit mots est dépassé.
\item \textbf{Collocations exactes} : on dit bien \eng{drop litter}, \eng{sort waste}, \eng{throw away} (phrasal verb séparable), \eng{dispose of} (avec la préposition \eng{of}).
\item \textbf{Impératifs} correctement formés pour les règles (\eng{Sort}, \eng{Don't drop}, \eng{Reuse}, \eng{Help}, \eng{Recycle}) et \eng{let's} pour l'appel à l'action final.
\item \textbf{Indénombrables respectés} : \eng{waste}, \eng{litter}, \eng{garbage} sans article ni pluriel.
\item \textbf{Slogan} court, rythmé et parallèle (\eng{sort it, bin it, recycle it}), facile à mémoriser — \eng{bin} utilisé comme verbe, usage courant en anglais britannique.
\end{itemize}
\end{exoresolu}

\courssub{Bilan de l'activité}

Cette activité a permis de mobiliser l'ensemble des savoir-faire de la leçon dans une tâche finale authentique et motivante :
\begin{itemize}
\item \textbf{Compréhension de l'écrit} : repérage du lexique thématique dans un document administratif réaliste (lettre du Principal).
\item \textbf{Vocabulaire} : réactivation et réemploi des champs lexicaux de la collecte, de l'élimination et du recyclage, avec un minimum de huit collocations par production.
\item \textbf{Grammaire} : pratique de l'impératif, de \eng{must/should} et de \eng{let's} dans des règles et un appel à l'action.
\item \textbf{Expression orale} : pitch chronométré, répartition de la parole dans le groupe, travail de la prononciation des mots pièges.
\item \textbf{Expression écrite} : conception d'un document authentique (affiche) respectant les contraintes du genre (titre, slogan, règles, schéma).
\item \textbf{Compétence citoyenne} : sensibilisation à l'hygiène et à la protection de l'environnement dans le contexte réel du lycée camerounais.
\end{itemize}

Le vote collectif et la correction au tableau renforcent l'apprentissage par les pairs : les erreurs de faux amis (\eng{disposal} $\neq$ « disposition »), d'indénombrables (\eng{rubbish} sans pluriel) et de prononciation (\eng{bin}, \eng{waste}, \eng{hazardous}) sont corrigées une dernière fois, de manière mémorable, avant l'évaluation. L'affiche gagnante peut être réellement affichée dans la cour du lycée, donnant à la production des élèves une visibilité et une utilité concrètes.

\newpage

\coursec{Méthodes et exercices résolus}

\coursec{Lesson 35 — Vocabulary: Words and expressions related to garbage collection, disposal, and recycling services}

\courssub{Waste and its types: the core vocabulary}

\begin{textebox}[Authentic context: A municipal notice in Buea]
The Buea Council wishes to inform residents that \textbf{household waste} collection will take place every Monday and Thursday. Please place your \textbf{wheelie bins} at the roadside by 6:00 a.m. \textbf{Littering} in public spaces is an offence punishable by a fine. \textbf{Hazardous waste} (batteries, paint, chemicals) must not be mixed with \textbf{general waste}; it should be taken to the \textbf{recycling centre} at Muea. Let us keep our township clean and green.
\end{textebox}

\begin{definition}[Key vocabulary: Types of waste]
\begin{itemize}
\item \textbf{Waste} (uncountable) — general term for unwanted materials (\emph{déchets, ordures}). \eng{Household waste} (ordures ménagères), \eng{industrial waste} (déchets industriels), \eng{hazardous waste} (déchets dangereux), \eng{organic waste} (déchets organiques), \eng{e-waste} (déchets électroniques).
\item \textbf{Rubbish} (BrE, uncountable) — synonym of \eng{waste} (\emph{ordures, détritus}). \eng{Garbage} (AmE) and \eng{trash} (AmE) are common in Cameroon but \eng{rubbish} is the standard British form used in exams.
\item \textbf{Litter} (uncountable) — waste thrown on the ground in public places (\emph{détritus, déchets abandonnés}). Verb: \eng{to litter} (jeter des détritus). \eng{Littering} (noun/action) is an offence.
\item \textbf{Refuse} (formal, uncountable) — household or industrial waste (\emph{rebuts, ordures}). Often used in \eng{refuse collection}, \eng{refuse collector}.
\item \textbf{Sewage} (uncountable) — waste water and excrement carried in sewers (\emph{eaux usées, égouts}). Not to be confused with \eng{sewerage} (system of sewers).
\end{itemize}
\end{definition}

\begin{propriete}[Countable vs Uncountable — Critical for the Bac]
\begin{center}
\begin{tabularx}{\linewidth}{|X|X|X|}
\hline
\rowcolor{popblueL} \textbf{Word} & \textbf{Grammar} & \textbf{Correct usage} \\
\hline
Waste / Rubbish / Garbage / Trash / Refuse / Litter / Sewage & Uncountable & \eng{There is a lot of waste.} (singular verb) \\
\hline
A piece of waste / An item of rubbish / A bit of litter & Partitive (countable) & \eng{He picked up a piece of litter.} \\
\hline
Bin / Dustbin / Wheelie bin / Container / Bag / Sack / Truck / Lorry & Countable & \eng{Two bins}, \eng{three lorries}. \\
\hline
\end{tabularx}
\end{center}
\end{propriete}

\begin{attention}[Piège Bac] N'écris jamais \eng{wastes}, \eng{rubbishes}, \eng{litters}, \eng{sewages} au sens « ordures ». Le pluriel \eng{wastes} n'existe que dans des contextes très techniques (ex. \eng{nuclear wastes} — types de déchets nucléaires), hors programme.
\end{attention}

\courssub{Collection services: people, vehicles and places}

\begin{definition}[People, vehicles and places]
\begin{center}
\begin{tabularx}{\linewidth}{|X|X|X|}
\hline
\rowcolor{popblueL} \textbf{English (BrE standard)} & \textbf{French} & \textbf{Notes / AmE variants} \\
\hline
Refuse collector / Dustman / Binman & Éboueur & \eng{Garbage collector} (AmE), \eng{sanitation worker} (formel) \\
\hline
Bin lorry / Dustcart / Refuse truck & Camion-poubelle & \eng{Garbage truck} (AmE) \\
\hline
Wheelie bin / Dustbin / Bin & Poubelle (sur roues / classique) & \eng{Trash can} (AmE), \eng{recycling bin} (poubelle de tri) \\
\hline
Landfill (site) / Landfill site / Dump & Décharge (contrôlée / sauvage) & \eng{Dump} souvent péjoratif (décharge sauvage) ; \eng{landfill} = site d'enfouissement technique \\
\hline
Recycling centre / Recycling plant / Bottle bank & Centre de tri / usine de recyclage / conteneur à verre & \eng{Recycling centre} = \eng{civic amenity site} (UK) \\
\hline
Incinerator / Incineration plant & Incinérateur / usine d'incinération & Verbe : \eng{to incinerate} \\
\hline
Compost heap / Compost bin & Tas de compost / composteur & Verbe : \eng{to compost} \\
\hline
\end{tabularx}
\end{center}
\end{definition}

\begin{exemplebox}[Collocations with collection verbs]
\begin{itemize}
\item \eng{To collect waste / rubbish / bins} (collecter) — \eng{The council collects the bins every Monday.}
\item \eng{To empty the bin / the dustbin} (vider la poubelle) — \eng{Please empty the bin before it overflows.}
\item \eng{To take out the rubbish / the bins} (sortir les poubelles) — \eng{It's your turn to take out the rubbish.}
\item \eng{To pick up litter} (ramasser les détritus) — \eng{Volunteers picked up litter along the beach.}
\item \eng{To dump waste (illegally)} (déverser/décharger sauvagement) — \eng{Companies must not dump toxic waste in rivers.}
\end{itemize}
\end{exemplebox}

\courssub{Disposal, recycling and the 3 Rs}

\begin{definition}[Disposal, treatment and the 3 Rs]
\begin{itemize}
\item \textbf{Disposal} (noun) — \emph{élimination, traitement final}. Collocation: \eng{waste disposal}, \eng{disposal of hazardous waste}, \eng{final disposal site}. Verb: \eng{to dispose of} (préposition \textbf{of} obligatoire). \eng{We must dispose of batteries safely.}
\item \textbf{Recycling} (noun) — \emph{recyclage}. Verb: \eng{to recycle}. Adj: \eng{recyclable} (recyclable), \eng{recycled} (recyclé). \eng{Recyclable materials} (matériaux recyclables).
\item \textbf{The 3 Rs} (Reduce, Reuse, Recycle) — Les 3 R.
\begin{itemize}
\item \eng{Reduce} (réduire) : \eng{Reduce waste at source}, \eng{reduce packaging}.
\item \eng{Reuse} (réutiliser) : \eng{Reuse plastic bags}, \eng{reusable bottle} (bouteille réutilisable).
\item \eng{Recycle} (recycler) : \eng{Recycle paper, glass, metal}.
\end{itemize}
\item \textbf{Composting} (noun) — \emph{compostage}. Verb: \eng{to compost}. \eng{Organic waste can be composted.}
\item \textbf{Incineration} (noun) — \emph{incinération}. Verb: \eng{to incinerate}. \eng{Incineration reduces volume but produces emissions.}
\item \textbf{Separation / Sorting} — \eng{To sort waste} (trier les déchets), \eng{source separation} (tri à la source), \eng{sorting centre} (centre de tri).
\item \textbf{Pollution} — \eng{Air pollution}, \eng{water pollution}, \eng{soil pollution}. Verb: \eng{to pollute}. Adj: \eng{polluted}.
\end{itemize}
\end{definition}

\begin{savaistu}[Le Cameroun et la gestion des déchets]
Au Cameroun, HYSACAM (Hygiène et Salubrité du Cameroun) gère la collecte dans les grandes villes (Yaoundé, Douala, Bafoussam, Bamenda, etc.). Le tri à la source reste peu répandu ; la plupart des déchets finissent dans des \eng{landfills} (ex. Nkolfoulou à Yaoundé, Kpambara à Douala). Des initiatives locales (ONG, start-ups comme \eng{Green Vision} ou \eng{Plastic Man}) développent le recyclage du plastique et la valorisation des déchets organiques en compost ou biogaz. Le vocabulaire de cette leçon te permet de comprendre les rapports d'activité de HYSACAM, les articles de presse (Cameroon Tribune, The Guardian Post) et de t'exprimer sur l'assainissement de ton quartier.
\end{savaistu}

\courssub{Collocations and idiomatic expressions}

\begin{propriete}[Collocations incontournables (à apprendre par cœur)]
\begin{center}
\begin{tabularx}{\linewidth}{|X|X|X|}
\hline
\rowcolor{popblueL} \textbf{Collocation (English)} & \textbf{Traduction} & \textbf{Exemple} \\
\hline
\eng{Waste management} & Gestion des déchets & \eng{Effective waste management is a priority.} \\
\hline
\eng{To dispose of waste} & Éliminer / traiter des déchets & \eng{Hospitals must dispose of medical waste carefully.} \\
\hline
\eng{To throw away / To throw out} & Jeter & \eng{Don't throw away that box; reuse it.} \\
\hline
\eng{To sort (waste) into categories} & Trier (les déchets) en catégories & \eng{Residents are asked to sort waste into three bins.} \\
\hline
\eng{To raise awareness (about/among)} & Sensibiliser (à / auprès de) & \eng{NGOs raise awareness about plastic pollution.} \\
\hline
\eng{To protect the environment from} & Protéger l'environnement contre & \eng{We must protect the environment from degradation.} \\
\hline
\eng{To be littered with} & Être jonché de / encombré de & \eng{The streets were littered with plastic bags.} \\
\hline
\eng{Overflowing bins} & Poubelles qui débordent & \eng{Overflowing bins attract rats and flies.} \\
\hline
\eng{Illegal dumping / Fly-tipping (UK)} & Dépôt sauvage d'ordures & \eng{The council fights illegal dumping with cameras.} \\
\hline
\eng{Zero waste} & Zéro déchet & \eng{A zero-waste lifestyle is difficult but possible.} \\
\hline
\end{tabularx}
\end{center}
\end{propriete}

\begin{exemplebox}[Idiomatic expressions related to waste and cleaning]
\begin{itemize}
\item \eng{To talk rubbish / to talk garbage} — Dire des bêtises, parler pour ne rien dire. \eng{Stop talking rubbish and listen!}
\item \eng{To clean up one's act} — Se corriger, améliorer son comportement. \eng{The factory was forced to clean up its act after the scandal.}
\item \eng{To sweep something under the carpet} — Cacher un problème au lieu de le résoudre. \eng{The council cannot sweep the waste crisis under the carpet.}
\item \eng{One man's trash is another man's treasure} — Ce qui est déchet pour l'un est trésor pour l'autre (proverbe).
\item \eng{To go to waste} — Être gâché. \eng{All that food went to waste.}
\end{itemize}
\end{exemplebox}

\courssub{Word families, spelling and false friends}

\begin{propriete}[Word families — Derivation table]
\begin{center}
\begin{tabular}{|l|l|l|l|}
\hline
\rowcolor{popblueL} \textbf{Verb} & \textbf{Noun (action/state)} & \textbf{Noun (person/agent)} & \textbf{Adjective} \\
\hline
Pollute & Pollution & Polluter & Polluted / Polluting \\
\hline
Recycle & Recycling & Recycler & Recyclable / Recycled \\
\hline
Dispose & Disposal & — & Disposable (jetable) \\
\hline
Collect & Collection & Collector & Collective \\
\hline
Incinerate & Incineration & Incinerator & — \\
\hline
Compost & Composting & — & Compostable \\
\hline
Contaminate & Contamination & Contaminant & Contaminated \\
\hline
Degrade & Degradation & — & Degradable / Biodegradable \\
\hline
\end{tabular}
\end{center}
\end{propriete}

\begin{attention}[Orthographe : les 4 pièges du Bac]
\begin{enumerate}
\item \textbf{Environment} — un \emph{n} avant \emph{-ment} (pas \eng{enviroment}).
\item \textbf{Disposal} — un seul \emph{s} après \emph{di-} (pas \eng{dispOsal} avec deux \emph{s} ; le verbe est \eng{dispose} avec \emph{s}, le nom \emph{disposal} avec \emph{s} simple).
\item \textbf{Beginning} — double \emph{n} (règle du \emph{n} final accentué : \eng{begin} $\rightarrow$ \eng{beginning}).
\item \textbf{Recycling} — un seul \emph{c} après \emph{re-} (pas \eng{recycling} avec \emph{cc}).
\end{enumerate}
\end{attention}

\begin{definition}[False friends — Les 6 paires à ne jamais confondre]
\begin{center}
\begin{tabularx}{\linewidth}{|X|X|X|}
\hline
\rowcolor{popblueL} \textbf{English word} & \textbf{Wrong French guess} & \textbf{Correct meaning} \\
\hline
Garbage & Garage & Ordures (AmE) \\
\hline
Dump & Dune & Décharge (souvent sauvage) ; verbe : jeter/déverser \\
\hline
Sensitize & Sensibiliser & Rendre sensible (médical/technique) — \textbf{Utiliser : \eng{raise awareness}} \\
\hline
Quarter & Quartier & Quart (1/4), quartier militaire (\eng{quarters}) — \textbf{Utiliser : \eng{neighbourhood, district, area}} \\
\hline
Launch & Lancer (jeter) & Lancer (produit, fusée, bateau) — \textbf{Utiliser : \eng{throw away, dump}} \\
\hline
Dirt & Déchet & Saleté, terre, boue (pas synonyme de \eng{waste} en gestion urbaine) \\
\hline
\end{tabularx}
\end{center}
\end{definition}

\begin{exemplebox}[Prononciation repère (lettres françaises)]
\begin{itemize}
\item \eng{Waste} — « ouéiste » (comme \eng{waist} « taille », homophones).
\item \eng{Garbage} — « gâr-bidj » (\emph{g} dur, \emph{ge} final = \emph{dj}).
\item \eng{Sewage} — « sou-idj » (\emph{ew} = \emph{ou}, \emph{age} = \emph{idj}).
\item \eng{Litter} — « li-teur » (\emph{tt} = \emph{t} simple, \emph{er} = \emph{eur}).
\item \eng{Recycle} — « ri-saï-keul » (\emph{cy} = \emph{saï}).
\item \eng{Disposal} — « dis-pô-zal » (\emph{po} = \emph{pô}, \emph{s} = \emph{z}).
\item \eng{Incinerator} — « in-si-né-rei-teur » (accent sur \emph{né}).
\end{itemize}
\end{exemplebox}

\courssub{Using the vocabulary in context — Savoir-faire Bac}

\begin{methode}[Identifier le bon mot du champ lexical des déchets]
Pour choisir le mot anglais exact dans un texte ou un exercice de vocabulaire, procède par étapes :
\begin{enumerate}
\item \textbf{Repère la catégorie} : s'agit-il du déchet lui-même (\eng{waste, rubbish, litter, garbage}), de la personne (\eng{refuse collector, dustman, scavenger}), du véhicule (\eng{bin lorry, garbage truck}), du lieu (\eng{landfill, dumpsite, recycling centre}) ou de l'action (\eng{to collect, to dump, to sort, to recycle}) ?
\item \textbf{Distingue les registres} : \eng{rubbish} et \eng{dustbin} sont britanniques ; \eng{garbage} et \eng{trash can} sont américains. Au Cameroun, on rencontre les deux normes, mais reste cohérent dans une même copie (privilégie le BrE).
\item \textbf{Vérifie la nature grammaticale} : \eng{waste} est généralement indénombrable (on dit \eng{a lot of waste}, jamais \eng{wastes} au pluriel pour des ordures) ; \eng{litter} désigne les déchets jetés dans la rue, pas la litière des animaux dans ce contexte.
\item \textbf{Contrôle la préposition associée} : \eng{to dispose \textbf{of} waste}, \eng{to throw something \textbf{away}}, \eng{to put rubbish \textbf{in} the bin}, \eng{to sort waste \textbf{into} categories}.
\item \textbf{Reformule pour vérifier} : remplace le mot par sa définition simple ; si la phrase reste correcte, le choix est bon.
\end{enumerate}
\end{methode}

\begin{exoresolu}[Compléter un texte à trous — type Bac]
Énoncé. Complete the following passage with the correct words from the box: \emph{landfill — litter — refuse collectors — recycle — bins — waste}.

« In Yaoundé, household \ldots\ldots\ldots\ is collected twice a week by \ldots\ldots\ldots\ who empty the \ldots\ldots\ldots\ into large trucks. Unfortunately, many people still drop \ldots\ldots\ldots\ in the streets instead of using the containers provided. Most of the rubbish ends up in a \ldots\ldots\ldots\ outside the city, although some materials could be sent to factories that \ldots\ldots\ldots\ plastic and metal. »
\tcblower
Solution détaillée.
\begin{itemize}
\item Trou 1 : « household \ldots\ is collected » — il faut un nom indénombrable désignant les ordures ménagères : \textbf{waste}. \eng{Litter} est impossible car il désigne les détritus jetés au sol, pas les ordures collectées.
\item Trou 2 : « collected by \ldots\ who empty » — le pronom relatif \eng{who} exige une personne : \textbf{refuse collectors} (éboueurs).
\item Trou 3 : « empty the \ldots\ into large trucks » — on vide des conteneurs : \textbf{bins} (poubelles).
\item Trou 4 : « drop \ldots\ in the streets » — jeter des détritus par terre : \textbf{litter} (collocation : \eng{to drop litter}).
\item Trou 5 : « ends up in a \ldots\ outside the city » — un lieu d'enfouissement : \textbf{landfill} (décharge contrôlée).
\item Trou 6 : « factories that \ldots\ plastic and metal » — une action de transformation : \textbf{recycle}.
\end{itemize}
Texte complété : « In Yaoundé, household \textbf{waste} is collected twice a week by \textbf{refuse collectors} who empty the \textbf{bins} into large trucks. Unfortunately, many people still drop \textbf{litter} in the streets instead of using the containers provided. Most of the rubbish ends up in a \textbf{landfill} outside the city, although some materials could be sent to factories that \textbf{recycle} plastic and metal. »
Conclusion : toujours exploiter les indices grammaticaux (article, pronom relatif, verbe) avant de choisir le mot.
\end{exoresolu}

\begin{methode}[Réemployer le vocabulaire dans une production personnelle]
Pour produire un paragraphe ou une prise de parole riche et correcte sur l'environnement :
\begin{enumerate}
\item \textbf{Prépare une banque de mots} : avant d'écrire, liste 8 à 10 mots du thème (ex. \eng{household waste, dustbin, bin lorry, landfill, to sort, to recycle, compost, pollution}).
\item \textbf{Structure ton propos} : situation (le problème) $\rightarrow$ causes $\rightarrow$ solutions (les 3 R : \eng{reduce, reuse, recycle}) $\rightarrow$ conclusion personnelle.
\item \textbf{Utilise des collocations exactes} : \eng{to collect rubbish}, \eng{to reduce waste}, \eng{to raise awareness}, \eng{to protect the environment}, \eng{waste management}. Évite les combinaisons inventées.
\item \textbf{Varie les structures} : alterne voix active et passive (\eng{The council collects the bins} / \eng{The bins are collected every Monday}).
\item \textbf{Relis-toi} : vérifie l'accord sujet-verbe, les articles et les prépositions ; un mot de vocabulaire juste mais mal inséré perd des points.
\end{enumerate}
\end{methode}

\begin{exoresolu}[Production écrite guidée — type Bac]
Énoncé. Write a paragraph of about 80 words on the following topic: « How can your school improve waste management? » Use at least five words or expressions from the lesson.
\tcblower
Solution détaillée.
Démarche : on suit le plan problème $\rightarrow$ solutions $\rightarrow$ conclusion, en insérant le lexique de la leçon de façon naturelle.
Réponse rédigée :
« Our school produces a lot of \textbf{household waste}, especially plastic bottles and paper. At the moment, everything is thrown into the same \textbf{bins} and later burnt, which pollutes the air. To improve \textbf{waste management}, the school should provide separate containers so that students can \textbf{sort} paper, plastic and organic matter. The organic waste could be turned into \textbf{compost} for the school garden. In addition, the administration should \textbf{raise awareness} among students about the three Rs: \textbf{reduce, reuse and recycle}. If everyone makes an effort, our school will become cleaner and healthier. »
Justification : cinq éléments du lexique sont employés avec leurs collocations correctes (\eng{to sort waste}, \eng{to raise awareness}, \eng{waste management}) ; la voix passive et les modaux (\eng{should, could}) structurent l'argumentation. Conclusion : un paragraphe court mais dense, avec lexique précis et grammaire maîtrisée, obtient la meilleure note.
\end{exoresolu}

\begin{methode}[Déjouer les faux amis du domaine des déchets]
\begin{enumerate}
\item \textbf{Méfie-toi des mots transparents} : un mot anglais qui ressemble au français n'a pas toujours le même sens. \eng{Garbage} ne veut pas dire « garage » ; \eng{a dump} n'est pas « une dune ».
\item \textbf{Apprends les paires piégeuses} : « un déchet » se dit \eng{waste} ou \eng{rubbish}, jamais \eng{a dirt} ; « la propreté » se dit \eng{cleanliness}, pas \eng{cleanness} dans ce contexte ; « jeter » se dit \eng{to throw away}, pas \eng{to launch}.
\item \textbf{Évite les calques syntaxiques} : on dit \eng{to take out the rubbish} (sortir la poubelle), pas \eng{to go out the rubbish} ; \eng{to pick up litter} (ramasser), pas \eng{to take litter}.
\item \textbf{Vérifie dans un dictionnaire bilingue} les mots que tu n'as jamais vus en contexte.
\item \textbf{Mémorise les expressions entières}, pas les mots isolés : \eng{waste disposal} (élimination des déchets), \eng{sewage} (eaux usées), \eng{hazardous waste} (déchets dangereux).
\end{enumerate}
\end{methode}

\begin{exoresolu}[Corriger les calques — type Bac]
Énoncé. Each sentence contains a mistake caused by French interference. Correct it.
\begin{enumerate}
\item The garbage is cleaned every morning in our quarter.
\item We must launch our rubbish in the dustbin.
\item The government should sensitize the population about recycling.
\item The wastes are collected on Tuesdays.
\end{enumerate}
\tcblower
Solution détaillée.
\begin{enumerate}
\item « The garbage is cleaned » est un calque de « les ordures sont nettoyées ». En anglais, on \emph{collecte} les ordures : \eng{The rubbish is collected every morning in our neighbourhood.} (\eng{quarter} est aussi un faux ami : « quartier » se dit \eng{neighbourhood} ou \eng{district}.)
\item \eng{To launch} signifie « lancer » (un produit, une fusée), pas « jeter » : \eng{We must throw our rubbish into the dustbin.}
\item \eng{To sensitize} existe mais est médical/technique ; pour « sensibiliser » on dit \eng{to raise awareness among} : \eng{The government should raise awareness among the population about recycling.}
\item \eng{Waste} est indénombrable au sens d'« ordures » : pas de pluriel et verbe au singulier : \eng{Waste is collected on Tuesdays.}
\end{enumerate}
Conclusion : chaque correction repose sur une règle précise (collocation, faux ami, indénombrable) ; justifier ses choix est la clé des exercices de correction au Bac.
\end{exoresolu}

\begin{methode}[Construire et employer les familles de mots]
\begin{enumerate}
\item \textbf{Identifie la racine} : \eng{pollute} (verbe) $\rightarrow$ \eng{pollution} (nom) $\rightarrow$ \eng{polluted} (adjectif) ; \eng{recycle} $\rightarrow$ \eng{recycling} (nom) $\rightarrow$ \eng{recyclable} (adjectif).
\item \textbf{Repère la nature demandée} dans la phrase : après un article il faut un nom (\eng{the pollution}), après un auxiliaire un participe ou un infinitif, avant un nom un adjectif (\eng{recyclable materials}).
\item \textbf{Soigne l'orthographe des suffixes} : \eng{environment} (avec un \emph{n} avant \emph{-ment}), \eng{disposal} (un seul \emph{s}, double \emph{p} dans \eng{disposed}), \eng{beginning} (double \emph{n}).
\item \textbf{Travaille la prononciation} en lettres françaises : \eng{garbage} « gâr-bidj », \eng{sewage} « sou-idj », \eng{litter} « li-teur », \eng{recycle} « ri-saï-keul ». Attention : \eng{waste} se prononce « ouéiste » et se confond à l'oreille avec \eng{waist} (la taille).
\item \textbf{Teste la dérivation} : si tu hésites entre deux formes, remplace par un mot français de la même famille (« polluer / pollution / pollué ») pour vérifier la nature grammaticale.
\end{enumerate}
\end{methode}

\begin{exoresolu}[Dérivation et orthographe — type Bac]
Énoncé. Complete each sentence with the correct form of the word in brackets.
\begin{enumerate}
\item Plastic bags are not \ldots\ldots\ldots\ and stay in the soil for centuries. (recycle)
\item The \ldots\ldots\ldots\ of waste is a serious problem in big cities. (dispose)
\item Air \ldots\ldots\ldots\ causes many respiratory diseases. (pollute)
\item We should protect our \ldots\ldots\ldots\ for future generations. (environment)
\end{enumerate}
\tcblower
Solution détaillée.
\begin{enumerate}
\item Après \eng{are not}, il faut un adjectif : \textbf{recyclable} (« recyclable »). Phrase : \eng{Plastic bags are not recyclable and stay in the soil for centuries.}
\item Après l'article \eng{the}, il faut un nom : \textbf{disposal} (attention : un seul \emph{s}). Collocation : \eng{waste disposal}. Phrase : \eng{The disposal of waste is a serious problem in big cities.}
\item Sujet de la phrase, donc nom : \textbf{pollution} (verbe \eng{pollute} + suffixe \emph{-tion}). Phrase : \eng{Air pollution causes many respiratory diseases.}
\item Après le possessif \eng{our}, il faut un nom : \textbf{environment} (ne pas oublier le \emph{n} final avant \emph{-ment}). Phrase : \eng{We should protect our environment for future generations.}
\end{enumerate}
Conclusion : la position du mot dans la phrase (après article, possessif, auxiliaire) indique toujours sa nature grammaticale ; c'est le réflexe n°1 des exercices de dérivation.
\end{exoresolu}

\begin{attention}[Les 5 erreurs qui coûtent des points au Bac]
\begin{enumerate}
\item \textbf{Mettre \eng{waste} au pluriel} : \eng{waste} est indénombrable. Écris \eng{a lot of waste is produced}, jamais \eng{wastes are produced}.
\item \textbf{Calquer « sensibiliser »} : ne dis pas \eng{to sensitize the population} ; dis \eng{to raise awareness among the population} ou \eng{to make people aware of the problem}.
\item \textbf{Confondre les faux amis} : \eng{garbage} (ordures) n'est pas « garage » ; « quartier » se dit \eng{neighbourhood}, pas \eng{quarter} ; « jeter » se dit \eng{to throw away}, pas \eng{to launch}.
\item \textbf{Fautes d'orthographe sur les mots-clés} : \eng{environment} (avec \emph{n}), \eng{disposal} (un seul \emph{s}), \eng{beginning} (double \emph{n}), \eng{recycling} (un seul \emph{c} après \eng{re-}).
\item \textbf{Oublier les prépositions des collocations} : on dit \eng{to dispose \textbf{of} waste}, \eng{to throw something \textbf{away}}, \eng{to put rubbish \textbf{in} the bin}, \eng{to protect the environment \textbf{from} pollution}. Une préposition fausse ou absente fait perdre le point de la phrase.
\end{enumerate}
\end{attention}

\begin{exercice}[Entraînement Bac — Vocabulary in Use]
\begin{enumerate}
\item Complete with the correct word: \eng{litter / waste / rubbish / garbage}.
\begin{enumerate}
\item \eng{Household \_\_\_\_\_\_ is collected on Mondays.}
\item \eng{Don't drop \_\_\_\_\_\_ in the park; use the bin.}
\item \eng{The \_\_\_\_\_\_ truck comes at 6 a.m. (AmE context).}
\item \eng{Industrial \_\_\_\_\_\_ pollutes the river.}
\end{enumerate}
\item Choose the correct preposition: \eng{of / away / in / into / from}.
\begin{enumerate}
\item We must dispose \_\_\_\_\_\_ chemical waste safely.
\item Throw the paper \_\_\_\_\_\_ the recycling bin.
\item Sort the waste \_\_\_\_\_\_ three categories.
\item Protect the forest \_\_\_\_\_\_ illegal dumping.
\end{enumerate}
\item Word formation: use the word in brackets to form a word that fits the gap.
\begin{enumerate}
\item The \_\_\_\_\_\_ (collect) of refuse is a municipal service.
\item This packaging is not \_\_\_\_\_\_ (recycle); it must be incinerated.
\item \_\_\_\_\_\_ (pollute) levels have dropped since the new law.
\item A \_\_\_\_\_\_ (scavenge) was seen sorting through the dumpsite.
\end{enumerate}
\item Translate into English (watch out for false friends!).
\begin{enumerate}
\item Les éboueurs vident les poubelles deux fois par semaine.
\item Il faut sensibiliser la population au tri des déchets.
\item Ce quartier est jonché de détritus.
\item Nous ne devons pas jeter nos ordures dans la rue.
\end{enumerate}
\end{enumerate}
\end{exercice}

\begin{corrige}[Exercice 1]
\begin{enumerate}
\item \textbf{waste} (uncountable, household waste).
\item \textbf{litter} (collocation \eng{drop litter}, public place).
\item \textbf{garbage} (AmE context \eng{garbage truck}).
\item \textbf{waste} (industrial waste, uncountable).
\end{enumerate}
\end{corrige}

\begin{corrige}[Exercice 2]
\begin{enumerate}
\item \textbf{of} (\eng{dispose of}).
\item \textbf{into} (mouvement vers l'intérieur \eng{throw into}).
\item \textbf{into} (\eng{sort into categories}).
\item \textbf{from} (\eng{protect from}).
\end{enumerate}
\end{corrige}

\begin{corrige}[Exercice 3]
\begin{enumerate}
\item \textbf{collection} (nom après \eng{the}).
\item \textbf{recyclable} (adjectif après \eng{not}, \eng{be}).
\item \textbf{Pollution} (sujet, nom en \emph{-tion}).
\item \textbf{scavenger} (nom d'agent, personne qui fouille les décharges).
\end{enumerate}
\end{corrige}

\begin{corrige}[Exercice 4]
\begin{enumerate}
\item \eng{Refuse collectors empty the bins twice a week.} (pas \eng{clean the bins}, pas \eng{quarter}).
\item \eng{We must raise awareness among the population about waste sorting.} (pas \eng{sensitize}, \eng{sorting} ou \eng{separation}).
\item \eng{This neighbourhood is littered with rubbish / litter.} (\eng{neighbourhood}, \eng{be littered with}).
\item \eng{We must not throw our rubbish in the street / on the ground.} (pas \eng{launch}, \eng{throw in/on}).
\end{enumerate}
\end{corrige}

\begin{experience}[Activité orale : Débat « Zero Waste School »]
\textbf{Contexte} : Le club environnement de ton lycée propose un plan « Zero Waste » pour l'an prochain.
\textbf{Rôles} :
\begin{itemize}
\item \textit{Student A (Club President)} : présente le plan (tri, compost, gourdes, fin du plastique jetable). Utilise : \eng{reduce, reuse, recycle, compost, raise awareness, waste management, bins, sorting}.
\item \textit{Student B (Sceptic)} : pose des problèmes (coût, habitudes, manque de poubelles, vendeurs ambulants). Utilise : \eng{litter, overflowing bins, illegal dumping, habits, cost, difficult}.
\item \textit{Student C (Mediator / Principal)} : résume, propose un vote, insiste sur \eng{step by step, pilot project, evaluation}.
\end{itemize}
\textbf{Consigne} : 5 min de préparation (notes en anglais), 5 min de débat en groupe de 3. Le prof note l'usage du vocabulaire ciblé (check-list au tableau).
\end{experience}

\newpage

\coursec{Exercices}

\courssub{Je vérifie mes connaissances}

\begin{exercice}[QCM : Vocabulaire de base]
Chaque phrase ne comporte qu'une seule réponse correcte. Choisissez la bonne option (a, b, c ou d).
\begin{enumerate}
\item The large truck that collects household waste in our neighbourhood early every Monday morning is called a \dots
\begin{itemize}
\item[a)] dustcart \item[b)] dustbin \item[c)] landfill \item[d)] litter
\end{itemize}
\item Plastic bottles, glass jars and aluminium cans should be placed in separate \dots to facilitate recycling.
\begin{itemize}
\item[a)] dumps \item[b)] bins \item[c)] tips \item[d)] heaps
\end{itemize}
\item The local council has decided to \dots the old industrial site because the soil is contaminated with toxic chemicals.
\begin{itemize}
\item[a)] recycle \item[b)] compost \item[c)] remediate \item[d)] litter
\end{itemize}
\item Household waste that cannot be recycled is usually buried in a \dots
\begin{itemize}
\item[a)] landfill \item[b)] compost heap \item[c)] recycling centre \item[d)] incinerator
\end{itemize}
\item In many Cameroonian cities, \dots collectors push carts through the streets to gather rubbish from households.
\begin{itemize}
\item[a)] garbage \item[b)] waste \item[c)] refuse \item[d)] scavenger
\end{itemize}
\end{enumerate}
\end{exercice}

\begin{exercice}[Vrai ou Faux justifié]
Lisez le texte ci-dessous et dites si les affirmations sont vraies (V) ou fausses (F). Justifiez chaque réponse par une phrase du texte.

\begin{textebox}[Waste Management in Buea]
Buea Municipal Council has launched a new waste management strategy. Residents are now required to sort their rubbish into three categories: biodegradable waste goes into green bins, recyclable materials such as paper, plastic and metal go into yellow bins, and residual waste goes into grey bins. Collection takes place twice a week. A new composting plant has been built on the outskirts of town to treat organic waste and produce fertiliser for local farmers. The council warns that anyone caught dumping waste illegally will face a heavy fine.
\end{textebox}

\begin{enumerate}
\item Residents must separate their waste into four different categories.
\item Organic waste is collected in yellow bins.
\item The composting plant is located in the town centre.
\item Illegal dumping is punishable by a fine.
\item Waste is collected once a week.
\end{enumerate}
\end{exercice}

\begin{exercice}[Vocabulaire : Collocations et expressions]
Complétez les phrases suivantes avec le mot approprié choisi dans la liste ci-dessous. Il y a deux mots en trop.
\medskip
\noindent \textit{fly-tipping \quad kerbside \quad hazardous \quad landfill \quad litter \quad recycle \quad skip \quad bin \quad collect \quad dump}
\begin{enumerate}
\item Construction companies often hire a large \dots to remove rubble from building sites.
\item \dots waste, such as batteries and paint, must never be thrown in the ordinary dustbin.
\item The council provides a \dots collection service for paper and glass every Tuesday morning.
\item Please do not \dots chewing gum on the pavement; use the litter bin.
\item Volunteers gathered at the beach to \dots plastic waste washed up by the tide.
\item The mayor promised to crack down on \dots, which spoils the countryside.
\end{enumerate}
\end{exercice}

\begin{exercice}[Familles de mots : Formes dérivées]
Complétez le tableau suivant en trouvant les formes manquantes (nom, verbe, adjectif, personne). La première ligne est donnée en exemple.

\begin{center}
\begin{tabularx}{\linewidth}{|X|X|X|X|}
\hline
\rowcolor{popblueL} \textbf{Verbe} & \textbf{Nom (chose/action)} & \textbf{Nom (personne/agent)} & \textbf{Adjectif} \\
\hline
collect & collection & collector & collective \\
\hline
recycle & & & \\
\hline
dispose & & & \\
\hline
pollute & & & \\
\hline
reduce & & & \\
\hline
incinerate & & & \\
\hline
\end{tabularx}
\end{center}
\end{exercice}

\courssub{Je m'entraîne}

\begin{exercice}[Traduction : Du français vers l'anglais]
Traduisez les phrases suivantes en anglais en réemployant le vocabulaire de la leçon.
\begin{enumerate}
\item Les éboueurs passent ramasser les poubelles tous les mardis et vendredis.
\item Il est interdit de jeter ses déchets dans la rue ; on risque une amende.
\item La mairie a installé de nouveaux conteneurs pour le tri sélectif (verre, papier, plastique).
\item Les déchets dangereux doivent être traités dans une usine spéciale.
\item Nous devons réduire notre production de déchets et composter nos épluchures de légumes.
\end{enumerate}
\end{exercice}

\begin{exercice}[Transformation de phrases]
Réécrivez chaque phrase en commençant par les mots indiqués, sans en changer le sens.
\begin{enumerate}
\item They collect the rubbish early in the morning.
\begin{quote} The rubbish \dots
\end{quote}
\item "Don't throw your batteries in the bin," said the teacher.
\begin{quote} The teacher warned the students \dots
\end{quote}
\item Recycling paper saves trees.
\begin{quote} Trees \dots
\end{quote}
\item The landfill site is full. They must find a new one.
\begin{quote} Since the landfill site \dots
\end{quote}
\item It is compulsory to sort your waste in this city.
\begin{quote} You \dots
\end{quote}
\end{enumerate}
\end{exercice}

\begin{exercice}[Texte à trous : Gestion des déchets à Douala]
Complétez l'article ci-dessous avec les mots de la liste. Chaque mot s'utilise une seule fois. Il y a deux mots en trop.
\medskip
\noindent \textit{bins \quad collect \quad compost \quad dump \quad fine \quad hazardous \quad landfill \quad litter \quad recycle \quad sort \quad truck \quad waste}

\begin{textebox}[Douala tackles waste crisis]
Douala, the economic capital of Cameroon, generates thousands of tonnes of \textbf{(1)} every day. To address the crisis, HYSACAM has deployed new \textbf{(2)} trucks to \textbf{(3)} household rubbish more efficiently. Residents are urged to \textbf{(4)} their rubbish at home: organic matter can be used to \textbf{(5)} fertiliser, while plastic and metal should be placed in yellow \textbf{(6)} to be \textbf{(7)}. The main \textbf{(8)} site at PK10 is nearly full, prompting authorities to search for a new location. Illegal \textbf{(9)} of rubbish in streams causes flooding during the rainy season. Anyone caught throwing \textbf{(10)} on the street now risks a heavy \textbf{(11)}.
\end{textebox}
\end{exercice}

\begin{exercice}[Appariements : Définitions et termes]
Associez chaque terme (1--8) à sa définition correcte (a--h).
\begin{center}
\begin{tabularx}{\linewidth}{|X|X|}
\hline
\rowcolor{popblueL} \textbf{Terme} & \textbf{Définition} \\
\hline
1. Wheelie bin & a. A place where waste is buried in the ground. \\
2. Landfill & b. Waste that poses a threat to health or environment. \\
3. Kerbside collection & c. A large container for building waste. \\
4. Hazardous waste & d. A bin with wheels, used for household waste. \\
5. Skip & e. The process of turning organic waste into fertiliser. \\
6. Composting & f. Service where bins are emptied from the pavement. \\
7. Incinerator & g. A person who collects and sells recyclable materials. \\
8. Scavenger / Waste picker & h. A furnace for burning waste at high temperature. \\
\hline
\end{tabularx}
\end{center}
\end{exercice}

\courssub{Je me prépare au Bac}

\begin{exercice}[Compréhension de texte et questions de langue -- Type Bac]
Lisez attentivement le texte suivant et répondez aux questions qui suivent.

\begin{textebox}[The Challenge of Plastic in Yaoundé]
Yaoundé, the political capital of Cameroon, faces a mounting plastic waste crisis. Every day, tonnes of single-use plastic bags, bottles and packaging clog the city's drainage channels, causing severe flooding during the rainy season. Despite a government ban on non-biodegradable plastics enacted in 2014, enforcement remains weak. Street vendors still package hot food in thin polythene bags, and markets are flooded with imported plastic packaging.

A local NGO, "Green Horizon", has launched a community-based recycling initiative. They have set up collection points in five neighbourhoods where residents can bring clean plastic bottles in exchange for a small cash incentive. The collected plastic is shredded and sold to a factory in Douala that produces polyester fibres for the textile industry. "We turn a problem into a resource," explains the project coordinator. "But we need more sorting at source. If waste is mixed, it is too costly to clean."

The city council has promised to provide more public bins and to organise regular sensitisation campaigns in schools. However, experts argue that without a formalised waste-picker integration programme and stricter penalties for illegal dumping, the streets of Yaoundé will remain littered.
\end{textebox}

\textbf{Comprehension questions (answer in English, using your own words as far as possible):}
\begin{enumerate}
\item What are the main consequences of plastic waste blocking drainage channels in Yaoundé?
\item Why has the 2014 ban on non-biodegradable plastics not solved the problem?
\item Describe the "Green Horizon" initiative. How does it encourage participation?
\item What happens to the plastic collected by the NGO?
\item According to the text, what two measures are necessary to improve the situation in the long term?
\end{enumerate}

\textbf{Language work:}
\begin{enumerate}
\setcounter{enumi}{5}
\item Find in the text words or expressions that mean the same as:
\begin{itemize}
\item[a)] blocking / obstructing (paragraph 1)
\item[b)] application of a law / making people obey (paragraph 1)
\item[c)] cut into small pieces (paragraph 2)
\item[d)] at the origin / where the waste is produced (paragraph 2)
\item[e)] people who sell goods in the street (paragraph 1)
\end{itemize}
\item Rewrite the following sentences as indicated:
\begin{enumerate}
\item "We turn a problem into a resource," explains the coordinator. (Reported speech)
\item The city council has promised to provide more bins. (Passive voice)
\item If waste is mixed, it is too costly to clean. (Rewrite using "unless")
\end{enumerate}
\item Give the correct form of the words in brackets:
\begin{enumerate}
\item The \dots (collect) of plastic waste creates jobs for many young people.
\item Burning plastic releases \dots (toxic) fumes into the atmosphere.
\item The government must \dots (force) the law on plastic bags.
\item A \dots (sustain) solution requires the participation of everyone.
\end{enumerate}
\end{enumerate}
\end{exercice}

\begin{exercice}[Traduction : Thème -- Type Bac]
Traduisez le passage suivant en anglais. (Vocabulaire imposé : \textit{waste management, landfill, sorting, recycling plant, fine, hazardous, collector, kerbside}).

« La gestion des déchets à Bamenda pose de sérieux problèmes. La décharge principale est saturée et les déchets dangereux sont souvent mélangés aux ordures ménagères. Le tri à la source reste rare, bien qu'une usine de recyclage ait ouvert récemment. Les éboueurs ramassent les poubelles au bord du trottoir deux fois par semaine. Ceux qui jettent leurs déchets n'importe où risquent une forte amende. Pour améliorer la situation, il faut éduquer la population et intégrer les récupérateurs informels dans le système officiel. »
\end{exercice}

\begin{exercice}[Expression écrite : Article de magazine -- Type Bac]
\textbf{Subject:} Your town is facing a serious waste management problem (overflowing bins, illegal dumping, blocked drains, health risks). Write an article for your school magazine proposing concrete solutions to the municipal authorities and the population.

\textbf{Instructions:}
\begin{itemize}
\item Length: 150--200 words.
\item Structure: Title, Introduction (description of the problem), Body (at least 3 distinct solutions with arguments), Conclusion (call to action).
\item Use appropriate vocabulary from the lesson (\cle{landfill}, \cle{sorting}, \cle{recycling}, \cle{composting}, \cle{fine}, \cle{kerbside collection}, \cle{hazardous waste}, \cle{litter}, \cle{waste picker}, \cle{incinerator}).
\item Pay attention to connectors (\cle{moreover}, \cle{consequently}, \cle{therefore}, \cle{in addition}, \cle{however}).
\item Check spelling, verb tenses and sentence structure.
\end{itemize}
\end{exercice}

\newpage

\coursec{Corrigés des exercices}

\begin{corrige}[Exercice 1 — QCM : Vocabulaire de base]
\textbf{Réponses : 1. a \quad 2. b \quad 3. c \quad 4. a \quad 5. c}

\textbf{Justifications :}
\begin{enumerate}
\item \textbf{a) dustcart} (ou \eng{refuse truck} en britannique, \eng{garbage truck} en américain) : c'est le camion qui ramasse les ordures. \eng{Dustbin} = la poubelle (le contenant) ; \eng{landfill} = la décharge ; \eng{litter} = les détritus jetés au sol. Attention à ne pas confondre le \emph{contenant} (\eng{bin}) et le \emph{véhicule} (\eng{cart} / \eng{truck}).
\item \textbf{b) bins} : on trie les matières recyclables dans des poubelles séparées (\eng{separate bins}). \eng{Dumps} et \eng{tips} désignent des décharges (sauvages ou publiques), \eng{heaps} des tas.
\item \textbf{c) remediate} : \eng{to remediate a site} = dépolluer / assainir un site contaminé. \eng{Recycle} (recycler), \eng{compost} (composter) et \eng{litter} (joncher de détritus) ne conviennent pas pour un sol contaminé par des produits chimiques.
\item \textbf{a) landfill} : les déchets non recyclables sont enfouis dans une décharge. Un \eng{incinerator} brûle les déchets (on ne les y « enterre » pas) ; un \eng{compost heap} ne reçoit que des déchets organiques.
\item \textbf{c) refuse} : \eng{refuse collectors} (également \eng{waste collectors}) = les éboueurs. \eng{Garbage collector} existe en anglais américain, mais dans la liste proposée seul \eng{refuse} forme la collocation britannique correcte avec \eng{collectors}. \eng{Scavenger} désigne un récupérateur informel, pas un employé municipal.
\end{enumerate}
\end{corrige}

\begin{corrige}[Exercice 2 — Vrai ou Faux justifié]
\begin{enumerate}
\item \textbf{F} — Le texte dit : \eng{“Residents are now required to sort their rubbish into three categories”} (trois catégories, pas quatre).
\item \textbf{F} — Le texte dit : \eng{“biodegradable waste goes into green bins”} ; les bacs jaunes (\eng{yellow bins}) sont réservés aux matières recyclables (\eng{recyclable materials such as paper, plastic and metal}).
\item \textbf{F} — Le texte dit : \eng{“A new composting plant has been built on the outskirts of town”} (\eng{on the outskirts} = à la périphérie, en banlieue), donc pas au centre-ville.
\item \textbf{V} — Le texte dit : \eng{“anyone caught dumping waste illegally will face a heavy fine”} (\eng{to face a heavy fine} = encourir une forte amende).
\item \textbf{F} — Le texte dit : \eng{“Collection takes place twice a week”} (deux fois par semaine, pas une fois).
\end{enumerate}
\textbf{Point de méthode :} au Bac, la justification doit être une citation exacte du texte, entre guillemets, sans paraphrase.
\end{corrige}

\begin{corrige}[Exercice 3 — Vocabulaire : Collocations et expressions]
\textbf{Réponses : 1. skip \quad 2. Hazardous \quad 3. kerbside \quad 4. dump \quad 5. collect \quad 6. fly-tipping}

\textbf{Phrases complètes et traductions :}
\begin{enumerate}
\item \eng{Construction companies often hire a large skip to remove rubble from building sites.} « Les entreprises de construction louent souvent une grande benne pour évacuer les gravats des chantiers. » (Collocation : \eng{to hire a skip}.)
\item \eng{Hazardous waste, such as batteries and paint, must never be thrown in the ordinary dustbin.} « Les déchets dangereux, comme les piles et la peinture, ne doivent jamais être jetés dans la poubelle ordinaire. » (Collocation : \eng{hazardous waste}.)
\item \eng{The council provides a kerbside collection service for paper and glass every Tuesday morning.} « La municipalité assure un ramassage en bordure de trottoir du papier et du verre chaque mardi matin. » (Collocation : \eng{kerbside collection}.)
\item \eng{Please do not dump chewing gum on the pavement; use the litter bin.} « Merci de ne pas jeter votre chewing-gum sur le trottoir ; utilisez la poubelle. »
\item \eng{Volunteers gathered at the beach to collect plastic waste washed up by the tide.} « Des bénévoles se sont rassemblés sur la plage pour ramasser les déchets plastiques rejetés par la marée. »
\item \eng{The mayor promised to crack down on fly-tipping, which spoils the countryside.} « Le maire a promis de sévir contre les dépôts sauvages d'ordures, qui défigurent la campagne. » (Collocation : \eng{to crack down on} = sévir contre.)
\end{enumerate}
\textbf{Mots en trop :} \eng{landfill}, \eng{litter}, \eng{recycle}, \eng{bin}.
\end{corrige}

\begin{corrige}[Exercice 4 — Familles de mots : Formes dérivées]
\begin{center}
\begin{tabularx}{\linewidth}{|X|X|X|X|}
\hline
\rowcolor{popblueL} \textbf{Verbe} & \textbf{Nom (chose/action)} & \textbf{Nom (personne/agent)} & \textbf{Adjectif} \\
\hline
collect & collection & collector & collective \\
\hline
recycle & recycling & recycler & recyclable / recycled \\
\hline
dispose & disposal & disposer & disposable \\
\hline
pollute & pollution & polluter & polluted / polluting \\
\hline
reduce & reduction & reducer & reduced / reducible \\
\hline
incinerate & incineration & incinerator & incinerated \\
\hline
\end{tabularx}
\end{center}
\textbf{Remarques :}
\begin{itemize}
\item \eng{Incinerator} désigne à la fois la machine (l'incinérateur, le four) et, par extension, l'installation : c'est un nom d'\emph{agent matériel}, pas de personne.
\item \eng{Disposable} (jetable) est un adjectif très fréquent : \eng{disposable nappies} = « couches jetables ». Ne pas confondre avec \eng{disposal} (nom : élimination, évacuation).
\item \eng{Recyclable} (recyclable) décrit la matière ; \eng{recycled} (recyclé) décrit le produit fini : \eng{recycled paper} = « papier recyclé ».
\item Prononciation : \eng{pollution} « po-lou-cheune », \eng{reduction} « ri-deuk-cheune », \eng{disposal} « diss-po-zeul ».
\end{itemize}
\end{corrige}

\begin{corrige}[Exercice 5 — Traduction : Du français vers l'anglais]
\begin{enumerate}
\item \eng{The refuse collectors come and empty the bins every Tuesday and Friday.} \\
Variantes acceptées : \eng{The dustmen collect the rubbish every Tuesday and Friday.} Attention : « tous les mardis et vendredis » se rend par \eng{every Tuesday and Friday} (singulier après \eng{every}) ou \eng{on Tuesdays and Fridays}.
\item \eng{It is forbidden to throw your waste in the street; you risk a fine.} \\
Variantes : \eng{Littering in the street is prohibited; offenders risk a fine.} / \eng{You can be fined for dropping litter in the street.}
\item \eng{The town council has installed new containers for waste sorting (glass, paper, plastic).} \\
Variante : \eng{... new recycling banks for glass, paper and plastic.} « Mairie » (l'administration) = \eng{town council} / \eng{municipality} ; \eng{town hall} désigne le bâtiment.
\item \eng{Hazardous waste must be treated in a special plant.} \\
Variante : \eng{Dangerous waste must be processed at a special facility.} Passif obligatoire : \eng{must be treated}.
\item \eng{We must reduce our waste production and compost our vegetable peelings.} \\
Variante : \eng{We need to cut down on the waste we produce and compost our vegetable scraps.}
\end{enumerate}
\textbf{Points de vigilance :} \eng{waste} est indénombrable (jamais \eng{wastes} au sens d'ordures ménagères) ; « ramasser les poubelles » se dit \eng{empty the bins} ou \eng{collect the rubbish}, pas \eng{collect the bins}.
\end{corrige}

\begin{corrige}[Exercice 6 — Transformation de phrases]
\begin{enumerate}
\item \eng{The rubbish is collected early in the morning.} \\
Passif au présent simple : auxiliaire \eng{is} + participe passé \eng{collected} ; l'agent (\eng{they}) disparaît car il est indéterminé.
\item \eng{The teacher warned the students not to throw their batteries in the bin.} \\
Discours rapporté avec \eng{warn} : \eng{warn somebody not to do something} ; l'impératif négatif devient \eng{not to + base verbale} ; \eng{your} devient \eng{their}.
\item \eng{Trees are saved by recycling paper.} ou mieux : \eng{Trees are saved when paper is recycled.} \\
Passif : \eng{are saved} ; on peut garder le gérondif \eng{by recycling paper} en complément de moyen.
\item \eng{Since the landfill site is full, they must find a new one.} \\
\eng{Since} (puisque) introduit la cause ; \eng{one} reprend \eng{landfill site} pour éviter la répétition.
\item \eng{You must sort your waste in this city.} ou \eng{You have to sort your waste in this city.} \\
\eng{It is compulsory to...} (il est obligatoire de) équivaut à \eng{must} / \eng{have to}.
\end{enumerate}
\textbf{Point de vigilance :} ne jamais écrire \eng{the rubbish is collect} : le participe passé est indispensable après l'auxiliaire \eng{be}.
\end{corrige}

\begin{corrige}[Exercice 7 — Texte à trous : Gestion des déchets à Douala]
\textbf{Note de relecture :} le sujet comportait une coquille au trou n° 2 : la phrase \eng{has deployed new \dots trucks} exige un nom épithète. Le mot de la liste à employer est \eng{garbage} (\eng{garbage trucks} = « camions à ordures ») ; \eng{truck} figurait dans la liste par erreur et doit être compté parmi les mots en trop.

\textbf{Réponses : 1. waste \quad 2. garbage \quad 3. collect \quad 4. sort \quad 5. compost \quad 6. bins \quad 7. recycle \quad 8. landfill \quad 9. dump \quad 10. litter \quad 11. fine}

\textbf{Texte complété et justifications :}
\begin{enumerate}
\item \eng{Douala generates thousands of tonnes of \textbf{waste} every day.} « Douala produit des milliers de tonnes de déchets chaque jour. » (\eng{waste} indénombrable.)
\item \eng{The city council has deployed new \textbf{garbage} trucks.} « La municipalité a mis en service de nouveaux camions à ordures. » (Nom épithète au singulier : \eng{garbage trucks}.)
\item \eng{...to \textbf{collect} household rubbish.} « ...pour ramasser les ordures ménagères. » (Infinitif après \eng{to}.)
\item \eng{Residents are encouraged to \textbf{sort} their rubbish.} « Les habitants sont encouragés à trier leurs ordures. »
\item \eng{Organic matter can be used to \textbf{compost}...} — lecture correcte : \eng{organic matter can be \textbf{composted} and turned into fertiliser} « la matière organique peut être compostée et transformée en engrais » ; si le sujet impose la forme de la liste, on écrit \eng{compost} après \eng{to}.
\item \eng{...yellow \textbf{bins} for recyclable materials.} « ...des bacs jaunes pour les matières recyclables. »
\item \eng{...plastic bottles are sent to be \textbf{recycle}d.} La forme correcte est le participe passé : \eng{recycled} « ...les bouteilles en plastique sont envoyées pour être recyclées. »
\item \eng{...the main \textbf{landfill} site at PK10.} « ...la décharge principale au PK10. »
\item \eng{...illegal \textbf{dump}ing of rubbish...} La forme correcte est le nom en \emph{-ing} : \eng{dumping} « ...le dépôt sauvage d'ordures... »
\item \eng{...throwing \textbf{litter} on the street.} « ...jeter des détritus dans la rue. » (\eng{litter} indénombrable.)
\item \eng{...offenders will pay a heavy \textbf{fine}.} « ...les contrevenants paieront une forte amende. » (Collocation : \eng{a heavy fine}.)
\end{enumerate}
\textbf{Mots en trop :} \eng{hazardous} et \eng{truck} (coquille du sujet, voir la note ci-dessus).

\textbf{Point de méthode :} dans un texte à trous, vérifier toujours la \emph{forme} exigée par la phrase (infinitif, participe passé, nom en \emph{-ing}) : le mot de la liste peut devoir être accordé (\eng{recycle} $\rightarrow$ \eng{recycled} ; \eng{dump} $\rightarrow$ \eng{dumping}).
\end{corrige}

\begin{corrige}[Exercice 8 — Appariements : Définitions et termes]
\textbf{Réponses : 1. d \quad 2. a \quad 3. f \quad 4. b \quad 5. c \quad 6. e \quad 7. h \quad 8. g}

\textbf{Justifications et traductions :}
\begin{itemize}
\item \textbf{1. Wheelie bin — d} : poubelle roulante (\eng{wheel} = roue) pour les ordures ménagères.
\item \textbf{2. Landfill — a} : décharge où les déchets sont enfouis dans le sol (\eng{to bury} = enfouir).
\item \textbf{3. Kerbside collection — f} : ramassage en bordure de trottoir (\eng{kerb} = bord du trottoir ; orthographe américaine : \eng{curb}).
\item \textbf{4. Hazardous waste — b} : déchets dangereux présentant une menace (\eng{threat}) pour la santé ou l'environnement.
\item \textbf{5. Skip — c} : grande benne à gravats utilisée sur les chantiers.
\item \textbf{6. Composting — e} : transformation des déchets organiques en engrais (\eng{fertiliser}).
\item \textbf{7. Incinerator — h} : four (\eng{furnace}) brûlant les déchets à haute température.
\item \textbf{8. Scavenger / Waste picker — g} : récupérateur qui collecte et revend des matières recyclables.
\end{itemize}
\end{corrige}

\begin{corrige}[Exercice 9 — Compréhension de texte et questions de langue (Type Bac)]
\textbf{Barème indicatif (20 points) :} compréhension 10 pts (2 pts par question) ; language work 10 pts (Q6 : 2,5 pts ; Q7 : 3 pts ; Q8 : 4,5 pts).

\textbf{Comprehension — réponses modèles :}
\begin{enumerate}
\item \eng{When plastic waste blocks the drainage channels, rainwater cannot flow away, which causes severe flooding in the city during the rainy season.} (Conséquence : inondations.)
\item \eng{The ban has not solved the problem because it is not properly enforced: enforcement remains weak, so street vendors still use polythene bags and markets are full of imported plastic packaging.}
\item \eng{“Green Horizon” has set up collection points in five neighbourhoods where residents can bring clean plastic bottles. It encourages participation by giving people a small cash incentive in exchange for the bottles they bring.}
\item \eng{The collected plastic is shredded and then sold to a factory in Douala, where it is turned into polyester fibres for the textile industry.}
\item \eng{According to the experts, the city needs a formalised waste-picker integration programme and stricter penalties for illegal dumping.}
\end{enumerate}

\textbf{Language work :}
\begin{enumerate}
\setcounter{enumi}{5}
\item a) \eng{clogging} ; b) \eng{enforcement} ; c) \eng{shredded} ; d) \eng{at source} ; e) \eng{street vendors}.
\item \begin{enumerate}
\item \eng{The coordinator explains that they turn a problem into a resource.} (Présent de vérité générale conservé ; \eng{we} devient \eng{they}.)
\item \eng{More bins have been promised by the city council.} (Passif au present perfect : \eng{have been promised}.)
\item \eng{Unless waste is sorted, it is too costly to clean.} (\eng{Unless} = \eng{if... not} : « si les déchets ne sont pas triés ».)
\end{enumerate}
\item \begin{enumerate}
\item \eng{The collection of plastic waste creates jobs for many young people.} (Nom : \eng{collection}.)
\item \eng{Burning plastic releases toxic fumes into the atmosphere.} (Adjectif : \eng{toxic} ; \eng{toxicity} est le nom.)
\item \eng{The government must enforce the law on plastic bags.} (Verbe : \eng{to enforce} ; collocation : \eng{enforce the law}.)
\item \eng{A sustainable solution requires the participation of everyone.} (Adjectif : \eng{sustainable} = durable.)
\end{enumerate}
\end{enumerate}
\textbf{Points de vigilance :} ne pas copier de longues phrases du texte dans les réponses de compréhension (paraphraser) ; accorder le participe passé au passif ; \eng{unless} exige une forme affirmative après lui.
\end{corrige}

\begin{corrige}[Exercice 10 — Traduction : Thème (Type Bac)]
\textbf{Barème indicatif :} 10 points — 1 point par segment correctement traduit (8 segments) + 2 points pour l'ensemble (cohérence, vocabulaire imposé employé).

\textbf{Traduction modèle :}

\eng{Waste management in Bamenda poses serious problems. The main landfill is saturated, and hazardous waste is often mixed with household rubbish. Sorting at source remains rare, although a recycling plant has recently opened. The collectors empty the bins at the kerbside twice a week. Those who dump their waste anywhere risk a heavy fine. To improve the situation, it is necessary to educate the population and to integrate informal waste pickers into the official system.}

\textbf{Points de vigilance :}
\begin{itemize}
\item « La décharge principale est saturée » : \eng{is saturated} ou \eng{is full to capacity} (orthographe : \eng{saturated}).
\item « bien qu'une usine de recyclage ait ouvert » : le subjonctif français se rend par l'indicatif en anglais : \eng{although a recycling plant has opened}.
\item « risquent une forte amende » : \eng{risk a heavy fine} ou \eng{face a heavy fine} (collocation : \eng{heavy fine}, pas \eng{strong fine}).
\item « les récupérateurs informels » : \eng{informal waste pickers} (ou \eng{informal scavengers}).
\item \eng{waste} indénombrable ; \eng{twice a week} = « deux fois par semaine ».
\end{itemize}
\end{corrige}

\begin{corrige}[Exercice 11 — Expression écrite : Article de magazine (Type Bac)]
\textbf{Barème indicatif (20 points) :} respect de la tâche et du format article 4 pts ; richesse et exactitude du vocabulaire du thème 6 pts ; correction grammaticale et orthographique 5 pts ; organisation et connecteurs 3 pts ; longueur respectée (150--200 mots) 2 pts.

\textbf{Proposition de rédaction modèle :} voir l'article ci-dessous.

\textbf{Points de vigilance :}
\begin{itemize}
\item Donner un \textbf{titre} à l'article et signer éventuellement (\eng{By...}) : c'est le format attendu.
\item Employer au moins trois solutions distinctes, chacune introduite par un connecteur (\eng{first}, \eng{moreover}, \eng{in addition}, \eng{finally}).
\item Éviter les calques : « sensibiliser » = \eng{to raise awareness} (dans un article formel ; \eng{sensitisation campaigns} est toutefois attesté dans l'usage camerounais) ; « il faut » = \eng{we must} / \eng{it is necessary to}.
\item Vérifier : présent simple pour les vérités générales, \eng{should} pour les recommandations, pas de \eng{s} à \eng{waste} ni à \eng{rubbish} (indénombrables).
\end{itemize}
\end{corrige}

\begin{textebox}[Model article — “Let's Clean Up Our Town!”]
Our town is drowning in rubbish. Bins overflow in every neighbourhood, illegal dumping is spreading along the roads, and blocked drains cause flooding whenever it rains. This situation is not only ugly: it is a serious threat to public health.

Concrete solutions exist. First, the municipal authorities should organise a reliable kerbside collection service, with enough wheelie bins for every household. Moreover, sorting at source must be encouraged: organic waste can be composted, while paper, glass and plastic should be sent to a recycling plant instead of the already saturated landfill. In addition, hazardous waste such as batteries and paint must be collected separately and treated in a special facility. However, none of this will work without discipline: anyone caught littering or fly-tipping should pay a heavy fine. Finally, informal waste pickers, who already recover tons of recyclable materials, should be integrated into the official system and given protective equipment.

Therefore, we call on the council and on every citizen to act now. A clean town is possible if we all reduce, reuse and recycle. Let us start today!
\end{textebox}

\newpage

\coursec{Fiche bilan}

\begin{popfigure}
\begin{center}\tcbox[colback=poporange,colframe=poporange,coltext=white,arc=9pt,boxsep=4pt,left=6pt,right=6pt]{\bfseries\small Waste Management}\end{center}
\vspace{-2pt}
\begin{tcbraster}[raster columns=3,raster equal height=rows,raster column skip=6pt,raster row skip=6pt,enhanced,blankest]
\begin{tcolorbox}[enhanced,colback=white,colframe=poporange,boxrule=0.9pt,arc=3pt,title={Types of Waste},fonttitle=\bfseries\scriptsize,coltitle=white,colbacktitle=poporange,left=4pt,right=4pt,top=2pt,bottom=2pt,boxsep=1pt,toptitle=1.5pt,bottomtitle=1.5pt,halign=left]
\begin{itemize}[nosep,leftmargin=*,itemsep=1pt,font=\scriptsize]
\item Household / Domestic
\item Industrial / Hazardous
\item Organic / Biodegradable
\item Recyclable / E-waste
\end{itemize}
\end{tcolorbox}
\begin{tcolorbox}[enhanced,colback=white,colframe=popgreen,boxrule=0.9pt,arc=3pt,title={Collection Services},fonttitle=\bfseries\scriptsize,coltitle=white,colbacktitle=popgreen,left=4pt,right=4pt,top=2pt,bottom=2pt,boxsep=1pt,toptitle=1.5pt,bottomtitle=1.5pt,halign=left]
\begin{itemize}[nosep,leftmargin=*,itemsep=1pt,font=\scriptsize]
\item Bin / Wheelie bin / Skip
\item Garbage truck / Lorry
\item Collector / Sanitation worker
\item Schedule / Pickup
\end{itemize}
\end{tcolorbox}
\begin{tcolorbox}[enhanced,colback=white,colframe=poppurple,boxrule=0.9pt,arc=3pt,title={Disposal Methods},fonttitle=\bfseries\scriptsize,coltitle=white,colbacktitle=poppurple,left=4pt,right=4pt,top=2pt,bottom=2pt,boxsep=1pt,toptitle=1.5pt,bottomtitle=1.5pt,halign=left]
\begin{itemize}[nosep,leftmargin=*,itemsep=1pt,font=\scriptsize]
\item Landfill / Dump
\item Incineration / Incinerator
\item Composting / Treatment plant
\item Leachate / Methane
\end{itemize}
\end{tcolorbox}
\begin{tcolorbox}[enhanced,colback=white,colframe=poppink,boxrule=0.9pt,arc=3pt,title={Recycling \& 3 Rs},fonttitle=\bfseries\scriptsize,coltitle=white,colbacktitle=poppink,left=4pt,right=4pt,top=2pt,bottom=2pt,boxsep=1pt,toptitle=1.5pt,bottomtitle=1.5pt,halign=left]
\begin{itemize}[nosep,leftmargin=*,itemsep=1pt,font=\scriptsize]
\item Reduce / Reuse / Recycle
\item Sorting / Separation
\item Facility / Plant / Depot
\item Upcycling / Circular economy
\end{itemize}
\end{tcolorbox}
\begin{tcolorbox}[enhanced,colback=white,colframe=popgold,boxrule=0.9pt,arc=3pt,title={People \& Places},fonttitle=\bfseries\scriptsize,coltitle=white,colbacktitle=popgold,left=4pt,right=4pt,top=2pt,bottom=2pt,boxsep=1pt,toptitle=1.5pt,bottomtitle=1.5pt,halign=left]
\begin{itemize}[nosep,leftmargin=*,itemsep=1pt,font=\scriptsize]
\item Local council / Authority
\item Transfer station
\item Recycling centre / Drop-off
\item Landfill site
\end{itemize}
\end{tcolorbox}
\begin{tcolorbox}[enhanced,colback=white,colframe=popdark,boxrule=0.9pt,arc=3pt,title={Collocations \& Idioms},fonttitle=\bfseries\scriptsize,coltitle=white,colbacktitle=popdark,left=4pt,right=4pt,top=2pt,bottom=2pt,boxsep=1pt,toptitle=1.5pt,bottomtitle=1.5pt,halign=left]
\begin{itemize}[nosep,leftmargin=*,itemsep=1pt,font=\scriptsize]
\item Take out the trash
\item Waste separation
\item Carbon footprint
\item Fly-tipping / Littering
\end{itemize}
\end{tcolorbox}
\begin{tcolorbox}[enhanced,colback=white,colframe=popblue!80!poppurple,boxrule=0.9pt,arc=3pt,title={Word Families \& False Friends},fonttitle=\bfseries\scriptsize,coltitle=white,colbacktitle=popblue!80!poppurple,left=4pt,right=4pt,top=2pt,bottom=2pt,boxsep=1pt,toptitle=1.5pt,bottomtitle=1.5pt,halign=left]
\begin{itemize}[nosep,leftmargin=*,itemsep=1pt,font=\scriptsize]
\item Dispose (of) $\neq$ Disposer
\item Refuse (n/ v) $\neq$ Refuser
\item Sensible $\neq$ Sensitive
\item Suffixes: -age, -ment, -tion
\end{itemize}
\end{tcolorbox}
\end{tcbraster}

\legende{Carte mentale récapitulative : Vocabulaire de la gestion des déchets (collecte, élimination, recyclage)}
\end{popfigure}

\begin{bilan}

\end{bilan}

\courssub{Lexique clé : Anglais / Français}

\begin{center}
\begin{tabularx}{\linewidth}{|X|X|X|}
\hline
\rowcolor{popblueL} \textbf{English} & \textbf{Français} & \textbf{Remarque / Collocation} \\
\hline
\cle{Household waste} / \cle{Domestic waste} & Déchets ménagers & \eng{Household waste collection} (collecte des ordures ménagères) \\
\cle{Industrial waste} & Déchets industriels & Souvent \eng{hazardous} (dangereux) \\
\cle{Organic waste} / \cle{Biodegradable waste} & Déchets organiques / biodégradables & \eng{Composting} (compostage) \\
\cle{Recyclables} & Matériaux recyclables & \eng{Paper, cardboard, plastic, glass, metal} \\
\cle{E-waste} / \cle{Electronic waste} & Déchets d'équipements électriques (DEEE) & \eng{Obsolete devices} (appareils obsolètes) \\
\cle{Litter} & Déchets sauvages / détritus & Verbe : \eng{to litter} (jeter par terre) ; \eng{Anti-littering campaign} \\
\cle{Fly-tipping} & Dépôt sauvage (illégal) & Spécifique UK ; \eng{Dumping} (US/général) \\
\hline
\end{tabularx}
\end{center}

\vspace{0.3cm}
\begin{center}
\begin{tabularx}{\linewidth}{|X|X|X|}
\hline
\rowcolor{poporangeL} \textbf{Collection \& People} & \textbf{Français} & \textbf{Remarque} \\
\hline
\cle{Wheelie bin} (UK) / \cle{Trash can} (US) & Poubelle à roulettes / Poubelle & \cle{Bin bag} / \cle{Bin liner} = sac poubelle \\
\cle{Skip} (UK) / \cle{Dumpster} (US) & Benne à ordures (chantier) & \eng{Hire a skip} (louer une benne) \\
\cle{Garbage truck} (US) / \cle{Dustcart} / \cle{Refuse lorry} (UK) & Camion-poubelle & \eng{Bin lorry} (UK courant) \\
\cle{Sanitation worker} / \cle{Refuse collector} / \cle{Dustman} (UK) & Éboueur / Agent de salubrité & Terme respectueux : \eng{sanitation worker} \\
\cle{Pickup} / \cle{Collection day} & Jour de ramassage & \eng{Weekly pickup} (ramassage hebdomadaire) \\
\cle{Transfer station} & Centre de transit / quai de transfert & Tri avant enfouissement/incinération \\
\hline
\end{tabularx}
\end{center}

\vspace{0.3cm}
\begin{center}
\begin{tabularx}{\linewidth}{|X|X|X|}
\hline
\rowcolor{popgreenL} \textbf{Disposal, Recycling \& 3 Rs} & \textbf{Français} & \textbf{Remarque} \\
\hline
\cle{Landfill} / \cle{Landfill site} & Décharge (contrôlée) & \eng{To landfill} (enfouir) ; \eng{Landfill tax} \\
\cle{Incineration} / \cle{Incinerator} & Incinération / Incinérateur & \eng{Waste-to-energy} (valorisation énergétique) \\
\cle{Composting} & Compostage & \eng{Compost heap/bin} (tas/bac de compost) \\
\cle{Reduce, Reuse, Recycle} & Réduire, Réutiliser, Recycler & Les \cle{3 Rs} (hiérarchie des déchets) \\
\cle{Sorting} / \cle{Waste separation} & Tri sélectif & \eng{Sorting bins} (poubelles de tri) \\
\cle{Recycling facility} / \cle{Plant} & Centre / Usine de recyclage & \cle{Materials Recovery Facility} (MRF) \\
\cle{Upcycling} & Surcyclage / Upcycling & Créer objet valeur supérieure \\
\cle{Circular economy} & Économie circulaire & Modèle durable vs linéaire \\
\hline
\end{tabularx}
\end{center}

\courssub{Formation des mots (Word Families) — Règles à connaître}

\begin{center}
\begin{tabularx}{\linewidth}{|X|X|X|X|}
\hline
\rowcolor{poppurpleL} \textbf{Verbe} & \textbf{Nom (action/résultat)} & \textbf{Nom (agent/lieu)} & \textbf{Adjectif / Participe} \\
\hline
\eng{collect} & \eng{collection} & \eng{collector} & \eng{collective} / \eng{collected} \\
\eng{dispose (of)} & \eng{disposal} & — & \eng{disposable} (jetable) \\
\eng{recycle} & \eng{recycling} & \eng{recycler} & \eng{recyclable} / \eng{recycled} \\
\eng{incinerate} & \eng{incineration} & \eng{incinerator} & \eng{incinerated} \\
\eng{compost} & \eng{composting} / \eng{compost} & \eng{composter} (bac) & \eng{compostable} \\
\eng{pollute} & \eng{pollution} & \eng{polluter} & \eng{polluted} / \cle{polluting} \\
\eng{contaminate} & \eng{contamination} & — & \eng{contaminated} \\
\eng{separate} & \eng{separation} & \eng{separator} & \eng{separate} / \eng{separated} \\
\hline
\end{tabularx}
\end{center}

\courssub{Préfixes productifs}
\begin{itemize}
\item \cle{re-} (à nouveau) : \eng{reuse, recycle, recover, reprocess, rethink}.
\item \cle{mis-} (mal) : \eng{mismanage, misplace, mis-sort} (mal trier).
\item \cle{non-} (non-) : \eng{non-recyclable, non-biodegradable, non-hazardous}.
\item \cle{over-} (excès) : \eng{overpackaging} (sur-emballage), \eng{overconsumption}.
\end{itemize}

\courssub{Méthodes express}

\begin{methode}[Comment décrire le circuit d'un déchet (Speaking/ Writing)]
\begin{enumerate}
\item \textbf{Identify the source} : \eng{“Households in Yaoundé generate…”}
\item \textbf{Name the container} : \eng{“Residents put waste in wheelie bins / skips.”}
\item \textbf{Describe collection} : \eng{“Garbage trucks collect bins on Tuesdays.”}
\item \textbf{Explain transfer} : \eng{“Waste is taken to a transfer station for sorting.”}
\item \textbf{State final destination} : \eng{“Recyclables go to a recycling plant; residual waste goes to landfill / incinerator.”}
\item \textbf{Use passive voice} : \eng{“Waste is sorted / is incinerated / is buried.”}
\end{enumerate}
\end{methode}

\begin{methode}[Rédaction d'une plainte formelle (Letter of complaint)]
\begin{enumerate}
\item \textbf{Header} : Your address / Date / Council address.
\item \textbf{Subject line} : \eng{Re: Missed garbage collection / Illegal dumping in [Quarter]}.
\item \textbf{Facts} : \eng{“The bins in [Street] have not been emptied for two weeks.”}
\item \textbf{Consequences} : \eng{“This attracts pests / causes bad smells / health hazard.”}
\item \textbf{Request} : \eng{“I urge you to schedule an immediate pickup / provide more bins.”}
\item \textbf{Closing} : \eng{“I look forward to your prompt action.”} + \eng{Yours faithfully,}
\end{enumerate}
\end{methode}

\courssub{Faux amis \& Pièges de prononciation}

\begin{center}
\begin{tabularx}{\linewidth}{|X|X|X|}
\hline
\rowcolor{poppinkL} \textbf{Mot anglais} & \textbf{Sens anglais} & \textbf{Piège français (Faux ami)} \\
\hline
\eng{Dispose (of)} & Se débarrasser de, éliminer & $\neq$ \emph{Disposer} (avoir à disposition / arranger). \eng{“How to dispose of batteries?”} \\
\eng{Refuse} (n.) & Ordures, déchets & $\neq$ \emph{Refuser} (verbe). Verbe \eng{refuse} = refuser. \\
\eng{Sensible} & Raisonnable, sensé & $\neq$ \emph{Sensible} (FR). FR \emph{sensible} = \eng{Sensitive}. \\
\eng{Rubbish} (UK) / \eng{Trash/Garbage} (US) & Déchets & \eng{Rubbish} (adj. familier) = nul, faux. \eng{“That's rubbish!”} \\
\eng{Tip} (UK) & Décharge & $\neq$ \emph{Pourboire} (tip) / \emph{Pointe}. \eng{“Take it to the tip.”} \\
\eng{Skip} & Benne & $\neq$ \emph{Sauter}. Verbe \eng{skip} = sauter / omettre. \\
\hline
\end{tabularx}
\end{center}

\textbf{Prononciation clé :}
\begin{itemize}
\item \eng{Waste} — \eng{Waist} (taille) \textbf{homophones}.
\item \eng{Recycle} — accent sur 2e syllabe.
\item \eng{Incinerator} — accent sur 2e syllabe.
\item \eng{Compost} — accent sur 1re.
\item \eng{Landfill} — accent sur 1re.
\end{itemize}



\begin{aretenir}[Checklist avant le Bac]
$\square$ Je distingue \cle{household}, \cle{industrial}, \cle{hazardous}, \cle{organic}, \cle{e-waste} et \cle{litter}.
$\square$ Je nomme les conteneurs : \eng{wheelie bin}, \eng{skip}, \eng{bin bag}, \eng{sorting bins}.
$\square$ Je cite les véhicules et acteurs : \eng{garbage truck / dustcart}, \eng{sanitation worker / refuse collector}.
$\square$ J'explique les 3 méthodes d'élimination : \eng{landfill}, \eng{incineration}, \eng{composting} (+ \eng{waste-to-energy}).
$\square$ J'applique la hiérarchie des \cle{3 Rs} : \eng{Reduce} $>$ \eng{Reuse} $>$ \eng{Recycle}.
$\square$ J'utilise les collocations : \eng{take out the trash}, \eng{waste separation}, \eng{fly-tipping}, \eng{carbon footprint}.
$\square$ Je maîtrise les familles : \eng{collect/collection/collector}, \eng{dispose/disposal/disposable}, \eng{recycle/recycling/recyclable}.
$\square$ J'évite les faux amis : \eng{dispose of} $\neq$ disposer, \eng{refuse} (n.) $\neq$ refuser, \eng{sensible} $\neq$ sensible.
$\square$ Je prononce correctement : \eng{waste}, \eng{recycle}, \eng{incinerator}, \eng{landfill}.
$\square$ Je sais rédiger une lettre de plainte formelle (\eng{missed collection}, \eng{illegal dumping}).
$\square$ Je peux décrire le circuit complet d'un déchet à l'oral (passif, connecteurs logiques).
$\square$ Je connais le vocabulaire local : \eng{HYSACAM}, \eng{decentralised collection}, \eng{informal sector / waste pickers}.
\end{aretenir}

\findecours

\end{document}
